% =====================================================================
%  Anatomía de un LLM — Fase 0 — Módulo 0.9
%  Estadística
%  Compilar: latexmk -r ../latexmkrc P09-estadistica.tex
% =====================================================================
\documentclass[11pt,a4paper,twoside,openany]{book}
\usepackage{../estilo-apuntes}

\newcommand{\ApuntesCurso}{Anatomía de un LLM}
\newcommand{\ApuntesModulo}{Módulo 0.9}

% Notación propia del módulo (compatible con el módulo 0.8)
\newcommand{\PP}{\mathbb{P}}
\newcommand{\N}{\mathbb{N}}
\DeclareMathOperator{\Var}{Var}
\DeclareMathOperator{\Cov}{Cov}
\DeclareMathOperator{\Cor}{Cor}
\DeclareMathOperator{\ECM}{ECM}
\DeclareMathOperator{\sesgo}{sesgo}
\DeclareMathOperator{\ee}{ee}
\DeclareMathOperator{\RIC}{RIC}
\DeclareMathOperator{\Bin}{Bin}
\DeclareMathOperator{\Ber}{Ber}
\DeclareMathOperator{\Pois}{Pois}
\DeclareMathOperator{\Expo}{Exp}
\DeclareMathOperator{\Unif}{U}
\DeclareMathOperator{\Geom}{Geom}
\DeclareMathOperator*{\argmax}{arg\,max}
\DeclareMathOperator*{\argmin}{arg\,min}
\newcommand{\Nor}{\mathcal{N}}
\newcommand{\Lik}{\mathcal{L}}
\newcommand{\Xbar}{\overline{X}}
\newcommand{\xbar}{\bar{x}}
\newcommand{\ybar}{\bar{y}}
\newcommand{\iid}{i.i.d.}
\newcommand{\ind}{\mathbf{1}}
\newcommand{\convp}{\xrightarrow{\;\PP\;}}

% Anclas de hyperref únicas para los enunciados (numerados por tema)
\providecommand{\theHdefinicion}{}\renewcommand{\theHdefinicion}{\thechapter.\arabic{definicion}}
\providecommand{\theHteorema}{}\renewcommand{\theHteorema}{\thechapter.\arabic{definicion}}
\providecommand{\theHproposicion}{}\renewcommand{\theHproposicion}{\thechapter.\arabic{definicion}}
\providecommand{\theHlema}{}\renewcommand{\theHlema}{\thechapter.\arabic{definicion}}
\providecommand{\theHcorolario}{}\renewcommand{\theHcorolario}{\thechapter.\arabic{definicion}}
\providecommand{\theHejemplo}{}\renewcommand{\theHejemplo}{\thechapter.\arabic{definicion}}
\providecommand{\theHobservacion}{}\renewcommand{\theHobservacion}{\thechapter.\arabic{definicion}}

\hypersetup{pdftitle={Módulo 0.9. Estadística},pdfauthor={Anatomía de un LLM}}

\begin{document}
\frontmatter
% ------------------------------------------------------------------ Portada
\begin{titlepage}
\centering
\vspace*{2.5cm}
{\large\scshape Anatomía de un LLM}\par
\vspace{0.4cm}
{\normalsize Fundamentos matemáticos de los modelos de lenguaje}\par
\vspace{2.2cm}
\rule{\textwidth}{0.8pt}\par\vspace{0.5cm}
{\small\scshape Fase 0 \textperiodcentered\ Preliminares matemáticos}\par\vspace{0.3cm}
{\Huge\bfseries Módulo 0.9\par}\vspace{0.3cm}
{\LARGE Estadística\par}
\vspace{0.5cm}\rule{\textwidth}{0.4pt}\par
\vspace{1.6cm}
{\large Apuntes del curso}\par
\vfill
{\small Edición de octubre de 2026}\par
\end{titlepage}

% ------------------------------------------------------- Presentación
\chapter*{Presentación del módulo}
\addcontentsline{toc}{chapter}{Presentación del módulo}
\markboth{Presentación del módulo}{}

La probabilidad parte de un modelo conocido y calcula qué datos cabe esperar; la estadística recorre el camino inverso: parte de unos datos observados y extrae conclusiones sobre el modelo que los ha generado. Entrenar un modelo de lenguaje es un problema de estimación (se buscan los parámetros que hacen más verosímil el corpus), y evaluarlo es un problema de inferencia (se mide la exactitud en un conjunto finito de preguntas y se quiere saber cuánto se parece esa cifra a la exactitud real y si la diferencia con otro modelo es auténtica o fruto del azar). Este módulo, el último de la fase de preliminares, presenta las herramientas básicas de la estadística con el nivel y la organización de un primer curso universitario, y ofrece una colección amplia de ejercicios numéricos y de demostración.

\section*{Objetivos}
Al finalizar el módulo, el estudiante debe ser capaz de:
\begin{enumerate}[label=\arabic*.]
  \item resumir un conjunto de datos con medidas de centralización, posición y dispersión, construir histogramas y diagramas de caja, y calcular e interpretar la covarianza y el coeficiente de correlación;
  \item definir estimador, sesgo, varianza, error estándar, error cuadrático medio y consistencia, y demostrar que la media muestral y la varianza muestral con denominador $n-1$ son estimadores insesgados;
  \item enunciar el teorema central del límite y aplicarlo a medias, sumas y proporciones, incluida la aproximación normal a la binomial;
  \item obtener estimadores por el método de los momentos y por máxima verosimilitud para los modelos de Bernoulli, Poisson, normal y exponencial, y construir intervalos de confianza para medias y proporciones, así como calcular el tamaño muestral necesario;
  \item plantear y resolver contrastes de hipótesis ($z$, $t$ y $\chi^{2}$ de bondad de ajuste), interpretar el $p$-valor, los errores de tipo~I y~II y la potencia, y corregir por comparaciones múltiples;
  \item deducir la recta de mínimos cuadrados, interpretar $R^{2}$ y los residuos, y distinguir correlación de causalidad.
\end{enumerate}

\section*{Requisitos previos}
Matemáticas de bachillerato y los módulos anteriores de esta fase, en especial el módulo~0.8 (combinatoria y probabilidad), cuya notación se mantiene: $\PP$ para la probabilidad, $\E$ para la esperanza y $\Var$ para la varianza. Se usan además derivadas e integrales de una variable (módulos~0.5 y~0.6) y derivadas parciales (módulo~0.7). Al comienzo del \cref{tema:2} se recuerdan las propiedades de la esperanza y la varianza que se emplean.

\section*{Organización}
\begin{center}
\small
\begin{tabular}{@{}cll@{}}
\toprule
Tema & Contenido & Fuentes principales \\
\midrule
1 & Estadística descriptiva & \cite{openstax}, \cite{rice}, \cite{mit1805} \\
2 & Muestreo y estimadores & \cite{wasserman}, \cite{bh}, \cite{stats200} \\
3 & Distribuciones muestrales y teorema central del límite & \cite{bh}, \cite{mit1805}, \cite{openstax} \\
4 & Estimación puntual e intervalos de confianza & \cite{mit1805}, \cite{wasserman}, \cite{stats200} \\
5 & Contrastes de hipótesis & \cite{mit1805}, \cite{wasserman}, \cite{stats200} \\
6 & Regresión lineal simple y correlación & \cite{openstax}, \cite{rice}, \cite{wasserman} \\
\bottomrule
\end{tabular}
\end{center}

Cada tema termina con una lista de lecturas y con una colección extensa de ejercicios graduados en tres niveles: $\star$ comprensión, $\star\star$ cálculo o práctica mecánica y $\star\star\star$ demostración. Los ejercicios marcados con \emph{Fuente} están adaptados de cursos y libros de texto universitarios; el enunciado se ha redactado de nuevo y, en algunos casos, se han cambiado los datos. Las soluciones completas, paso a paso, están en el \cref{ap:soluciones}, y las tablas de la distribución normal, de la $t$ de Student y de la $\chi^{2}$ que se necesitan para resolver los ejercicios a mano están en el \cref{ap:tablas}.

\section*{Dónde se usa en el curso}
\begin{center}
\small
\begin{tabularx}{\textwidth}{@{}>{\raggedright\arraybackslash}p{4.6cm}X@{}}
\toprule
Concepto & Módulos del curso que lo emplean \\
\midrule
Media, varianza y tipificación & M05 (momentos), M11 (normalización por capas y RMSNorm), M07 (medias exponenciales de Adam). \\
Estimador, sesgo, varianza, ECM & M05 (estimación de parámetros), M08 (descomposición sesgo--varianza), M22 y M23 (estimador insesgado de pass@$k$). \\
Máxima verosimilitud & M01 (EMV de los bigramas), M05 (EMV de la normal), M08 (mínimos cuadrados y regresión logística como EMV). \\
Teorema central del límite, error estándar & M05 (Monte Carlo), M07 (ruido del gradiente con minilotes), M22 y M23 (error estándar de una exactitud). \\
Intervalos de confianza & M05 (comparación de modelos), M16 (ajuste de leyes de escala), M22 y M23 (intervalos para benchmarks, regla del~3). \\
Contrastes y comparaciones múltiples & M22 y M23 (comparar modelos en muchos benchmarks), M21 (evaluaciones por preferencias). \\
Regresión lineal y $R^{2}$ & M08 (ecuaciones normales), M16 (ajuste log-log de leyes de potencia), M23 (sondas lineales). \\
Correlación frente a causalidad & M23 (intervenciones causales frente a correlaciones en la interpretabilidad). \\
\bottomrule
\end{tabularx}
\end{center}

\section*{Sobre las fuentes}
El texto es de redacción propia. La estadística descriptiva sigue el capítulo~2 de \emph{Introductory Statistics} de OpenStax~\cite{openstax} y el capítulo~10 de Rice~\cite{rice}. La teoría de estimadores, máxima verosimilitud, contrastes y regresión sigue los capítulos~6, 9, 10 y~13 de Wasserman~\cite{wasserman}, las lecturas del curso 18.05 del MIT de Orloff y Bloom~\cite{mit1805} y las clases del curso Stats~200 de Stanford~\cite{stats200}. Los resultados de probabilidad (momentos muestrales, teorema central del límite, distribuciones $\chi^{2}$ y $t$) siguen los capítulos~6, 7 y~10 de Blitzstein y Hwang~\cite{bh}, texto del curso Stat~110 de Harvard.

\section*{Notación}
\begin{center}
\small
\begin{tabularx}{\textwidth}{@{}>{\raggedright\arraybackslash}p{5.2cm}X@{}}
\toprule
Símbolo & Significado \\
\midrule
$x_1,\dots,x_n$; $x_{(1)}\le\dots\le x_{(n)}$ & datos observados y datos ordenados \\
$\xbar$, $s^{2}$, $s$ & media, varianza y desviación típica muestrales (observadas) \\
$X_1,\dots,X_n$ \iid & muestra aleatoria simple (variables independientes e idénticamente distribuidas) \\
$\Xbar$, $S^{2}$ & media y varianza muestrales como variables aleatorias \\
$\mu$, $\sigma^{2}$, $p$, $\lambda$, $\theta$ & parámetros de la población \\
$\widehat\theta$ & estimador (o estimación) de $\theta$ \\
$\PP$, $\E$, $\Var$, $\Cov$ & probabilidad, esperanza, varianza y covarianza \\
$\Nor(\mu,\sigma^{2})$, $\Phi$ & distribución normal y función de distribución de la normal estándar \\
$\Ber(p)$, $\Bin(n,p)$, $\Pois(\lambda)$, $\Expo(\lambda)$ & Bernoulli, binomial, Poisson y exponencial de tasa $\lambda$ \\
$z_{q}$, $t_{k;q}$, $\chi^{2}_{k;q}$ & cuantiles de orden $q$ de $\Nor(0,1)$, de $t_k$ y de $\chi^{2}_k$ \\
$H_0$, $H_1$, $\alpha$, $\beta$ & hipótesis nula y alternativa, nivel de significación y probabilidad de error de tipo~II \\
$\ee(\widehat\theta)$, $\ECM(\widehat\theta)$ & error estándar y error cuadrático medio \\
$\convp$ & convergencia en probabilidad \\
\bottomrule
\end{tabularx}
\end{center}
Los cuantiles se escriben siempre por su probabilidad acumulada: $z_{0{,}975}=1{,}96$ es el valor que deja a su izquierda una probabilidad $0{,}975$.

\tableofcontents

\mainmatter
% =====================================================================
\chapter{Estadística descriptiva}
\label{tema:1}
% =====================================================================

Antes de construir modelos conviene mirar los datos: resumirlos con unos pocos números y dibujarlos. Este tema presenta las medidas de centralización (media, mediana), de posición (cuantiles y cuartiles), de dispersión (rango, recorrido intercuartílico, varianza y desviación típica), las representaciones gráficas básicas (histograma y diagrama de caja) y las medidas de relación entre dos variables (covarianza y correlación). Se sigue el capítulo~2 de \emph{Introductory Statistics} de OpenStax~\cite{openstax} (§2.2 a §2.7), el capítulo~10 de Rice~\cite{rice}, «Summarizing Data», y, para la covarianza y la correlación, la lectura~7b del curso 18.05 del MIT~\cite{mit1805} y la §7.3 de Blitzstein y Hwang~\cite{bh}.

\section{Datos, población y muestra}

\begin{definicion}[Población, muestra y variable]\label{def:t1-poblacion}
La \emph{población} es el conjunto de todos los individuos u objetos sobre los que se quiere obtener información; una \emph{muestra} es un subconjunto de la población que efectivamente se observa. Una \emph{variable} es una característica que se mide en cada individuo. La variable es \emph{cuantitativa} si toma valores numéricos con los que tiene sentido operar (discreta si toma valores aislados, como un número de errores; continua si puede tomar cualquier valor de un intervalo, como un tiempo) y \emph{cualitativa} o \emph{categórica} si toma valores en un conjunto de categorías.
\end{definicion}

Un conjunto de datos cuantitativos es una lista $x_1,\dots,x_n$ de números reales, donde $n$ es el \emph{tamaño} de la muestra. Los mismos datos ordenados de menor a mayor se escriben $x_{(1)}\le x_{(2)}\le\dots\le x_{(n)}$; $x_{(1)}$ es el mínimo y $x_{(n)}$, el máximo. En este tema los datos son números fijos y las definiciones son puramente descriptivas; a partir del \cref{tema:2} se interpretarán como valores observados de variables aleatorias.

\begin{ejemplo}[Datos de referencia]\label{ej:t1-tokens}
Se miden las longitudes, en tokens, de ocho respuestas de un modelo de lenguaje:
\[
12,\quad 15,\quad 11,\quad 18,\quad 14,\quad 40,\quad 13,\quad 16.
\]
Ordenados: $11,12,13,14,15,16,18,40$. Este conjunto se usará como ejemplo a lo largo del tema; el valor $40$ es claramente mayor que los demás.
\end{ejemplo}

\section{Medidas de centralización}

\begin{definicion}[Media, mediana y moda]\label{def:t1-centro}
Sean $x_1,\dots,x_n$ datos reales.
\begin{enumerate}
  \item La \emph{media} (aritmética) es $\displaystyle\xbar=\frac1n\sum_{i=1}^{n}x_i$.
  \item La \emph{mediana} es el valor central de los datos ordenados: $x_{((n+1)/2)}$ si $n$ es impar, y $\tfrac12\bigl(x_{(n/2)}+x_{(n/2+1)}\bigr)$ si $n$ es par.
  \item La \emph{moda} es el valor que aparece con mayor frecuencia (puede no ser única).
\end{enumerate}
\end{definicion}

\begin{ejemplo}\label{ej:t1-media-mediana}
En el \cref{ej:t1-tokens}, $\sum x_i=139$ y $\xbar=139/8=17{,}375$. Como $n=8$ es par, la mediana es la media de los valores cuarto y quinto ordenados: $(14+15)/2=14{,}5$. Si se elimina el valor $40$, la media pasa a $99/7\approx 14{,}14$ y la mediana a $14$: la media se ha desplazado más de tres unidades y la mediana, media unidad.
\end{ejemplo}

El ejemplo ilustra que la media es \emph{sensible a los valores atípicos} y la mediana es \emph{robusta}. Cuando la distribución de los datos es asimétrica con cola larga a la derecha (como las longitudes de respuesta, los tiempos o los ingresos), la media suele ser mayor que la mediana~\cite[§2.6]{openstax}. Las dos medidas tienen una caracterización variacional que explica ese comportamiento.

\begin{proposicion}[Propiedades extremales de la media y la mediana]\label{prop:t1-extremal}
Sean $x_1,\dots,x_n$ datos reales.
\begin{enumerate}
  \item $\sum_{i=1}^{n}(x_i-\xbar)=0$.
  \item Para todo $a\in\R$, $\displaystyle\sum_{i=1}^{n}(x_i-a)^{2}=\sum_{i=1}^{n}(x_i-\xbar)^{2}+n(\xbar-a)^{2}$. En consecuencia, la función $a\mapsto\sum_i(x_i-a)^{2}$ alcanza su mínimo únicamente en $a=\xbar$.
  \item La función $a\mapsto\sum_i\abs{x_i-a}$ alcanza su mínimo en la mediana.
\end{enumerate}
\end{proposicion}

\begin{proof}
(i) $\sum_i(x_i-\xbar)=\sum_i x_i-n\xbar=n\xbar-n\xbar=0$.

(ii) Se escribe $x_i-a=(x_i-\xbar)+(\xbar-a)$ y se desarrolla el cuadrado:
\[
\sum_i(x_i-a)^{2}=\sum_i(x_i-\xbar)^{2}+2(\xbar-a)\sum_i(x_i-\xbar)+n(\xbar-a)^{2}.
\]
El término central es nulo por (i). El último sumando es no negativo y solo se anula si $a=\xbar$.

(iii) Se demuestra en el \cref{eje:t1-mediana-min}.
\end{proof}

\section{Medidas de posición: cuantiles y diagrama de caja}

La mediana divide los datos ordenados en dos mitades. Los cuartiles los dividen en cuatro partes. Hay varias convenciones para calcularlos que dan resultados ligeramente distintos con pocos datos (los programas estadísticos suelen interpolar); en este módulo se usa la siguiente, que es la de~\cite[§2.3]{openstax} y la habitual en el cálculo a mano.

\begin{definicion}[Cuartiles y recorrido intercuartílico]\label{def:t1-cuartiles}
Sean $x_{(1)}\le\dots\le x_{(n)}$ los datos ordenados. La \emph{mitad inferior} está formada por los $\lfloor n/2\rfloor$ primeros datos y la \emph{mitad superior}, por los $\lfloor n/2\rfloor$ últimos (si $n$ es impar, la mediana no pertenece a ninguna de las dos). El \emph{primer cuartil} $Q_1$ es la mediana de la mitad inferior; el \emph{segundo cuartil} $Q_2$ es la mediana; el \emph{tercer cuartil} $Q_3$ es la mediana de la mitad superior. El \emph{recorrido intercuartílico} es $\RIC=Q_3-Q_1$ y el \emph{rango} es $x_{(n)}-x_{(1)}$.
\end{definicion}

Más en general, el \emph{cuantil} de orden $q\in(0,1)$ (o \emph{percentil} $100q$) es un valor que deja a su izquierda aproximadamente una fracción $q$ de los datos; $Q_1$, $Q_2$ y $Q_3$ son los cuantiles de órdenes $0{,}25$, $0{,}5$ y $0{,}75$.

\begin{definicion}[Resumen de cinco números y diagrama de caja]\label{def:t1-caja}
El \emph{resumen de cinco números} es $(x_{(1)},Q_1,Q_2,Q_3,x_{(n)})$. Un dato se considera \emph{atípico} (criterio de Tukey) si está fuera del intervalo
\[
\bigl[\,Q_1-1{,}5\,\RIC,\;Q_3+1{,}5\,\RIC\,\bigr],
\]
cuyos extremos se llaman \emph{vallas}. El \emph{diagrama de caja} dibuja una caja de $Q_1$ a $Q_3$ con una línea en la mediana y dos «bigotes» que llegan hasta el dato más extremo no atípico de cada lado; los atípicos se marcan por separado~\cite[§2.4]{openstax}.
\end{definicion}

\begin{ejemplo}\label{ej:t1-cuartiles}
En el \cref{ej:t1-tokens}, $n=8$, la mitad inferior es $11,12,13,14$ y la superior $15,16,18,40$. Entonces $Q_1=(12+13)/2=12{,}5$, $Q_2=14{,}5$, $Q_3=(16+18)/2=17$ y $\RIC=4{,}5$. Las vallas son $12{,}5-6{,}75=5{,}75$ y $17+6{,}75=23{,}75$, de modo que $40$ es atípico. El bigote superior llega hasta $18$.
\end{ejemplo}

\begin{figure}[ht]
\centering
\begin{tikzpicture}[x=0.075cm,y=0.33cm]
  % Histograma
  \foreach \a/\h in {30/1,40/3,50/3,60/7,70/4,80/5,90/8}{
    \draw[fill=gray!25] (\a,0) rectangle (\a+10,\h);
  }
  \draw[->] (28,0) -- (106,0) node[right]{\small nota};
  \draw[->] (28,0) -- (28,9.2) node[above]{\small frecuencia};
  \foreach \a in {30,40,...,100}{\draw (\a,0) -- (\a,-0.25) node[below,text height=1.4ex,text depth=0.2ex]{\scriptsize\a};}
  \foreach \h in {2,4,6,8}{\draw (28,\h) -- (27.2,\h) node[left]{\scriptsize\h};}
  % Diagrama de caja
  \begin{scope}[yshift=-2.3cm]
    \draw (33,0) -- (61,0); \draw (90,0) -- (100,0);
    \draw (33,-0.6) -- (33,0.6); \draw (100,-0.6) -- (100,0.6);
    \draw[fill=gray!15] (61,-1) rectangle (90,1);
    \draw[thick] (73,-1) -- (73,1);
    \foreach \a in {30,40,...,100}{\draw (\a,-1.6) -- (\a,-1.85) node[below,text height=1.4ex,text depth=0.2ex]{\scriptsize\a};}
    \draw (28,-1.6) -- (106,-1.6);
    \node[above] at (61,1) {\scriptsize $Q_1$};
    \node[above] at (73,1) {\scriptsize $Q_2$};
    \node[above] at (90,1) {\scriptsize $Q_3$};
  \end{scope}
\end{tikzpicture}
\caption{Histograma (clases de amplitud 10) y diagrama de caja de las 31 notas de examen del ejemplo~2.33 de~\cite{openstax}: mínimo 33, $Q_1=61$, mediana 73, $Q_3=90$, máximo 100. La cola izquierda más larga se aprecia en ambos gráficos.}\label{fig:t1-hist-caja}
\end{figure}

\begin{definicion}[Histograma]\label{def:t1-histograma}
Se divide un intervalo que contiene los datos en \emph{clases} contiguas $[a_0,a_1),[a_1,a_2),\dots,[a_{k-1},a_k]$ y se cuenta la \emph{frecuencia absoluta} $f_j$ de datos en cada clase. El \emph{histograma} dibuja sobre cada clase un rectángulo de altura proporcional a $f_j$ (si las clases tienen igual amplitud) o a $f_j/(a_j-a_{j-1})$ (si no la tienen, de modo que el \emph{área} sea proporcional a la frecuencia). La \emph{frecuencia relativa} es $f_j/n$.
\end{definicion}

El histograma muestra la forma de la distribución (simetría, asimetría, número de modas), mientras que el diagrama de caja resume la posición y la dispersión y señala los atípicos; la \cref{fig:t1-hist-caja} muestra ambos para un mismo conjunto de datos.

\section{Medidas de dispersión}

\begin{definicion}[Varianza y desviación típica muestrales]\label{def:t1-varianza}
La \emph{varianza muestral} de $x_1,\dots,x_n$ ($n\ge 2$) es
\[
s^{2}=\frac{1}{n-1}\sum_{i=1}^{n}(x_i-\xbar)^{2},
\]
y la \emph{desviación típica muestral} es $s=\sqrt{s^{2}}$. La cantidad $\sum_i(x_i-\xbar)^{2}$ se llama \emph{suma de cuadrados} de los datos.
\end{definicion}

\begin{observacion}[¿Por qué $n-1$?]\label{obs:t1-n-1}
El denominador $n-1$ en lugar de $n$ tiene dos justificaciones. La primera es algebraica: las $n$ desviaciones $x_i-\xbar$ suman cero (\cref{prop:t1-extremal}), de modo que solo $n-1$ de ellas pueden elegirse libremente; se dice que la suma de cuadrados tiene $n-1$ \emph{grados de libertad}. La segunda es estadística: si los datos son una muestra aleatoria de una población con varianza $\sigma^{2}$, la versión con $n-1$ acierta «en promedio», es decir, es un estimador insesgado de $\sigma^{2}$, mientras que la versión con $n$ la subestima sistemáticamente. Esto se demuestra en el \cref{tema:2} (\cref{teo:t2-s2}). Con denominador $n$ se obtiene la \emph{varianza poblacional} de los propios datos, $\frac1n\sum_i(x_i-\xbar)^{2}$, que es la que se usa cuando los datos \emph{son} toda la población o cuando se normaliza un vector, como en la normalización por capas de una red (módulo~11).
\end{observacion}

\begin{proposicion}[Fórmula de cálculo]\label{prop:t1-abreviada}
$\displaystyle\sum_{i=1}^{n}(x_i-\xbar)^{2}=\sum_{i=1}^{n}x_i^{2}-n\xbar^{2}$, y por tanto $s^{2}=\dfrac{\sum_i x_i^{2}-n\xbar^{2}}{n-1}$.
\end{proposicion}

\begin{proof}
Es el caso $a=0$ de la \cref{prop:t1-extremal}(ii): $\sum_i x_i^{2}=\sum_i(x_i-\xbar)^{2}+n\xbar^{2}$.
\end{proof}

\begin{ejemplo}\label{ej:t1-varianza}
En el \cref{ej:t1-tokens}, $\sum x_i^{2}=3035$ y $n\xbar^{2}=8\cdot 17{,}375^{2}=2415{,}125$. La suma de cuadrados es $3035-2415{,}125=619{,}875$, de donde $s^{2}=619{,}875/7\approx 88{,}55$ y $s\approx 9{,}41$. Sin el dato $40$ se obtiene $s\approx 2{,}41$: la desviación típica, como la media, es muy sensible a los atípicos, mientras que el $\RIC$ apenas cambia (de $4{,}5$ a $4$).
\end{ejemplo}

\begin{proposicion}[Cambios de escala y de origen]\label{prop:t1-lineal}
Si $y_i=ax_i+b$ con $a,b\in\R$, entonces $\ybar=a\xbar+b$, $s_y^{2}=a^{2}s_x^{2}$ y $s_y=\abs{a}\,s_x$. Si $a>0$, la mediana y los cuartiles se transforman igual que la media y $\RIC_y=a\,\RIC_x$.
\end{proposicion}

\begin{proof}
$\ybar=\frac1n\sum_i(ax_i+b)=a\xbar+b$. Entonces $y_i-\ybar=a(x_i-\xbar)$, y elevando al cuadrado y sumando, $\sum_i(y_i-\ybar)^{2}=a^{2}\sum_i(x_i-\xbar)^{2}$. Si $a>0$, la transformación conserva el orden, luego transforma los datos ordenados en datos ordenados y las medianas de cualquier subconjunto se transforman por $t\mapsto at+b$.
\end{proof}

\begin{definicion}[Tipificación]\label{def:t1-z}
Si $s>0$, la \emph{puntuación tipificada} (o puntuación $z$) del dato $x_i$ es $z_i=(x_i-\xbar)/s$. Por la \cref{prop:t1-lineal}, los datos tipificados tienen media $0$ y desviación típica $1$; $z_i$ mide a cuántas desviaciones típicas está $x_i$ de la media.
\end{definicion}

\section{Datos bivariantes: covarianza y correlación}

Cuando en cada individuo se miden dos variables se obtienen pares $(x_1,y_1),\dots,(x_n,y_n)$, que se representan en un \emph{diagrama de dispersión}.

\begin{definicion}[Covarianza y correlación muestrales]\label{def:t1-cov}
La \emph{covarianza muestral} y el \emph{coeficiente de correlación} (de Pearson) son
\[
s_{xy}=\frac{1}{n-1}\sum_{i=1}^{n}(x_i-\xbar)(y_i-\ybar),
\qquad
r=\frac{s_{xy}}{s_x\,s_y}=\frac{\sum_i(x_i-\xbar)(y_i-\ybar)}{\sqrt{\sum_i(x_i-\xbar)^{2}}\,\sqrt{\sum_i(y_i-\ybar)^{2}}},
\]
donde se supone $s_x,s_y>0$.
\end{definicion}

El signo de cada producto $(x_i-\xbar)(y_i-\ybar)$ es positivo cuando el punto está en el cuadrante superior derecho o inferior izquierdo respecto de $(\xbar,\ybar)$ y negativo en los otros dos; $s_{xy}>0$ indica que predominan los primeros, es decir, una tendencia creciente. La covarianza depende de las unidades; la correlación no (\cref{eje:t1-r-invariante}).

\begin{teorema}[Acotación de la correlación]\label{teo:t1-r}
Se cumple $-1\le r\le 1$. Además, $\abs{r}=1$ si y solo si existen $a\neq 0$ y $b$ tales que $y_i=ax_i+b$ para todo $i$, y entonces el signo de $r$ es el de $a$.
\end{teorema}

\begin{proof}
Sean $u_i=x_i-\xbar$ y $v_i=y_i-\ybar$, y $\mathbf u,\mathbf v\in\R^{n}$ los vectores correspondientes, ambos no nulos porque $s_x,s_y>0$. Entonces $r=\langle\mathbf u,\mathbf v\rangle/(\norm{\mathbf u}\norm{\mathbf v})$, y la desigualdad de Cauchy--Schwarz, $\abs{\langle\mathbf u,\mathbf v\rangle}\le\norm{\mathbf u}\norm{\mathbf v}$, da $\abs r\le 1$. La igualdad en Cauchy--Schwarz se da exactamente cuando $\mathbf v=a\mathbf u$ para algún $a\neq0$, es decir, $y_i-\ybar=a(x_i-\xbar)$ para todo $i$, que equivale a $y_i=ax_i+b$ con $b=\ybar-a\xbar$. En ese caso $\langle\mathbf u,\mathbf v\rangle=a\norm{\mathbf u}^{2}$ tiene el signo de $a$. Recíprocamente, si $y_i=ax_i+b$ entonces $v_i=au_i$ y $r=a\norm{\mathbf u}^{2}/(\norm{\mathbf u}\,\abs a\norm{\mathbf u})=a/\abs a=\pm1$.
\end{proof}

\begin{ejemplo}\label{ej:t1-r}
Para los pares $(1,2),(2,4),(3,5),(4,4),(5,5)$: $\xbar=3$, $\ybar=4$, $\sum(x_i-\xbar)^{2}=10$, $\sum(y_i-\ybar)^{2}=6$ y $\sum(x_i-\xbar)(y_i-\ybar)=(-2)(-2)+(-1)(0)+0\cdot1+1\cdot0+2\cdot1=6$. Entonces $s_{xy}=6/4=1{,}5$ y $r=6/\sqrt{60}\approx 0{,}775$: asociación lineal positiva moderadamente fuerte.
\end{ejemplo}

Los mismos conceptos existen para variables aleatorias. Si $(X,Y)$ es un vector aleatorio con varianzas finitas y positivas, su \emph{covarianza} y su \emph{correlación} son
\[
\Cov(X,Y)=\E\bigl[(X-\E X)(Y-\E Y)\bigr]=\E[XY]-\E[X]\E[Y],
\qquad
\Cor(X,Y)=\frac{\Cov(X,Y)}{\sqrt{\Var X\,\Var Y}},
\]
y también $\abs{\Cor(X,Y)}\le1$~\cite[§7.3]{bh}. Las fórmulas muestrales de la \cref{def:t1-cov} son las análogas con los datos (salvo el factor $n/(n-1)$, que se cancela en $r$).

\begin{observacion}[Errores frecuentes]\label{obs:t1-errores}
(a) Una correlación cercana a cero no significa que no haya relación, sino que no hay relación \emph{lineal}: para los puntos $(-2,4),(-1,1),(0,0),(1,1),(2,4)$, en los que $y=x^{2}$, se tiene $r=0$ (\cref{eje:t6-cuadratica}).
(b) Una correlación alta no implica causalidad (\cref{tema:6}).
(c) Calcular $s^{2}$ con denominador $n$ cuando se pide la varianza muestral, o confundir $s^{2}$ con $s$.
(d) Olvidar ordenar los datos antes de calcular la mediana o los cuartiles.
(e) Comparar medias de grupos muy asimétricos o con atípicos sin mirar también la mediana.
\end{observacion}

\begin{observacion}[Conexión con el curso]
La normalización por capas (\emph{LayerNorm}) del módulo~11 y de los Transformers del módulo~14 tipifica el vector de activaciones de cada token: le resta su media y lo divide por su desviación típica (con denominador $n$, como en la \cref{obs:t1-n-1}). Las tablas de resultados de los benchmarks de los módulos~22 y~23 son, en esencia, medias de variables $0$/$1$ (acierto o fallo), y los diagramas de caja son la forma habitual de mostrar la variabilidad entre semillas o entre ejecuciones de un mismo entrenamiento.
\end{observacion}

\begin{lecturas}
\item OpenStax~\cite{openstax}, cap.~2: §2.2 «Histograms, Frequency Polygons, and Time Series Graphs», §2.3 «Measures of the Location of the Data», §2.4 «Box Plots», §2.5 «Measures of the Center of the Data», §2.6 «Skewness and the Mean, Median, and Mode» y §2.7 «Measures of the Spread of the Data».
\item Rice~\cite{rice}, cap.~10 «Summarizing Data».
\item Orloff y Bloom~\cite{mit1805}, lectura~7b «Covariance and Correlation».
\item Blitzstein y Hwang~\cite{bh}, §7.3 «Covariance and correlation».
\end{lecturas}

\section*{Ejercicios del tema 1}
\addcontentsline{toc}{section}{Ejercicios del tema 1}

\begin{ejercicio}{1}{eje:t1-vf}
Indique si cada afirmación es verdadera o falsa y justifíquelo brevemente.
(a) Si se multiplica por~10 el mayor dato de un conjunto, la mediana no cambia.
(b) La desviación típica de un conjunto de datos puede ser negativa.
(c) Si todos los datos son iguales, $s=0$ y $\RIC=0$.
(d) Si $r=0{,}9$, los puntos están exactamente sobre una recta.
(e) Sumar $5$ a todos los datos no cambia ni la varianza ni el $\RIC$.
\end{ejercicio}

\begin{ejercicio}{2}{eje:t1-basico}
Para los datos $4,\,7,\,2,\,9,\,5,\,6,\,3,\,25,\,6$ (número de llamadas a herramientas de un agente en nueve tareas), calcule la media, la mediana, la moda, los cuartiles, el $\RIC$, las vallas de Tukey, la varianza y la desviación típica muestrales. ¿Hay atípicos? Repita el cálculo de la media, la mediana y $s$ sin el atípico y comente el resultado.
\end{ejercicio}

\begin{ejercicio}{2}{eje:t1-edades}
Las edades (en años) de veinte alumnos de una clase son
$9;\allowbreak\ 9{,}5;\allowbreak\ 9{,}5;\allowbreak\ 10;\allowbreak\ 10;\allowbreak\ 10;\allowbreak\ 10;\allowbreak\ 10{,}5;\allowbreak\ 10{,}5;\allowbreak\ 10{,}5;\allowbreak\ 10{,}5;\allowbreak\ 11;\allowbreak\ 11;\allowbreak\ 11;\allowbreak\ 11;\allowbreak\ 11;\allowbreak\ 11;\allowbreak\ 11{,}5;\allowbreak\ 11{,}5;\allowbreak\ 11{,}5$.
Calcule la media, la mediana, los cuartiles y la desviación típica muestral. ¿Qué fracción de los datos está a menos de una desviación típica de la media?
\fuente{adaptado de OpenStax, \emph{Introductory Statistics}~2e~\cite{openstax}, §2.7, ejemplo~2.32.}
\end{ejercicio}

\begin{ejercicio}{2}{eje:t1-cajas}
Los resúmenes de cinco números de las notas de un examen en un grupo de mañana y otro de tarde son $(32;\allowbreak\,56;\allowbreak\,74{,}5;\allowbreak\,82{,}5;\allowbreak\,99)$ y $(25{,}5;\allowbreak\,78;\allowbreak\,81;\allowbreak\,89;\allowbreak\,98)$. (a) Calcule el rango y el $\RIC$ de cada grupo. (b) ¿Qué grupo tiene la mitad central de los datos más dispersa? (c) Calcule las vallas de Tukey y determine si el mínimo de cada grupo es atípico. (d) Esboce los dos diagramas de caja sobre un mismo eje.
\fuente{adaptado de OpenStax, \emph{Introductory Statistics}~2e~\cite{openstax}, §2.4, ejemplo~2.24.}
\end{ejercicio}

\begin{ejercicio}{2}{eje:t1-agrupados}
Una tabla de frecuencias agrupa 26 datos en clases con marcas de clase (puntos medios) $1,4,7,10,13,16$ y frecuencias $1,6,10,7,0,2$, respectivamente. Aproxime la media y la desviación típica muestral sustituyendo cada dato por la marca de su clase.
\fuente{adaptado de OpenStax, \emph{Introductory Statistics}~2e~\cite{openstax}, §2.7, ejemplo~2.34.}
\end{ejercicio}

\begin{ejercicio}{2}{eje:t1-abreviada}
De un conjunto de $n=10$ datos se conocen $\sum x_i=50$ y $\sum x_i^{2}=290$. (a) Calcule $\xbar$ y $s^{2}$. (b) Se añade un undécimo dato igual a $5$. Calcule la nueva media y la nueva varianza muestral sin recalcular desde cero.
\end{ejercicio}

\begin{ejercicio}{2}{eje:t1-celsius}
Las temperaturas máximas de una semana, en grados Celsius, son $18,21,25,19,22,30,24$. (a) Calcule $\xbar$, la mediana, $s$ y el $\RIC$. (b) Sin repetir los cálculos, obtenga esas mismas medidas en grados Fahrenheit ($F=1{,}8\,C+32$). (c) ¿Cambian las puntuaciones $z$?
\end{ejercicio}

\begin{ejercicio}{2}{eje:t1-zscore}
Juan obtiene un $2{,}85$ en una escala en la que la media de su centro es $3{,}0$ y la desviación típica $0{,}7$. Alí obtiene un $77$ en otra escala cuya media es $80$ y cuya desviación típica es $10$. ¿Quién tiene una nota mejor en relación con su centro?
\fuente{adaptado de OpenStax, \emph{Introductory Statistics}~2e~\cite{openstax}, §2.7, ejemplo~2.35.}
\end{ejercicio}

\begin{ejercicio}{2}{eje:t1-cov-tabla}
Las variables $X$ e $Y$ toman los valores $1$ y $-1$, cada una con probabilidad $1/2$, y $c=\PP(X=1,Y=1)$. (a) Escriba la distribución conjunta en función de $c$ y calcule $\Cov(X,Y)$ y $\Cor(X,Y)$. (b) ¿Para qué valores de $c$ son independientes? ¿Para cuáles es $\Cor(X,Y)=1$? ¿Y $-1$?
\fuente{adaptado de Orloff y Bloom, MIT 18.05~\cite{mit1805}, \emph{problem set}~5, problema~3.}
\end{ejercicio}

\begin{ejercicio}{2}{eje:t1-agrupar}
Un grupo de $10$ alumnos tiene media $70$ y desviación típica muestral $8$ en un examen; otro de $15$ alumnos tiene media $80$ y desviación típica $6$. Calcule la media y la desviación típica muestral de los $25$ alumnos juntos.
\end{ejercicio}

\begin{ejercicio}{3}{eje:t1-mediana-min}
Demuestre que la función $g(a)=\sum_{i=1}^{n}\abs{x_i-a}$ alcanza su mínimo en la mediana. (\emph{Indicación:} estudie cómo cambia $g$ al mover $a$ en un intervalo entre dos datos ordenados consecutivos.)
\end{ejercicio}

\begin{ejercicio}{3}{eje:t1-r-invariante}
Sean $x_i'=ax_i+b$ e $y_i'=cy_i+d$ con $a,c\neq0$. Demuestre que $s_{x'y'}=ac\,s_{xy}$ y que $r_{x'y'}=\operatorname{signo}(ac)\,r_{xy}$. Deduzca que la correlación no depende de las unidades de medida.
\end{ejercicio}

\begin{ejercicio}{3}{eje:t1-chebyshev}
Sean $x_1,\dots,x_n$ datos con desviación típica muestral $s>0$ y sea $k>0$. Demuestre que la fracción de datos con $\abs{x_i-\xbar}\ge ks$ es a lo sumo $\dfrac{n-1}{n k^{2}}<\dfrac{1}{k^{2}}$. ¿Qué garantiza esta cota para $k=2$?
\end{ejercicio}

\begin{ejercicio}{3}{eje:t1-welford}
Sean $\xbar_m$ y $M_m=\sum_{i=1}^{m}(x_i-\xbar_m)^{2}$ la media y la suma de cuadrados de los $m$ primeros datos. Demuestre las fórmulas de actualización
\[
\xbar_m=\xbar_{m-1}+\frac{x_m-\xbar_{m-1}}{m},\qquad M_m=M_{m-1}+(x_m-\xbar_{m-1})(x_m-\xbar_m),
\]
que permiten calcular la media y la varianza en una sola pasada sobre los datos. Compruébelas con los datos $2,4,9$.
\end{ejercicio}

\begin{ejercicio}{3}{eje:t1-varianza-pares}
Demuestre la identidad
\[
s^{2}=\frac{1}{n(n-1)}\sum_{1\le i<j\le n}(x_i-x_j)^{2},
\]
que expresa la varianza muestral sin usar la media: es la mitad del promedio de los cuadrados de las diferencias entre todos los pares de datos distintos.
\end{ejercicio}

% =====================================================================
\chapter{Muestreo y estimadores}
\label{tema:2}
% =====================================================================

La estadística inferencial trata los datos como valores observados de variables aleatorias cuya distribución depende de parámetros desconocidos. Este tema introduce el modelo de muestra aleatoria simple, la noción de estimador y los criterios para juzgar su calidad: sesgo, varianza, error estándar, error cuadrático medio y consistencia. Se demuestran los dos resultados centrales del muestreo: la media muestral es un estimador insesgado de la media y la varianza muestral con denominador $n-1$ lo es de la varianza. Se sigue el capítulo~6 de Wasserman~\cite{wasserman}, «Models, Statistical Inference and Learning», la §6.3 de Blitzstein y Hwang~\cite{bh}, «Sample moments», la clase~12 del curso Stats~200 de Stanford~\cite{stats200} y la lectura~10a del curso 18.05 del MIT~\cite{mit1805}.

\section{Recordatorio: esperanza y varianza}

Se recogen las propiedades de la esperanza y la varianza de variables aleatorias (discretas o continuas) que se usarán continuamente~\cite[cap.~4 y §7.3]{bh}. Todas las variables que aparecen se suponen con varianza finita.

\begin{proposicion}[Propiedades de la esperanza y la varianza]\label{prop:t2-props}
Sean $X,Y,X_1,\dots,X_n$ variables aleatorias y $a,b\in\R$.
\begin{enumerate}
  \item Linealidad: $\E[aX+bY]=a\E[X]+b\E[Y]$ y, en general, $\E\bigl[\sum_iX_i\bigr]=\sum_i\E[X_i]$.
  \item $\Var(X)=\E[(X-\E X)^{2}]=\E[X^{2}]-(\E X)^{2}$, y $\Var(aX+b)=a^{2}\Var(X)$.
  \item $\Var\bigl(\sum_iX_i\bigr)=\sum_i\Var(X_i)+2\sum_{i<j}\Cov(X_i,X_j)$.
  \item Si $X$ e $Y$ son independientes, $\E[XY]=\E[X]\E[Y]$ y, por tanto, $\Cov(X,Y)=0$. En consecuencia, si $X_1,\dots,X_n$ son independientes, $\Var\bigl(\sum_iX_i\bigr)=\sum_i\Var(X_i)$.
  \item La covarianza es bilineal,
  \[
  \Cov\Bigl(\sum_ia_iX_i,\sum_jb_jY_j\Bigr)=\sum_{i,j}a_ib_j\Cov(X_i,Y_j),
  \]
  y $\Cov(X,X)=\Var(X)$.
\end{enumerate}
\end{proposicion}

\begin{proof}
(i) y la independencia en (iv) son propiedades de la esperanza que se suponen conocidas~\cite[§4.2 y §7.3]{bh}. (ii) Desarrollando el cuadrado y usando (i): $\E[(X-\mu)^{2}]=\E[X^{2}]-2\mu\E[X]+\mu^{2}=\E[X^{2}]-\mu^{2}$ con $\mu=\E X$; además $aX+b-\E[aX+b]=a(X-\mu)$. (v) Se sigue de la linealidad de la esperanza aplicada a $\bigl(\sum_ia_i(X_i-\E X_i)\bigr)\bigl(\sum_jb_j(Y_j-\E Y_j)\bigr)$. (iii) Es (v) con $Y_j=X_j$ y todos los coeficientes iguales a~1, separando los términos $i=j$ de los $i\neq j$. La última afirmación de (iv) se sigue de (iii).
\end{proof}

\section{Muestra aleatoria, parámetro y estadístico}

\begin{definicion}[Muestra aleatoria simple]\label{def:t2-mas}
Una \emph{muestra aleatoria simple} de tamaño $n$ de una distribución $F$ es una colección de variables aleatorias $X_1,\dots,X_n$ independientes y todas con distribución $F$; se dice que son \emph{independientes e idénticamente distribuidas} (\iid) y se escribe $X_1,\dots,X_n\sim F$ \iid\ Los datos $x_1,\dots,x_n$ son una realización (valores observados) de la muestra.
\end{definicion}

El modelo describe el muestreo con reemplazamiento de una población, o el muestreo sin reemplazamiento cuando la población es mucho mayor que la muestra, y también la repetición independiente de un experimento: lanzar una moneda, generar una respuesta de un modelo de lenguaje con una semilla distinta o elegir al azar una pregunta de un banco muy grande.

\begin{definicion}[Modelo paramétrico, parámetro y estadístico]\label{def:t2-parametro}
Un \emph{modelo paramétrico} es una familia de distribuciones $\{F_\theta:\theta\in\Theta\}$ indexada por un \emph{parámetro} $\theta$ (un número o un vector) que toma valores en el \emph{espacio paramétrico} $\Theta$. Un \emph{estadístico} es cualquier función $T=g(X_1,\dots,X_n)$ de la muestra que no depende de parámetros desconocidos; es, por tanto, una variable aleatoria, y su distribución se llama \emph{distribución muestral} de $T$.
\end{definicion}

Por ejemplo, $\{\Ber(p):p\in[0,1]\}$, $\{\Pois(\lambda):\lambda>0\}$ y $\{\Nor(\mu,\sigma^{2}):\mu\in\R,\sigma>0\}$ son modelos paramétricos. Los estadísticos más importantes son la \emph{media muestral} y la \emph{varianza muestral}
\[
\Xbar=\frac1n\sum_{i=1}^{n}X_i,\qquad S^{2}=\frac{1}{n-1}\sum_{i=1}^{n}(X_i-\Xbar)^{2},
\]
que son las cantidades del \cref{tema:1} vistas como variables aleatorias. Se reserva la mayúscula para la variable aleatoria y la minúscula ($\xbar$, $s^{2}$) para su valor observado.

\section{Estimadores y sus propiedades}

\begin{definicion}[Estimador, sesgo, error estándar y error cuadrático medio]\label{def:t2-estimador}
Un \emph{estimador} de un parámetro $\theta$ es un estadístico $\widehat\theta=\widehat\theta(X_1,\dots,X_n)$ que se usa para aproximar $\theta$; su valor con unos datos concretos es una \emph{estimación}. Se definen:
\begin{enumerate}
  \item el \emph{sesgo}: $\sesgo(\widehat\theta)=\E_\theta[\widehat\theta]-\theta$; el estimador es \emph{insesgado} si su sesgo es nulo para todo $\theta\in\Theta$;
  \item el \emph{error estándar}: $\ee(\widehat\theta)=\sqrt{\Var_\theta(\widehat\theta)}$;
  \item el \emph{error cuadrático medio}: $\ECM(\widehat\theta)=\E_\theta\bigl[(\widehat\theta-\theta)^{2}\bigr]$.
\end{enumerate}
El subíndice $\theta$ indica que las esperanzas se calculan con la distribución $F_\theta$.
\end{definicion}

El sesgo mide el error sistemático y el error estándar, la dispersión del estimador alrededor de su propia media. El error estándar suele depender de $\theta$, que es desconocido; al sustituir $\theta$ por su estimación se obtiene el \emph{error estándar estimado}, que se escribe $\widehat{\ee}$.

\begin{proposicion}[Descomposición sesgo--varianza]\label{prop:t2-ecm}
$\ECM(\widehat\theta)=\sesgo(\widehat\theta)^{2}+\Var(\widehat\theta)$.
\end{proposicion}

\begin{proof}
Sea $m=\E[\widehat\theta]$. Se escribe $\widehat\theta-\theta=(\widehat\theta-m)+(m-\theta)$, donde el segundo sumando es una constante. Entonces
\[
\E[(\widehat\theta-\theta)^{2}]=\E[(\widehat\theta-m)^{2}]+2(m-\theta)\E[\widehat\theta-m]+(m-\theta)^{2}=\Var(\widehat\theta)+0+\sesgo(\widehat\theta)^{2},
\]
porque $\E[\widehat\theta-m]=0$.
\end{proof}

Un estimador ligeramente sesgado puede tener menor $\ECM$ que uno insesgado si su varianza es bastante menor (\cref{eje:t2-ecm-normal}); este compromiso entre sesgo y varianza reaparece en el aprendizaje automático (módulo~8).

\section{La media muestral}

\begin{teorema}[Media y varianza de la media muestral]\label{teo:t2-media}
Sean $X_1,\dots,X_n$ \iid\ con media $\mu$ y varianza $\sigma^{2}$. Entonces
\[
\E[\Xbar]=\mu,\qquad \Var(\Xbar)=\frac{\sigma^{2}}{n},\qquad \ee(\Xbar)=\frac{\sigma}{\sqrt n}.
\]
En particular, $\Xbar$ es un estimador insesgado de $\mu$.
\end{teorema}

\begin{proof}
Por la linealidad de la esperanza, $\E[\Xbar]=\frac1n\sum_i\E[X_i]=\frac1n\cdot n\mu=\mu$. Por la \cref{prop:t2-props}(ii) y (iv), usando la independencia,
\[
\Var(\Xbar)=\frac{1}{n^{2}}\Var\Bigl(\sum_{i=1}^{n}X_i\Bigr)=\frac{1}{n^{2}}\sum_{i=1}^{n}\Var(X_i)=\frac{n\sigma^{2}}{n^{2}}=\frac{\sigma^{2}}{n}.\qedhere
\]
\end{proof}

El error estándar de la media decrece como $1/\sqrt n$: para dividirlo por dos hay que multiplicar el tamaño de la muestra por cuatro. Esta ley de la raíz cuadrada gobierna el coste de casi toda estimación por muestreo, desde las encuestas hasta la evaluación de modelos.

\begin{ejemplo}[Proporción muestral]\label{ej:t2-proporcion}
Si $X_1,\dots,X_n\sim\Ber(p)$ \iid\ (por ejemplo, $X_i=1$ si un modelo responde bien a la pregunta $i$), entonces $\mu=p$ y $\sigma^{2}=p(1-p)$. La \emph{proporción muestral} $\widehat p=\Xbar$ es la fracción de aciertos; es insesgada y $\ee(\widehat p)=\sqrt{p(1-p)/n}$, que se estima por $\widehat{\ee}=\sqrt{\widehat p(1-\widehat p)/n}$. Con $n=1000$ preguntas y $\widehat p=0{,}8$, $\widehat{\ee}=\sqrt{0{,}16/1000}\approx0{,}0126$, es decir, algo más de un punto porcentual.
\end{ejemplo}

\section{La varianza muestral}

\begin{teorema}[Insesgadez de la varianza muestral]\label{teo:t2-s2}
Sean $X_1,\dots,X_n$ \iid\ con media $\mu$ y varianza $\sigma^{2}$, con $n\ge2$. Entonces $\E[S^{2}]=\sigma^{2}$.
\end{teorema}

\begin{proof}
\emph{Paso 1.} Por la \cref{prop:t1-extremal}(ii) con $a=\mu$, aplicada a los valores $X_1,\dots,X_n$,
\[
\sum_{i=1}^{n}(X_i-\Xbar)^{2}=\sum_{i=1}^{n}(X_i-\mu)^{2}-n(\Xbar-\mu)^{2}.
\]
\emph{Paso 2.} Se toman esperanzas. Como $\E[(X_i-\mu)^{2}]=\Var(X_i)=\sigma^{2}$ y, por el \cref{teo:t2-media}, $\E[(\Xbar-\mu)^{2}]=\Var(\Xbar)=\sigma^{2}/n$,
\[
\E\Bigl[\sum_{i=1}^{n}(X_i-\Xbar)^{2}\Bigr]=n\sigma^{2}-n\cdot\frac{\sigma^{2}}{n}=(n-1)\sigma^{2}.
\]
\emph{Paso 3.} Dividiendo por $n-1$ se obtiene $\E[S^{2}]=\sigma^{2}$.
\end{proof}

La demostración explica el $n-1$: las desviaciones se miden respecto de $\Xbar$, que está más cerca de los datos que la verdadera media $\mu$ (es la que minimiza la suma de cuadrados), de modo que $\sum_i(X_i-\Xbar)^{2}$ es en promedio menor que $\sum_i(X_i-\mu)^{2}$, exactamente en una cantidad $\sigma^{2}$.

\begin{corolario}[Sesgo de la varianza con denominador $n$]\label{cor:t2-sesgo}
El estimador $\widehat\sigma^{2}=\frac1n\sum_i(X_i-\Xbar)^{2}=\frac{n-1}{n}S^{2}$ cumple $\E[\widehat\sigma^{2}]=\frac{n-1}{n}\sigma^{2}$, luego tiene sesgo $-\sigma^{2}/n$: subestima la varianza, aunque el sesgo tiende a cero cuando $n\to\infty$.
\end{corolario}

\begin{observacion}
Que $S^{2}$ sea insesgado para $\sigma^{2}$ no implica que $S$ lo sea para $\sigma$: en general $\E[S]<\sigma$ (\cref{eje:t2-s-sesgo}). La insesgadez no se conserva por transformaciones no lineales.
\end{observacion}

\section{Consistencia y ley de los grandes números}

Un estimador razonable debe acercarse al parámetro cuando crece el tamaño de la muestra. Se escribe $\widehat\theta_n$ para subrayar la dependencia de $n$.

\begin{definicion}[Consistencia]\label{def:t2-consistencia}
Una sucesión de variables aleatorias $Y_n$ \emph{converge en probabilidad} a una constante $c$, y se escribe $Y_n\convp c$, si para todo $\varepsilon>0$
\[
\lim_{n\to\infty}\PP\bigl(\abs{Y_n-c}>\varepsilon\bigr)=0.
\]
Un estimador $\widehat\theta_n$ es \emph{consistente} si $\widehat\theta_n\convp\theta$ para todo $\theta\in\Theta$.
\end{definicion}

\begin{lema}[Desigualdades de Markov y de Chebyshev]\label{lem:t2-chebyshev}
Si $Y\ge0$ y $a>0$, entonces $\PP(Y\ge a)\le\E[Y]/a$. En consecuencia, si $W$ tiene media $\mu_W$ y varianza finita, para todo $\varepsilon>0$
\[
\PP(\abs{W-\mu_W}\ge\varepsilon)\le\frac{\Var(W)}{\varepsilon^{2}}.
\]
\end{lema}

\begin{proof}
Sea $\ind_A$ la variable que vale $1$ si $Y\ge a$ y $0$ en otro caso. Como $Y\ge0$, se tiene $Y\ge a\ind_A$ en todos los casos, y tomando esperanzas $\E[Y]\ge a\PP(Y\ge a)$. Para la segunda desigualdad se aplica la primera a $Y=(W-\mu_W)^{2}$ y $a=\varepsilon^{2}$, observando que $\abs{W-\mu_W}\ge\varepsilon$ equivale a $(W-\mu_W)^{2}\ge\varepsilon^{2}$.
\end{proof}

\begin{proposicion}[Un $\ECM$ que tiende a cero implica consistencia]\label{prop:t2-ecm-consistencia}
Si $\ECM(\widehat\theta_n)\to0$ para todo $\theta$, entonces $\widehat\theta_n$ es consistente. En particular, basta que el sesgo y la varianza tiendan a cero.
\end{proposicion}

\begin{proof}
Por la desigualdad de Markov aplicada a $Y=(\widehat\theta_n-\theta)^{2}$ con $a=\varepsilon^{2}$, $\PP(\abs{\widehat\theta_n-\theta}>\varepsilon)\le\PP\bigl((\widehat\theta_n-\theta)^{2}\ge\varepsilon^{2}\bigr)\le\ECM(\widehat\theta_n)/\varepsilon^{2}\to0$. La segunda afirmación se sigue de la \cref{prop:t2-ecm}.
\end{proof}

\begin{teorema}[Ley débil de los grandes números]\label{teo:t2-lgn}
Si $X_1,X_2,\dots$ son \iid\ con media $\mu$ y varianza finita $\sigma^{2}$, entonces $\Xbar_n\convp\mu$; es decir, la media muestral es un estimador consistente de la media.
\end{teorema}

\begin{proof}
$\Xbar_n$ es insesgado y $\ECM(\Xbar_n)=\Var(\Xbar_n)=\sigma^{2}/n\to0$; se aplica la \cref{prop:t2-ecm-consistencia}. Explícitamente, $\PP(\abs{\Xbar_n-\mu}>\varepsilon)\le\sigma^{2}/(n\varepsilon^{2})$.
\end{proof}

La cota de Chebyshev es válida para cualquier distribución, pero es muy conservadora; en el \cref{tema:3} el teorema central del límite dará aproximaciones mucho más precisas.

\begin{observacion}[Errores frecuentes]\label{obs:t2-errores}
(a) Confundir el parámetro $\mu$ (un número fijo y desconocido) con el estadístico $\Xbar$ (una variable aleatoria) o con su valor observado $\xbar$ (un número conocido).
(b) Confundir la desviación típica de la población, $\sigma$, con el error estándar de la media, $\sigma/\sqrt n$.
(c) Pensar que «insesgado» significa «cercano al parámetro»: un estimador insesgado con mucha varianza puede ser casi siempre malo.
(d) Usar $\Var(X+Y)=\Var X+\Var Y$ sin comprobar la independencia (o, al menos, la incorrelación).
\end{observacion}

\begin{observacion}[Conexión con el curso]
El módulo~5 retoma estos conceptos (estimadores, sesgo y error cuadrático medio, ley de los grandes números y método de Monte Carlo) y el módulo~8 desarrolla la descomposición sesgo--varianza del error de predicción. En los módulos~22 y~23, la exactitud de un modelo en un benchmark es una proporción muestral (\cref{ej:t2-proporcion}) y la métrica pass@$k$ se estima con un estimador insesgado que se construye con un argumento de muestreo (\cref{eje:t2-passk}).
\end{observacion}

\begin{lecturas}
\item Wasserman~\cite{wasserman}, cap.~6 «Models, Statistical Inference and Learning» (estimación puntual, sesgo, error estándar, consistencia) y cap.~5 «Convergence of Random Variables».
\item Blitzstein y Hwang~\cite{bh}, §6.3 «Sample moments» y §10.2 «Law of large numbers».
\item Stats~200~\cite{stats200}, clase~12 «Parametric models, method of moments estimation» (sesgo, varianza y error cuadrático medio).
\item Orloff y Bloom~\cite{mit1805}, lectura~10a «Introduction to Statistics».
\end{lecturas}

\section*{Ejercicios del tema 2}
\addcontentsline{toc}{section}{Ejercicios del tema 2}

\begin{ejercicio}{1}{eje:t2-conceptos}
Se quiere conocer la exactitud real $p$ de un modelo de lenguaje en un banco de un millón de preguntas. Se eligen al azar $500$ preguntas, el modelo acierta $410$ y se informa de que «la exactitud es del $82\,\%$». Identifique la población, la muestra, el parámetro, el estimador y la estimación. ¿Es aleatorio el número $0{,}82$? ¿Y $\widehat p$?
\end{ejercicio}

\begin{ejercicio}{2}{eje:t2-enumeracion}
Una población está formada por los valores $1,3,5$, cada uno con probabilidad $1/3$. Se extraen muestras de tamaño $n=2$ con reemplazamiento. (a) Calcule $\mu$ y $\sigma^{2}$. (b) Enumere las nueve muestras posibles y obtenga la distribución muestral de $\Xbar$; compruebe que $\E[\Xbar]=\mu$ y $\Var(\Xbar)=\sigma^{2}/2$. (c) Obtenga la distribución de $S^{2}$ y compruebe que $\E[S^{2}]=\sigma^{2}$. ¿Cuánto vale $\E[\widehat\sigma^{2}]$ con denominador $n$?
\end{ejercicio}

\begin{ejercicio}{2}{eje:t2-tres-estimadores}
Sean $X_1,X_2,X_3$ \iid\ con media $\mu$ y varianza $\sigma^{2}$. Se consideran $T_1=\Xbar$, $T_2=(X_1+2X_2+X_3)/4$, $T_3=X_1+X_2-X_3$ y $T_4=(X_1+X_2+X_3)/4$. (a) Calcule el sesgo, la varianza y el $\ECM$ de cada uno. (b) ¿Cuál es el mejor de los insesgados? (c) ¿Para qué valores de $\mu$ tiene $T_4$ menor $\ECM$ que $T_1$?
\end{ejercicio}

\begin{ejercicio}{2}{eje:t2-poisson}
Sean $X_1,\dots,X_n\sim\Pois(\lambda)$ \iid, y $\widehat\lambda=\Xbar$. (a) Demuestre que $\widehat\lambda$ es insesgado y calcule su varianza, su error estándar y su $\ECM$. (b) En $n=50$ bloques de texto se cuentan los errores ortográficos y se obtiene $\xbar=3{,}2$. Dé la estimación de $\lambda$ y su error estándar estimado.
\fuente{adaptado de Stats~200, Stanford~\cite{stats200}, clase~12, ejemplo~12.5.}
\end{ejercicio}

\begin{ejercicio}{2}{eje:t2-expo-sesgo}
Sean $X_1,\dots,X_n\sim\Expo(\lambda)$ \iid\ ($n\ge2$), con media $1/\lambda$. Se admite que $\E[1/\Xbar]=\dfrac{n\lambda}{n-1}$. (a) Calcule el sesgo de $\widehat\lambda=1/\Xbar$. (b) Proponga un estimador insesgado de $\lambda$ de la forma $c/\Xbar$. (c) Con $n=5$ y $\sum x_i=10$, calcule ambas estimaciones.
\fuente{adaptado de Stats~200, Stanford~\cite{stats200}, clase~12, ejemplo~12.6.}
\end{ejercicio}

\begin{ejercicio}{2}{eje:t2-sumas-solapadas}
Sean $X_1,X_2,\dots$ \iid\ con distribución exponencial de tasa $2$ (varianza $1/4$) y sea $n\ge8$. Se definen $X=X_1+\dots+X_n$ e $Y=X_{n-7}+X_{n-6}+\dots+X_{2n-8}$. Calcule $\Cov(X,Y)$ y $\Cor(X,Y)$.
\fuente{adaptado de Orloff y Bloom, MIT 18.05~\cite{mit1805}, \emph{problem set}~5, problema~5.}
\end{ejercicio}

\begin{ejercicio}{2}{eje:t2-benchmark}
Un modelo acierta $300$ de $400$ preguntas elegidas al azar de un banco grande. (a) Estime la exactitud y su error estándar. (b) ¿Cuántas preguntas harían falta para reducir el error estándar a la mitad, si la exactitud estimada no cambia?
\end{ejercicio}

\begin{ejercicio}{2}{eje:t2-combinar}
Se dispone de dos estimadores independientes e insesgados $A$ y $B$ de un mismo parámetro, con $\Var(A)=4$ y $\Var(B)=1$. (a) Demuestre que $T_w=wA+(1-w)B$ es insesgado para todo $w\in\R$. (b) Halle el $w$ que minimiza $\Var(T_w)$ y el valor mínimo. (c) Generalice a varianzas $\sigma_A^{2}$ y $\sigma_B^{2}$.
\end{ejercicio}

\begin{ejercicio}{2}{eje:t2-chebyshev-n}
Se quiere estimar la exactitud $p$ de un modelo con una muestra de $n$ preguntas de modo que $\PP(\abs{\widehat p-p}\ge0{,}05)\le0{,}05$ sea cual sea $p$. Usando la desigualdad de Chebyshev y la cota $p(1-p)\le1/4$, halle un $n$ suficiente.
\end{ejercicio}

\begin{ejercicio}{3}{eje:t2-ecm-normal}
Para una muestra \iid\ de una normal se sabe que $\Var(S^{2})=2\sigma^{4}/(n-1)$. Calcule el $\ECM$ de $S^{2}$ y el de $\widehat\sigma^{2}=\frac{n-1}{n}S^{2}$, y demuestre que el estimador sesgado $\widehat\sigma^{2}$ tiene menor $\ECM$ para todo $n\ge2$.
\end{ejercicio}

\begin{ejercicio}{3}{eje:t2-uniforme}
Sean $X_1,\dots,X_n\sim\Unif(0,\theta)$ \iid\ y $M=\max_iX_i$. (a) Demuestre que $T_1=2\Xbar$ es insesgado y que $\Var(T_1)=\theta^{2}/(3n)$. (b) Demuestre que $\PP(M\le x)=(x/\theta)^{n}$ para $0\le x\le\theta$, calcule $\E[M]$ y $\E[M^{2}]$ y deduzca que $T_2=\frac{n+1}{n}M$ es insesgado con $\Var(T_2)=\dfrac{\theta^{2}}{n(n+2)}$. (c) Demuestre que $\Var(T_2)\le\Var(T_1)$, con igualdad solo si $n=1$.
\fuente{adaptado de Stats~200, Stanford~\cite{stats200}, \emph{homework}~5, problema~3, y de Orloff y Bloom, MIT 18.05~\cite{mit1805}, \emph{problem set}~6, problema~3.}
\end{ejercicio}

\begin{ejercicio}{3}{eje:t2-max-consistente}
Con la notación del \cref{eje:t2-uniforme}, demuestre que $M$ es un estimador consistente de $\theta$ de dos maneras: (a) calculando directamente $\PP(\abs{M-\theta}>\varepsilon)$; (b) mediante la \cref{prop:t2-ecm-consistencia}.
\end{ejercicio}

\begin{ejercicio}{3}{eje:t2-cov-desv}
Sean $X_1,\dots,X_n$ \iid\ con varianza $\sigma^{2}$. Demuestre que $\Cov(\Xbar,X_i-\Xbar)=0$ para cada $i$. (Para datos normales, este hecho es la clave de la independencia entre $\Xbar$ y $S^{2}$.)
\end{ejercicio}

\begin{ejercicio}{3}{eje:t2-s-sesgo}
Demuestre que $\E[S]\le\sigma$, con igualdad solo si $S$ es constante con probabilidad~1. (\emph{Indicación:} $\Var(S)\ge0$.)
\end{ejercicio}

\begin{ejercicio}{3}{eje:t2-passk}
Un modelo genera $n$ soluciones independientes a un problema, cada una correcta con probabilidad $p$, y sea $C\sim\Bin(n,p)$ el número de correctas. Para $1\le k\le n$, la métrica pass@$k$ es la probabilidad de que al menos una de $k$ soluciones independientes sea correcta, $1-(1-p)^{k}$. Demuestre que
\[
\widehat{\pi}_k=1-\binom{n-C}{k}\Big/\binom{n}{k}
\]
es un estimador insesgado de pass@$k$. (\emph{Indicación:} interprete $\binom{n-C}{k}/\binom nk$ como la probabilidad de que un subconjunto de $k$ soluciones elegido al azar no contenga ninguna correcta.)
\end{ejercicio}

% =====================================================================
\chapter{Distribuciones muestrales y teorema central del límite}
\label{tema:3}
% =====================================================================

Un estadístico es una variable aleatoria y, para hacer inferencia, hace falta conocer (al menos aproximadamente) su distribución. Este tema estudia la distribución de la media muestral: exacta cuando los datos son normales y aproximada, gracias al teorema central del límite, en cualquier otro caso. Se incluyen la aproximación normal a la binomial y las distribuciones $\chi^{2}$ y $t$ de Student, que se usarán en los intervalos de confianza y en los contrastes. Se sigue el capítulo~10 de Blitzstein y Hwang~\cite{bh} (§10.3 «Central limit theorem» y §10.4 «Chi-Square and Student-t»), la lectura~6b del curso 18.05 del MIT~\cite{mit1805}, «Central Limit Theorem and the Law of Large Numbers», y el capítulo~7 de OpenStax~\cite{openstax} (§7.1 y §7.2).

\section{La distribución normal}

\begin{definicion}[Distribución normal]\label{def:t3-normal}
Una variable $X$ tiene distribución \emph{normal} de media $\mu\in\R$ y varianza $\sigma^{2}>0$, y se escribe $X\sim\Nor(\mu,\sigma^{2})$, si su densidad es
\[
f(x)=\frac{1}{\sigma\sqrt{2\pi}}\exp\Bigl(-\frac{(x-\mu)^{2}}{2\sigma^{2}}\Bigr),\qquad x\in\R.
\]
La \emph{normal estándar} es $Z\sim\Nor(0,1)$; su función de distribución se denota $\Phi(z)=\PP(Z\le z)$ y su cuantil de orden $q$, $z_q=\Phi^{-1}(q)$.
\end{definicion}

La función $\Phi$ no tiene expresión mediante funciones elementales y se tabula (\cref{tab:phi}). Por la simetría de la densidad respecto de~$0$, $\Phi(-z)=1-\Phi(z)$ y $z_{1-q}=-z_q$.

\begin{proposicion}[Tipificación]\label{prop:t3-tipificacion}
Si $X\sim\Nor(\mu,\sigma^{2})$ y $a\neq0$, entonces $aX+b\sim\Nor(a\mu+b,a^{2}\sigma^{2})$. En particular, $Z=(X-\mu)/\sigma\sim\Nor(0,1)$ y
\[
\PP(X\le x)=\Phi\Bigl(\frac{x-\mu}{\sigma}\Bigr).
\]
\end{proposicion}

La demostración es el \cref{eje:t3-lineal-normal}. Con la tabla se obtiene la conocida regla empírica: una normal toma valores a menos de $1$, $2$ y $3$ desviaciones típicas de su media con probabilidades $2\Phi(1)-1\approx0{,}683$, $2\Phi(2)-1\approx0{,}954$ y $2\Phi(3)-1\approx0{,}997$.

\begin{teorema}[Combinaciones lineales de normales independientes]\label{teo:t3-suma-normales}
Si $X_1,\dots,X_n$ son independientes con $X_i\sim\Nor(\mu_i,\sigma_i^{2})$, entonces $\sum_ia_iX_i\sim\Nor\bigl(\sum_ia_i\mu_i,\sum_ia_i^{2}\sigma_i^{2}\bigr)$ para cualesquiera $a_i$ no todos nulos. En particular, si $X_1,\dots,X_n\sim\Nor(\mu,\sigma^{2})$ son \iid,
\[
\Xbar\sim\Nor\Bigl(\mu,\frac{\sigma^{2}}{n}\Bigr).
\]
\end{teorema}

La media y la varianza se siguen de la \cref{prop:t2-props}; que la distribución siga siendo normal se demuestra con funciones generadoras de momentos~\cite[§6.6]{bh} y se admite aquí sin demostración.

\section{El teorema central del límite}

Cuando los datos no son normales, la distribución exacta de $\Xbar$ suele ser complicada. El teorema central del límite afirma que, sin embargo, para $n$ grande es aproximadamente normal, sea cual sea la distribución de partida, con tal de que tenga varianza finita.

\begin{teorema}[Teorema central del límite]\label{teo:t3-tcl}
Sean $X_1,X_2,\dots$ \iid\ con media $\mu$ y varianza $0<\sigma^{2}<\infty$, y sea $S_n=X_1+\dots+X_n$. Entonces, para todo $z\in\R$,
\[
\lim_{n\to\infty}\PP\Bigl(\frac{\Xbar_n-\mu}{\sigma/\sqrt n}\le z\Bigr)=\lim_{n\to\infty}\PP\Bigl(\frac{S_n-n\mu}{\sigma\sqrt n}\le z\Bigr)=\Phi(z).
\]
\end{teorema}

Las dos variables tipificadas son la misma (basta multiplicar numerador y denominador por $n$). La demostración, con funciones generadoras de momentos, está en~\cite[§10.3]{bh}; aquí interesa su uso.

\begin{observacion}[Uso práctico]\label{obs:t3-uso}
Para $n$ «grande» se aproxima
\[
\Xbar\approx\Nor\Bigl(\mu,\frac{\sigma^{2}}{n}\Bigr),\qquad S_n\approx\Nor(n\mu,n\sigma^{2}).
\]
Una regla habitual es $n\ge30$~\cite[§7.1]{openstax}, pero la calidad de la aproximación depende de la forma de la distribución: con datos simétricos basta un $n$ pequeño y con datos muy asimétricos hace falta bastante más. El teorema exige varianza finita: para distribuciones de cola muy pesada, como la de Cauchy (o la $t$ con un grado de libertad), la media muestral no se estabiliza~\cite[§10.3]{bh}.
\end{observacion}

\begin{ejemplo}\label{ej:t3-tcl}
La longitud de las respuestas de un modelo tiene media $\mu=200$ tokens y desviación típica $\sigma=80$, con distribución asimétrica. Para una muestra de $n=64$ respuestas, el error estándar de la media es $80/\sqrt{64}=10$, y por el teorema central del límite
\[
\PP(\Xbar>215)\approx1-\Phi\Bigl(\frac{215-200}{10}\Bigr)=1-\Phi(1{,}5)=1-0{,}9332=0{,}0668.
\]
Una respuesta individual supera los $215$ tokens con mucha mayor probabilidad: la media de $64$ respuestas fluctúa ocho veces menos que cada respuesta.
\end{ejemplo}

\section{Aproximación normal a la binomial}

Si $X\sim\Bin(n,p)$, entonces $X$ tiene la misma distribución que la suma de $n$ variables $\Ber(p)$ \iid, con $\mu=p$ y $\sigma^{2}=p(1-p)$. El teorema central del límite da el resultado clásico de De Moivre y Laplace.

\begin{corolario}[Aproximación normal a la binomial]\label{cor:t3-binomial}
Si $X\sim\Bin(n,p)$ con $0<p<1$, para $n$ grande
\[
X\approx\Nor\bigl(np,\,np(1-p)\bigr),\qquad \widehat p=\frac Xn\approx\Nor\Bigl(p,\,\frac{p(1-p)}{n}\Bigr).
\]
Una regla habitual para usarla es $np\ge10$ y $n(1-p)\ge10$.
\end{corolario}

\begin{definicion}[Corrección por continuidad]\label{def:t3-cc}
Al aproximar una variable con valores enteros por una continua, cada valor entero $k$ se identifica con el intervalo $[k-\frac12,k+\frac12]$. Así, para enteros $a\le b$,
\[
\PP(a\le X\le b)\approx\Phi\Bigl(\frac{b+\frac12-np}{\sqrt{np(1-p)}}\Bigr)-\Phi\Bigl(\frac{a-\frac12-np}{\sqrt{np(1-p)}}\Bigr).
\]
\end{definicion}

\begin{ejemplo}\label{ej:t3-binomial}
Sea $X\sim\Bin(100;\,0{,}5)$: $np=50$ y $\sqrt{np(1-p)}=5$. Sin corrección, $\PP(X\ge60)\approx1-\Phi(2)=0{,}0228$; con corrección, $\PP(X\ge60)=\PP(X\ge59{,}5)\approx1-\Phi(1{,}9)=0{,}0287$. El valor exacto es $0{,}0284$: la corrección mejora mucho la aproximación.
\end{ejemplo}

\section{Distribuciones \texorpdfstring{$\chi^{2}$}{chi cuadrado} y \texorpdfstring{$t$}{t} de Student}

Cuando $\sigma$ es desconocida y se sustituye por $S$, la variable $(\Xbar-\mu)/(S/\sqrt n)$ ya no es normal. Su distribución exacta, para datos normales, se describe con dos familias derivadas de la normal~\cite[§10.4]{bh}.

\begin{definicion}[Distribuciones $\chi^{2}$ y $t$]\label{def:t3-chi-t}
Sean $Z_1,\dots,Z_k\sim\Nor(0,1)$ \iid\ La distribución de $V=Z_1^{2}+\dots+Z_k^{2}$ se llama \emph{$\chi^{2}$ con $k$ grados de libertad} y se escribe $V\sim\chi^{2}_k$. Si $Z\sim\Nor(0,1)$ y $V\sim\chi^{2}_k$ son independientes, la distribución de
\[
T=\frac{Z}{\sqrt{V/k}}
\]
se llama \emph{$t$ de Student con $k$ grados de libertad} y se escribe $T\sim t_k$. Sus cuantiles se denotan $\chi^{2}_{k;q}$ y $t_{k;q}$ (\cref{tab:t,tab:chi2}).
\end{definicion}

La $\chi^{2}_k$ toma valores positivos, tiene media $k$ y varianza $2k$ (\cref{eje:t3-chi-momentos}). La $t_k$ es simétrica respecto de~$0$ (\cref{eje:t3-t-simetrica}), tiene colas más pesadas que la normal y converge a ella cuando $k\to\infty$: por ejemplo, $t_{5;0{,}975}=2{,}571$, $t_{30;0{,}975}=2{,}042$ y $z_{0{,}975}=1{,}960$.

\begin{teorema}[Teorema de Fisher]\label{teo:t3-fisher}
Sean $X_1,\dots,X_n\sim\Nor(\mu,\sigma^{2})$ \iid\ con $n\ge2$. Entonces $\Xbar$ y $S^{2}$ son independientes, $\Xbar\sim\Nor(\mu,\sigma^{2}/n)$ y
\[
\frac{(n-1)S^{2}}{\sigma^{2}}\sim\chi^{2}_{n-1}.
\]
\end{teorema}

Se admite sin demostración~\cite[§10.4]{bh}. Los $n-1$ grados de libertad son los de la \cref{obs:t1-n-1}: la suma de cuadrados respecto de $\Xbar$ equivale a una suma de $n-1$ cuadrados de normales estándar independientes. Como $\E[\chi^{2}_{n-1}]=n-1$, el teorema reproduce, para datos normales, el resultado $\E[S^{2}]=\sigma^{2}$ del \cref{teo:t2-s2}, y como $\Var(\chi^{2}_{n-1})=2(n-1)$, da $\Var(S^{2})=2\sigma^{4}/(n-1)$, que se usó en el \cref{eje:t2-ecm-normal}.

\begin{corolario}[Estadístico $t$]\label{cor:t3-t}
En las condiciones del \cref{teo:t3-fisher},
\[
T=\frac{\Xbar-\mu}{S/\sqrt n}\sim t_{n-1}.
\]
\end{corolario}

\begin{proof}
Sean $Z=(\Xbar-\mu)/(\sigma/\sqrt n)\sim\Nor(0,1)$ y $V=(n-1)S^{2}/\sigma^{2}\sim\chi^{2}_{n-1}$, independientes por el \cref{teo:t3-fisher}. Entonces
\[
\frac{Z}{\sqrt{V/(n-1)}}=\frac{(\Xbar-\mu)\sqrt n/\sigma}{S/\sigma}=\frac{\Xbar-\mu}{S/\sqrt n},
\]
que tiene distribución $t_{n-1}$ por la \cref{def:t3-chi-t}.
\end{proof}

\begin{observacion}[Errores frecuentes]\label{obs:t3-errores}
(a) El teorema central del límite no dice que los \emph{datos} se vuelvan normales al aumentar $n$, sino que lo hace la distribución de su \emph{media} (o de su suma).
(b) Usar la desviación típica de la población, $\sigma$, en lugar del error estándar, $\sigma/\sqrt n$, al tipificar una media.
(c) Olvidar la corrección por continuidad con $n$ moderado o confundir su sentido: $\PP(X>59)=\PP(X\ge60)\approx\PP(Y\ge59{,}5)$, no $\PP(Y\ge60{,}5)$.
(d) Usar la $t$ para datos claramente no normales y $n$ pequeño: el \cref{cor:t3-t} exige normalidad; para $n$ grande, la $t$ y la normal prácticamente coinciden.
\end{observacion}

\begin{observacion}[Conexión con el curso]
El módulo~5 enuncia el teorema central del límite y lo aplica al método de Monte Carlo. En el módulo~7, el gradiente calculado con un minilote de $B$ ejemplos es una media muestral: su ruido decrece como $1/\sqrt B$. En los módulos~22 y~23, el error estándar de una exactitud medida sobre $n$ preguntas es $\sqrt{p(1-p)/n}$, que con $n=1000$ y $p=0{,}8$ vale unos $1{,}3$ puntos porcentuales; las diferencias de un punto entre modelos en un benchmark de ese tamaño están, por tanto, dentro del ruido de muestreo.
\end{observacion}

\begin{lecturas}
\item Blitzstein y Hwang~\cite{bh}, §5.4 «Normal», §10.3 «Central limit theorem» y §10.4 «Chi-Square and Student-t».
\item Orloff y Bloom~\cite{mit1805}, lecturas~6b «Central Limit Theorem and the Law of Large Numbers» y 5c «Gallery of Continuous Random Variables».
\item OpenStax~\cite{openstax}, cap.~7: §7.1 «The Central Limit Theorem for Sample Means (Averages)» y §7.2 «The Central Limit Theorem for Sums».
\item Wasserman~\cite{wasserman}, cap.~5 «Convergence of Random Variables».
\end{lecturas}

\section*{Ejercicios del tema 3}
\addcontentsline{toc}{section}{Ejercicios del tema 3}

\begin{ejercicio}{1}{eje:t3-vf}
Indique si cada afirmación es verdadera o falsa. (a) Por el teorema central del límite, si $n=1000$ los datos de una muestra de una exponencial son aproximadamente normales. (b) Si los datos son normales, $\Xbar$ es exactamente normal para todo $n$. (c) El error estándar de la media de $100$ observaciones es la décima parte de la desviación típica de cada observación. (d) La $t_{10}$ tiene más probabilidad en las colas que la $\Nor(0,1)$. (e) Si $X\sim\Bin(20;\,0{,}02)$, la aproximación normal es adecuada.
\end{ejercicio}

\begin{ejercicio}{1}{eje:t3-tabla}
Con la \cref{tab:phi,tab:zq}, calcule: (a) $\Phi(1{,}23)$; (b) $\PP(Z>-0{,}5)$; (c) $\PP(-1<Z<2)$; (d) $z_{0{,}95}$ y $z_{0{,}99}$; (e) el valor $c>0$ tal que $\PP(\abs Z\le c)=0{,}90$.
\end{ejercicio}

\begin{ejercicio}{2}{eje:t3-normal-basico}
Una puntuación $X$ sigue una distribución $\Nor(100,15^{2})$. Calcule $\PP(X>130)$, $\PP(85<X<115)$ y el percentil $90$ de $X$.
\end{ejercicio}

\begin{ejercicio}{2}{eje:t3-openstax-media}
Una población tiene media $90$ y desviación típica $15$. Se toma una muestra de tamaño $n=25$. (a) Calcule aproximadamente $\PP(85<\Xbar<92)$. (b) ¿Qué valor de $\Xbar$ está dos errores estándar por encima de la media?
\fuente{adaptado de OpenStax, \emph{Introductory Statistics}~2e~\cite{openstax}, §7.1, ejemplo~7.1.}
\end{ejercicio}

\begin{ejercicio}{2}{eje:t3-openstax-suma}
Con la población del \cref{eje:t3-openstax-media} y una muestra de tamaño $n=80$: (a) calcule aproximadamente la probabilidad de que la suma de los $80$ valores supere $7500$; (b) halle el valor de la suma que está $1{,}5$ desviaciones típicas por encima de su media.
\fuente{adaptado de OpenStax, \emph{Introductory Statistics}~2e~\cite{openstax}, §7.2, ejemplo~7.5.}
\end{ejercicio}

\begin{ejercicio}{2}{eje:t3-elecciones}
El $50\,\%$ de una población apoya al candidato A y el $30\,\%$ a otros candidatos menores. Se encuesta a $400$ personas elegidas al azar. Con el teorema central del límite, estime: (a) la probabilidad de que al menos el $52{,}5\,\%$ de los encuestados apoye a A; (b) la probabilidad de que menos del $31\,\%$ apoye a los candidatos menores.
\fuente{adaptado de Orloff y Bloom, MIT 18.05~\cite{mit1805}, \emph{problem set}~4, problema~2.}
\end{ejercicio}

\begin{ejercicio}{2}{eje:t3-redondeo}
Un comercio redondea el importe de cada pedido a los $5$ céntimos más próximos. Se supone que el error de redondeo de cada pedido toma los valores $-2,-1,0,1,2$ céntimos con igual probabilidad, de forma independiente. Con $1000$ pedidos diarios, estime la probabilidad de que el error total del día supere los $100$ céntimos en valor absoluto.
\fuente{adaptado de Orloff y Bloom, MIT 18.05~\cite{mit1805}, \emph{problem set}~4, problema~3.}
\end{ejercicio}

\begin{ejercicio}{2}{eje:t3-binomial}
Sea $X\sim\Bin(100;\,0{,}5)$ el número de caras en $100$ lanzamientos. (a) Aproxime $\PP(X=50)$ con la corrección por continuidad. (b) Aproxime $\PP(40\le X\le60)$ con y sin corrección. (c) Compare con los valores exactos $\PP(X=50)=0{,}0796$ y $\PP(40\le X\le60)=0{,}9648$.
\end{ejercicio}

\begin{ejercicio}{2}{eje:t3-benchmark}
Un modelo tiene exactitud real $p=0{,}8$ en un banco de preguntas muy grande. Se evalúa con $n=1000$ preguntas elegidas al azar. (a) Calcule el error estándar de $\widehat p$. (b) Estime $\PP(\widehat p\ge0{,}82)$. (c) Halle un intervalo centrado en $0{,}8$ que contenga a $\widehat p$ con probabilidad aproximada $0{,}95$.
\end{ejercicio}

\begin{ejercicio}{2}{eje:t3-tamano}
(a) ¿Qué tamaño muestral garantiza, según el teorema central del límite, que $\PP(\abs{\Xbar-\mu}<0{,}1\,\sigma)\ge0{,}95$? (b) Aplíquelo para estimar una exactitud con error menor que $0{,}05$ con probabilidad $0{,}95$, sea cual sea $p$. Compare con el \cref{eje:t2-chebyshev-n}.
\end{ejercicio}

\begin{ejercicio}{2}{eje:t3-t-cuantiles}
Con la \cref{tab:t}: (a) obtenga $t_{9;0{,}975}$, $t_{29;0{,}975}$ y $t_{120;0{,}975}$ y compárelos con $z_{0{,}975}$; (b) si $T\sim t_{15}$, calcule $\PP(T>2{,}131)$ y $\PP(\abs T\le1{,}753)$; (c) si $V\sim\chi^{2}_{10}$, halle $c$ con $\PP(V>c)=0{,}05$.
\end{ejercicio}

\begin{ejercicio}{3}{eje:t3-lineal-normal}
Demuestre la \cref{prop:t3-tipificacion}: si $X\sim\Nor(\mu,\sigma^{2})$ y $a\ne0$, entonces $Y=aX+b\sim\Nor(a\mu+b,a^{2}\sigma^{2})$. (\emph{Indicación:} calcule $\PP(Y\le y)$ distinguiendo $a>0$ y $a<0$, y derive.)
\end{ejercicio}

\begin{ejercicio}{3}{eje:t3-chi1}
Sea $Z\sim\Nor(0,1)$ y $V=Z^{2}$. Demuestre que $\PP(V\le v)=2\Phi(\sqrt v)-1$ para $v>0$ y deduzca que la densidad de la $\chi^{2}_1$ es
\[
f_V(v)=\frac{1}{\sqrt{2\pi v}}\,e^{-v/2},\qquad v>0.
\]
\end{ejercicio}

\begin{ejercicio}{3}{eje:t3-chi-momentos}
(a) Demuestre, integrando por partes, que si $Z\sim\Nor(0,1)$ entonces $\E[Z^{4}]=3\E[Z^{2}]=3$. (b) Deduzca que si $V\sim\chi^{2}_k$, entonces $\E[V]=k$ y $\Var(V)=2k$.
\end{ejercicio}

\begin{ejercicio}{3}{eje:t3-t-simetrica}
Sea $T\sim t_k$. Demuestre que $T$ y $-T$ tienen la misma distribución y deduzca que $t_{k;1-q}=-t_{k;q}$ y que $\PP(\abs T\le t_{k;1-\alpha/2})=1-\alpha$. (Se admite que la distribución de $T$ es continua.)
\end{ejercicio}

% =====================================================================
\chapter{Estimación puntual e intervalos de confianza}
\label{tema:4}
% =====================================================================

El \cref{tema:2} enseñó a juzgar un estimador; este tema enseña a construirlo. Se presentan los dos métodos generales de estimación puntual, el método de los momentos y el de máxima verosimilitud, y se aplican a los modelos de Bernoulli, Poisson, normal y exponencial. Después se pasa de la estimación puntual a la estimación por intervalos: intervalos de confianza para la media con varianza conocida y desconocida, para proporciones y cálculo del tamaño muestral. Se sigue el capítulo~9 de Wasserman~\cite{wasserman}, «Parametric Inference» (método de los momentos y máxima verosimilitud), las clases~12, 13 y~18 de Stats~200~\cite{stats200}, las lecturas~10b, 22 y~23a del curso 18.05 del MIT~\cite{mit1805} y el capítulo~8 de OpenStax~\cite{openstax}.

\section{El método de los momentos}

Sea $X_1,\dots,X_n$ una muestra \iid\ de un modelo $F_\theta$ con $\theta=(\theta_1,\dots,\theta_k)$. El \emph{momento poblacional} de orden $j$ es $m_j(\theta)=\E_\theta[X^{j}]$ y el \emph{momento muestral} de orden $j$ es $\widehat m_j=\frac1n\sum_iX_i^{j}$. Por la ley de los grandes números, $\widehat m_j\convp m_j(\theta)$.

\begin{definicion}[Estimador por el método de los momentos]\label{def:t4-momentos}
El \emph{estimador por el método de los momentos} $\widehat\theta$ es la solución del sistema de $k$ ecuaciones
\[
m_j(\widehat\theta)=\widehat m_j,\qquad j=1,\dots,k,
\]
es decir, el valor del parámetro con el que los primeros $k$ momentos del modelo coinciden con los de la muestra~\cite[cap.~9]{wasserman}.
\end{definicion}

\begin{ejemplo}[Momentos en los modelos básicos]\label{ej:t4-momentos}
\leavevmode
\begin{enumerate}
  \item Bernoulli: $\E X=p$, luego $\widehat p=\Xbar$.
  \item Poisson: $\E X=\lambda$, luego $\widehat\lambda=\Xbar$.
  \item Exponencial de tasa $\lambda$: $\E X=1/\lambda$, luego $\widehat\lambda=1/\Xbar$.
  \item Normal: $\E X=\mu$ y $\E X^{2}=\sigma^{2}+\mu^{2}$. De $\widehat\mu=\Xbar$ y $\widehat\sigma^{2}+\widehat\mu^{2}=\frac1n\sum_iX_i^{2}$ se obtiene $\widehat\sigma^{2}=\frac1n\sum_iX_i^{2}-\Xbar^{2}=\frac1n\sum_i(X_i-\Xbar)^{2}$, la varianza con denominador $n$.
  \item Uniforme en $[0,\theta]$: $\E X=\theta/2$, luego $\widehat\theta=2\Xbar$ (\cref{eje:t2-uniforme}).
\end{enumerate}
\end{ejemplo}

El método es sencillo y produce estimadores consistentes en condiciones generales, pero no siempre los mejores y, a veces, incompatibles con los datos (\cref{eje:t4-uniforme-ab}).

\section{Máxima verosimilitud}

\begin{definicion}[Verosimilitud y estimador de máxima verosimilitud]\label{def:t4-emv}
Sean $x_1,\dots,x_n$ los datos observados y $f(x\mid\theta)$ la función de probabilidad (caso discreto) o la densidad (caso continuo) del modelo. La \emph{función de verosimilitud} y la \emph{log-verosimilitud} son
\[
\Lik(\theta)=\prod_{i=1}^{n}f(x_i\mid\theta),\qquad \ell(\theta)=\ln\Lik(\theta)=\sum_{i=1}^{n}\ln f(x_i\mid\theta).
\]
El \emph{estimador de máxima verosimilitud} (EMV) es el valor $\widehat\theta=\argmax_{\theta\in\Theta}\Lik(\theta)$~\cite[cap.~9]{wasserman}.
\end{definicion}

La verosimilitud es la probabilidad (o densidad) de los datos observados vista como función del parámetro, con los datos fijos. Como el logaritmo es estrictamente creciente, maximizar $\Lik$ equivale a maximizar $\ell$, que es una suma y se deriva con más facilidad. El procedimiento habitual es resolver la \emph{ecuación de verosimilitud} $\ell'(\theta)=0$ y comprobar que la solución es un máximo global.

\begin{teorema}[EMV de Bernoulli]\label{teo:t4-emv-bernoulli}
Si $x_1,\dots,x_n\in\{0,1\}$ son observaciones \iid\ de $\Ber(p)$ y $k=\sum_ix_i$ es el número de unos, con $0<k<n$, el EMV es $\widehat p=k/n=\xbar$.
\end{teorema}

\begin{proof}
La función de probabilidad es $f(x\mid p)=p^{x}(1-p)^{1-x}$ para $x\in\{0,1\}$. Entonces $\Lik(p)=p^{k}(1-p)^{n-k}$ y, para $0<p<1$,
\[
\ell(p)=k\ln p+(n-k)\ln(1-p),\qquad \ell'(p)=\frac kp-\frac{n-k}{1-p}.
\]
La ecuación $\ell'(p)=0$ equivale a $k(1-p)=(n-k)p$, es decir, $p=k/n$. Como $\ell''(p)=-k/p^{2}-(n-k)/(1-p)^{2}<0$, $\ell$ es estrictamente cóncava en $(0,1)$ y el punto crítico es el máximo global. (En los extremos, $\Lik(0)=\Lik(1)=0$ cuando $0<k<n$.)
\end{proof}

Si $k=0$ o $k=n$, el máximo está en $\widehat p=0$ o $\widehat p=1$, respectivamente, y la fórmula $\widehat p=k/n$ sigue siendo válida.

\begin{teorema}[EMV de la normal]\label{teo:t4-emv-normal}
Si $x_1,\dots,x_n$ son observaciones \iid\ de $\Nor(\mu,\sigma^{2})$, no todas iguales, el EMV de $(\mu,\sigma^{2})$ es
\[
\widehat\mu=\xbar,\qquad\widehat\sigma^{2}=\frac1n\sum_{i=1}^{n}(x_i-\xbar)^{2}.
\]
\end{teorema}

\begin{proof}
Con $v=\sigma^{2}>0$, la log-verosimilitud es
\[
\ell(\mu,v)=-\frac n2\ln(2\pi)-\frac n2\ln v-\frac{1}{2v}\sum_{i=1}^{n}(x_i-\mu)^{2}.
\]
\emph{Paso 1 ($\mu$).} Para cada $v$ fijo, maximizar $\ell$ en $\mu$ equivale a minimizar $\sum_i(x_i-\mu)^{2}$, cuyo mínimo está en $\mu=\xbar$ (\cref{prop:t1-extremal}). Por tanto $\ell(\mu,v)\le\ell(\xbar,v)$ para todo $(\mu,v)$.

\emph{Paso 2 ($v$).} Sea $Q=\sum_i(x_i-\xbar)^{2}>0$ y $h(v)=\ell(\xbar,v)=c-\frac n2\ln v-\frac{Q}{2v}$. Entonces $h'(v)=-\frac{n}{2v}+\frac{Q}{2v^{2}}=\frac{Q-nv}{2v^{2}}$, positiva para $v<Q/n$ y negativa para $v>Q/n$. Luego $h$ alcanza su máximo en $v=Q/n$.

\emph{Paso 3.} Combinando: $\ell(\mu,v)\le\ell(\xbar,v)\le\ell(\xbar,Q/n)$ para todo $(\mu,v)$.
\end{proof}

El EMV de $\sigma^{2}$ coincide con el estimador por momentos y es sesgado (\cref{cor:t2-sesgo}); por eso, para datos, se usa casi siempre $S^{2}$.

\begin{ejemplo}[Un caso sin derivadas]\label{ej:t4-uniforme}
Si $x_1,\dots,x_n$ son observaciones de $\Unif(0,\theta)$, entonces $f(x\mid\theta)=1/\theta$ si $0\le x\le\theta$ y $0$ en otro caso, luego
\[
\Lik(\theta)=\begin{cases}\theta^{-n}&\text{si }\theta\ge\max_ix_i,\\0&\text{si }\theta<\max_ix_i.\end{cases}
\]
$\Lik$ es decreciente en $[\max_ix_i,\infty)$, de modo que $\widehat\theta=\max_ix_i$. La ecuación $\ell'(\theta)=0$ no tiene solución: el máximo está en un extremo del dominio.
\end{ejemplo}

\begin{observacion}[Propiedades del EMV]\label{obs:t4-propiedades}
Se enuncian sin demostración, como ampliación~\cite[cap.~9]{wasserman}: (a) \emph{invariancia}: si $\widehat\theta$ es el EMV de $\theta$, el EMV de $g(\theta)$ es $g(\widehat\theta)$; por ejemplo, el EMV de $\sigma$ es $\sqrt{\widehat\sigma^{2}}$; (b) en condiciones de regularidad (que no se cumplen en el \cref{ej:t4-uniforme}), el EMV es consistente y aproximadamente normal para $n$ grande, con un error estándar que se puede estimar a partir de la curvatura de $\ell$ en el máximo.
\end{observacion}

\section{Intervalos de confianza}

Una estimación puntual no informa de su precisión. Un intervalo de confianza acompaña la estimación de un margen de error con una garantía probabilística.

\begin{definicion}[Intervalo de confianza]\label{def:t4-ic}
Sea $0<\alpha<1$. Un \emph{intervalo de confianza} de nivel $1-\alpha$ para $\theta$ es un intervalo $[L,U]$ cuyos extremos $L=L(X_1,\dots,X_n)$ y $U=U(X_1,\dots,X_n)$ son estadísticos y tal que
\[
\PP_\theta(L\le\theta\le U)\ge1-\alpha\qquad\text{para todo }\theta\in\Theta.
\]
La probabilidad $\PP_\theta(L\le\theta\le U)$ se llama \emph{cobertura}~\cite[cap.~6]{wasserman}.
\end{definicion}

\begin{observacion}[Interpretación]\label{obs:t4-interpretacion}
Lo aleatorio es el intervalo, no el parámetro. Si se repitiera muchas veces el experimento y se calculara cada vez el intervalo de nivel $95\,\%$, aproximadamente el $95\,\%$ de esos intervalos contendría el verdadero $\theta$. Una vez observados los datos, el intervalo concreto (por ejemplo, $[0{,}79;\allowbreak\,0{,}84]$) contiene o no contiene a $\theta$; no es correcto decir que $\theta$ está en él «con probabilidad $0{,}95$» en el sentido frecuentista~\cite[lectura~23a]{mit1805}.
\end{observacion}

El método general para construir intervalos usa un \emph{pivote}: una función $Q(X_1,\dots,X_n;\theta)$ cuya distribución no depende de $\theta$.

\begin{teorema}[Intervalo para la media con $\sigma$ conocida]\label{teo:t4-ic-z}
Sean $X_1,\dots,X_n\sim\Nor(\mu,\sigma^{2})$ \iid\ con $\sigma$ conocida. El intervalo
\[
\Bigl[\Xbar-z_{1-\alpha/2}\frac{\sigma}{\sqrt n},\ \Xbar+z_{1-\alpha/2}\frac{\sigma}{\sqrt n}\Bigr]
\]
tiene cobertura exactamente $1-\alpha$. Si los datos no son normales pero $n$ es grande, la cobertura es aproximadamente $1-\alpha$.
\end{teorema}

\begin{proof}
El pivote $Z=(\Xbar-\mu)/(\sigma/\sqrt n)$ tiene distribución $\Nor(0,1)$ para todo $\mu$ (\cref{teo:t3-suma-normales,prop:t3-tipificacion}). Con $c=z_{1-\alpha/2}$, $\PP(-c\le Z\le c)=\Phi(c)-\Phi(-c)=2\Phi(c)-1=1-\alpha$. Despejando $\mu$ en las desigualdades,
\[
-c\le\frac{\Xbar-\mu}{\sigma/\sqrt n}\le c\iff\Xbar-c\frac{\sigma}{\sqrt n}\le\mu\le\Xbar+c\frac{\sigma}{\sqrt n}.
\]
Los dos sucesos son el mismo y tienen la misma probabilidad. Para datos no normales se aplica el mismo argumento con la aproximación del \cref{teo:t3-tcl}.
\end{proof}

La cantidad $E=z_{1-\alpha/2}\,\sigma/\sqrt n$ es el \emph{margen de error} (la semiamplitud del intervalo). Los valores más usados son $z_{0{,}95}=1{,}645$ (nivel $90\,\%$), $z_{0{,}975}=1{,}960$ ($95\,\%$) y $z_{0{,}995}=2{,}576$ ($99\,\%$).

\begin{corolario}[Tamaño muestral]\label{cor:t4-tamano}
Para que el margen de error del \cref{teo:t4-ic-z} sea a lo sumo $E$ basta tomar
\[
n\ge\Bigl(\frac{z_{1-\alpha/2}\,\sigma}{E}\Bigr)^{2},
\]
redondeando al entero superior.
\end{corolario}

\begin{proof}
$z_{1-\alpha/2}\sigma/\sqrt n\le E\iff\sqrt n\ge z_{1-\alpha/2}\sigma/E$, y se elevan al cuadrado ambos miembros, que son positivos.
\end{proof}

\section{Intervalo para la media con \texorpdfstring{$\sigma$}{sigma} desconocida}

En la práctica $\sigma$ casi nunca se conoce. Al sustituirla por $S$ hay que reemplazar la normal por la $t$ de Student.

\begin{teorema}[Intervalo $t$]\label{teo:t4-ic-t}
Sean $X_1,\dots,X_n\sim\Nor(\mu,\sigma^{2})$ \iid\ con $\sigma$ desconocida y $n\ge2$. El intervalo
\[
\Bigl[\Xbar-t_{n-1;1-\alpha/2}\frac{S}{\sqrt n},\ \Xbar+t_{n-1;1-\alpha/2}\frac{S}{\sqrt n}\Bigr]
\]
tiene cobertura exactamente $1-\alpha$.
\end{teorema}

\begin{proof}
Por el \cref{cor:t3-t}, el pivote $T=(\Xbar-\mu)/(S/\sqrt n)$ tiene distribución $t_{n-1}$, que no depende de $\mu$ ni de $\sigma$. Con $c=t_{n-1;1-\alpha/2}$, la simetría de la $t$ (\cref{eje:t3-t-simetrica}) da $\PP(-c\le T\le c)=1-\alpha$, y se despeja $\mu$ como en el \cref{teo:t4-ic-z}.
\end{proof}

Como $t_{n-1;1-\alpha/2}>z_{1-\alpha/2}$, el intervalo $t$ es más ancho: refleja la incertidumbre adicional por estimar $\sigma$. Para $n$ grande (digamos $n\ge30$ o más) ambos cuantiles casi coinciden y, por el teorema central del límite, el intervalo es aproximadamente válido aunque los datos no sean normales.

\begin{ejemplo}\label{ej:t4-t}
Se miden los tiempos de respuesta (en segundos) de un sistema en $n=10$ consultas y se obtiene $\xbar=2{,}40$ y $s=0{,}50$. Con $t_{9;0{,}975}=2{,}262$, el intervalo de nivel $95\,\%$ es $2{,}40\pm2{,}262\cdot0{,}50/\sqrt{10}=2{,}40\pm0{,}358$, es decir, $[2{,}04;\allowbreak\ 2{,}76]$. Si se hubiera usado $1{,}96$, el margen habría sido $0{,}310$, un $13\,\%$ menor.
\end{ejemplo}

\section{Intervalos para proporciones}

Si $X_1,\dots,X_n\sim\Ber(p)$, el \cref{cor:t3-binomial} da $\widehat p\approx\Nor\bigl(p,p(1-p)/n\bigr)$. El error estándar depende de $p$, desconocido, y se estima por $\sqrt{\widehat p(1-\widehat p)/n}$.

\begin{definicion}[Intervalo de Wald para una proporción]\label{def:t4-wald}
El intervalo de confianza aproximado de nivel $1-\alpha$ para $p$ es
\[
\widehat p\pm z_{1-\alpha/2}\sqrt{\frac{\widehat p(1-\widehat p)}{n}}.
\]
Como $p(1-p)\le1/4$, el intervalo \emph{conservador} $\widehat p\pm z_{1-\alpha/2}/(2\sqrt n)$ es siempre al menos tan ancho~\cite[§8.3]{openstax}.
\end{definicion}

La aproximación es buena si $n\widehat p$ y $n(1-\widehat p)$ son ambos al menos del orden de $10$; con muestras pequeñas o con $\widehat p$ cercano a $0$ o a $1$ la cobertura real es inferior a la nominal. Existen correcciones sencillas, como el método de «más cuatro» de~\cite[§8.3]{openstax} (sumar dos éxitos y dos fracasos antes de calcular) o el intervalo de Wilson, que quedan fuera de este módulo.

\begin{corolario}[Tamaño muestral para una proporción]\label{cor:t4-tamano-p}
Para que el margen de error sea a lo sumo $E$ sea cual sea $p$, basta $n\ge z_{1-\alpha/2}^{2}/(4E^{2})$. Si se dispone de una estimación previa $p^{*}$, basta $n\ge z_{1-\alpha/2}^{2}\,p^{*}(1-p^{*})/E^{2}$.
\end{corolario}

\begin{proof}
Es el \cref{cor:t4-tamano} con $\sigma=\sqrt{p(1-p)}$, acotado por $1/2$ o sustituido por $\sqrt{p^{*}(1-p^{*})}$.
\end{proof}

\begin{ejemplo}[Exactitud en un benchmark]\label{ej:t4-benchmark}
Un modelo acierta $820$ de $1000$ preguntas. Entonces $\widehat p=0{,}82$, $\widehat{\ee}=\sqrt{0{,}82\cdot0{,}18/1000}\approx0{,}01215$ y el intervalo de nivel $95\,\%$ es $0{,}82\pm1{,}96\cdot0{,}01215=0{,}82\pm0{,}0238$, es decir, $[0{,}796;\allowbreak\ 0{,}844]$. Para reducir el margen a un punto porcentual ($E=0{,}01$) harían falta $1{,}96^{2}\cdot0{,}82\cdot0{,}18/0{,}01^{2}\approx5671$ preguntas.
\end{ejemplo}

\begin{observacion}[Errores frecuentes]\label{obs:t4-errores}
(a) Interpretar un intervalo de confianza del $95\,\%$ como «hay un $95\,\%$ de probabilidad de que $\theta$ esté en este intervalo concreto» o como «el $95\,\%$ de los datos está en el intervalo» (\cref{obs:t4-interpretacion}).
(b) Usar $z$ en lugar de $t$ con $n$ pequeño y $\sigma$ estimada.
(c) Olvidar que el margen de error solo recoge el error de muestreo: no protege contra sesgos (por ejemplo, preguntas del benchmark que estaban en los datos de entrenamiento).
(d) Para comparar dos modelos, mirar solo si sus intervalos se solapan; dos intervalos que se solapan ligeramente pueden corresponder a una diferencia significativa (\cref{tema:5}).
(e) En el EMV, derivar sin comprobar que el punto crítico es un máximo, o derivar cuando el soporte depende del parámetro (\cref{ej:t4-uniforme}).
\end{observacion}

\begin{observacion}[Conexión con el curso]
El módulo~1 obtiene el EMV de un modelo de bigramas, que generaliza el \cref{teo:t4-emv-bernoulli} a distribuciones categóricas: la frecuencia relativa. Entrenar un modelo de lenguaje minimizando la entropía cruzada es exactamente maximizar la log-verosimilitud del corpus (módulos~1, 5 y~8). En los módulos~22 y~23 se construyen intervalos de confianza para la exactitud en un benchmark (\cref{ej:t4-benchmark}) y la regla del~3 (\cref{eje:t4-regla3}) para acotar la tasa de fallos graves cuando no se ha observado ninguno; en el módulo~16, los exponentes de las leyes de escala se estiman con su incertidumbre.
\end{observacion}

\begin{lecturas}
\item Wasserman~\cite{wasserman}, cap.~9 «Parametric Inference» (método de los momentos y máxima verosimilitud) y cap.~6 (intervalos de confianza).
\item Stats~200~\cite{stats200}, clases~12 «Parametric models, method of moments estimation», 13 «Maximum likelihood estimation» y 18 «Confidence intervals».
\item Orloff y Bloom~\cite{mit1805}, lecturas~10b «Maximum Likelihood Estimates», 22 «Confidence intervals based on normal data», 23a «Confidence Intervals: Three Views» y 23b «Confidence Intervals for the Mean of Non-normal Data».
\item OpenStax~\cite{openstax}, cap.~8: §8.1 «A Single Population Mean using the Normal Distribution», §8.2 «A Single Population Mean using the Student t Distribution» y §8.3 «A Population Proportion».
\item Rice~\cite{rice}, cap.~8 «Estimation of Parameters and Fitting of Probability Distributions».
\end{lecturas}

\section*{Ejercicios del tema 4}
\addcontentsline{toc}{section}{Ejercicios del tema 4}

\begin{ejercicio}{1}{eje:t4-interpretacion}
Un estudio concluye que un nuevo método de ajuste fino mejora la exactitud media en $15$ puntos, con un intervalo de confianza del $95\,\%$ de $[10,20]$ puntos. Indique cuáles de estas afirmaciones se deducen del intervalo: (i) hay una probabilidad del $95\,\%$ de que la mejora real esté entre $10$ y $20$; (ii) la probabilidad de que el método mejore la exactitud es al menos del $95\,\%$; (iii) el intervalo da una estimación de la mejora real; (iv) la hipótesis de que no hay mejora es probablemente falsa; (v) si se repitiera el estudio muchas veces, alrededor del $95\,\%$ de los intervalos contendría la mejora real; (vi) un contraste bilateral al $5\,\%$ rechazaría la hipótesis de mejora nula.
\fuente{adaptado de Orloff y Bloom, MIT 18.05~\cite{mit1805}, \emph{problem set}~10, problema~5(c).}
\end{ejercicio}

\begin{ejercicio}{2}{eje:t4-momentos-basico}
(a) Los tiempos entre peticiones a un servidor, en segundos, son $0{,}5;\allowbreak\ 1{,}2;\allowbreak\ 0{,}3;\allowbreak\ 2{,}0$. Suponiendo un modelo exponencial, estime $\lambda$ por el método de los momentos. (b) Para los datos $2,4,4,5,7$ y un modelo normal, obtenga las estimaciones de $\mu$ y $\sigma^{2}$ por el método de los momentos y compárelas con $s^{2}$.
\end{ejercicio}

\begin{ejercicio}{2}{eje:t4-geometrica}
Sean $X_1,\dots,X_n$ \iid\ con distribución geométrica en $\{1,2,3,\dots\}$, $\PP(X=x)=p(1-p)^{x-1}$, cuya media es $1/p$. (a) Obtenga el estimador de $p$ por el método de los momentos. (b) Obtenga el EMV. (c) Calcule ambos con los datos $2,1,4,3,1,1$ (número de intentos hasta que un agente resuelve una tarea).
\fuente{adaptado de Stats~200, Stanford~\cite{stats200}, \emph{homework}~5, problema~1.}
\end{ejercicio}

\begin{ejercicio}{2}{eje:t4-gamma}
Una variable positiva tiene densidad $f(x\mid a)=\dfrac{a^{5}}{4!}\,x^{4}e^{-ax}$ para $x>0$, con $a>0$. (a) Obtenga el EMV de $a$ a partir de $n$ observaciones \iid. (b) Calcúlelo para los datos $2{,}0;\allowbreak\ 3{,}0;\allowbreak\ 2{,}5;\allowbreak\ 4{,}5$.
\fuente{adaptado de Orloff y Bloom, MIT 18.05~\cite{mit1805}, \emph{problem set}~6, problema~1.}
\end{ejercicio}

\begin{ejercicio}{2}{eje:t4-uniforme-ab}
Los datos $1{,}2;\allowbreak\ 2{,}1;\allowbreak\ 1{,}3;\allowbreak\ 10{,}5;\allowbreak\ 5$ proceden de una distribución $\Unif(a,b)$. (a) Obtenga el EMV de $a$ y $b$ razonando directamente sobre la verosimilitud, sin derivar. (b) Obtenga los estimadores por el método de los momentos, sabiendo que la media es $(a+b)/2$ y la varianza $(b-a)^{2}/12$. (c) ¿Son compatibles con los datos los estimadores de (b)?
\fuente{adaptado de Orloff y Bloom, MIT 18.05~\cite{mit1805}, \emph{problem set}~6, problema~3.}
\end{ejercicio}

\begin{ejercicio}{2}{eje:t4-poisson-datos}
El número de alucinaciones detectadas en ocho documentos generados es $2,0,3,1,4,2,1,3$. Suponiendo un modelo de Poisson, (a) estime $\lambda$; (b) construya un intervalo aproximado del $95\,\%$ usando que $\Var(\Xbar)=\lambda/n$.
\end{ejercicio}

\begin{ejercicio}{2}{eje:t4-soda}
Una máquina llena latas con volúmenes $\Nor(\mu,\sigma^{2})$. Se miden nueve latas (en mL): $352,351,361,353,352,358,360,358,359$. (a) Con $\sigma=3$ conocida, construya los intervalos de confianza del $95\,\%$ y del $98\,\%$ para $\mu$. (b) Repita el apartado suponiendo $\sigma$ desconocida y compare.
\fuente{adaptado de Orloff y Bloom, MIT 18.05~\cite{mit1805}, \emph{problem set}~10, problema~2.}
\end{ejercicio}

\begin{ejercicio}{2}{eje:t4-alturas}
Se mide la estatura de $20$ personas, que se suponen normales, y se obtiene $\xbar=69{,}55$ y $s^{2}=14{,}26$ (en pulgadas y pulgadas cuadradas). (a) Construya un intervalo $t$ del $90\,\%$ para $\mu$. (b) Suponiendo conocida $\sigma=3{,}77$, construya el intervalo $z$ del $90\,\%$. (c) En la situación de (b), ¿cuántas personas hay que medir para que el intervalo del $90\,\%$ tenga una amplitud total de una pulgada?
\fuente{adaptado de Orloff y Bloom, MIT 18.05~\cite{mit1805}, \emph{problem set}~11, problema~1.}
\end{ejercicio}

\begin{ejercicio}{2}{eje:t4-openstax-z}
(a) Una población tiene $\sigma=3$; una muestra de $n=36$ da $\xbar=68$. Construya el intervalo del $90\,\%$ para $\mu$. (b) Si $\sigma=15$, ¿qué tamaño muestral hace falta para que el margen de error del $95\,\%$ sea $2$?
\fuente{adaptado de OpenStax, \emph{Introductory Statistics}~2e~\cite{openstax}, §8.1, ejemplos~8.2 y~8.7.}
\end{ejercicio}

\begin{ejercicio}{2}{eje:t4-proporcion}
(a) En una encuesta a $500$ personas, $421$ declaran usar un asistente de IA a diario. Construya el intervalo del $95\,\%$ para la proporción. (b) ¿Qué tamaño muestral garantiza un margen de error del $3\,\%$ con confianza del $90\,\%$, sin información previa sobre $p$?
\fuente{adaptado de OpenStax, \emph{Introductory Statistics}~2e~\cite{openstax}, §8.3, ejemplos~8.10 y~8.14.}
\end{ejercicio}

\begin{ejercicio}{2}{eje:t4-margen}
Se usa el intervalo conservador $\widehat p\pm z_{1-\alpha/2}/(2\sqrt n)$. (a) ¿Cuántas preguntas debe tener un benchmark para que el margen sea $0{,}01$ al $95\,\%$? Hágalo con $z=1{,}96$ y con la aproximación $z\approx2$. (b) ¿Y al $80\,\%$? (c) Una encuesta anuncia un margen de error de $\pm4$ puntos al $95\,\%$; estime cuántas personas se encuestaron.
\fuente{adaptado de Orloff y Bloom, MIT 18.05~\cite{mit1805}, problemas de la clase~23, pregunta conceptual~1 y problema~1.}
\end{ejercicio}

\begin{ejercicio}{3}{eje:t4-emv-poisson}
Sean $x_1,\dots,x_n$ observaciones \iid\ de $\Pois(\lambda)$ con $\sum_ix_i>0$. Demuestre que el EMV de $\lambda$ es $\xbar$, comprobando que es un máximo global, y que coincide con el estimador por momentos.
\end{ejercicio}

\begin{ejercicio}{3}{eje:t4-emv-expo}
Sean $x_1,\dots,x_n>0$ observaciones \iid\ de $\Expo(\lambda)$, con densidad $\lambda e^{-\lambda x}$. (a) Demuestre que el EMV de $\lambda$ es $1/\xbar$. (b) Reparametrice el modelo por la media $\beta=1/\lambda$ y demuestre directamente que el EMV de $\beta$ es $\xbar$, de acuerdo con la propiedad de invariancia.
\end{ejercicio}

\begin{ejercicio}{3}{eje:t4-laplace}
La distribución de Laplace con parámetro de posición $\mu$ tiene densidad $f(x\mid\mu)=\frac12e^{-\abs{x-\mu}}$. Demuestre que el EMV de $\mu$ es la mediana muestral. (\emph{Indicación:} \cref{eje:t1-mediana-min}.) Compare con el modelo normal, cuyo EMV de la media es $\xbar$.
\end{ejercicio}

\begin{ejercicio}{3}{eje:t4-regla3}
Se prueban $n$ entradas adversarias independientes contra un modelo y ninguna provoca un fallo grave. Sea $p$ la probabilidad de fallo en cada entrada. (a) Demuestre que el conjunto de valores de $p$ para los que observar cero fallos tiene probabilidad al menos $\alpha$ es $[0,p_{\max}]$ con $p_{\max}=1-\alpha^{1/n}$. (b) Demuestre que, para $\alpha=0{,}05$, $p_{\max}\le-\ln(0{,}05)/n\approx3/n$ («regla del~3»). (c) Calcule $p_{\max}$ y $3/n$ para $n=300$.
\end{ejercicio}

\begin{ejercicio}{3}{eje:t4-conservador}
Sea $\widehat p$ la proporción muestral de $n$ ensayos de Bernoulli. Usando la aproximación normal $\widehat p\approx\Nor(p,p(1-p)/n)$, demuestre que el intervalo conservador $\widehat p\pm z_{1-\alpha/2}/(2\sqrt n)$ tiene cobertura aproximada al menos $1-\alpha$ para todo $p\in(0,1)$, con igualdad solo si $p=1/2$. ¿Qué ocurre con la cobertura aproximada cuando $p=0{,}1$, $n=100$ y $\alpha=0{,}05$?
\end{ejercicio}

% =====================================================================
\chapter{Contrastes de hipótesis}
\label{tema:5}
% =====================================================================

Un contraste de hipótesis decide, a partir de los datos, si hay evidencia suficiente contra una afirmación sobre la población: que una moneda es equilibrada, que un nuevo modelo no mejora al anterior o que unos recuentos siguen una distribución dada. Este tema presenta el marco de Neyman y Pearson (hipótesis nula y alternativa, región de rechazo, errores de tipo~I y~II, potencia), el $p$-valor, los contrastes $z$ y $t$ para medias y proporciones, el contraste $\chi^{2}$ de bondad de ajuste y el problema de las comparaciones múltiples, con su aplicación a la evaluación de modelos de lenguaje en muchos benchmarks. Se sigue el capítulo~10 de Wasserman~\cite{wasserman}, «Hypothesis Testing and p-values», las lecturas~17b, 18 y~19 del curso 18.05 del MIT~\cite{mit1805}, las clases~5, 7 y~11 de Stats~200~\cite{stats200} y los capítulos~9 y~11 de OpenStax~\cite{openstax}.

\section{Planteamiento de un contraste}

\begin{definicion}[Hipótesis, estadístico y región de rechazo]\label{def:t5-contraste}
Sea $\{F_\theta:\theta\in\Theta\}$ un modelo y $\Theta_0\subset\Theta$. Se contrastan la \emph{hipótesis nula} $H_0:\theta\in\Theta_0$ frente a la \emph{hipótesis alternativa} $H_1:\theta\in\Theta\setminus\Theta_0$. Una hipótesis es \emph{simple} si especifica un único valor de $\theta$ y \emph{compuesta} en otro caso. Un \emph{contraste} consiste en un \emph{estadístico de contraste} $T=T(X_1,\dots,X_n)$ y una \emph{región de rechazo} $R$: si $T\in R$ se rechaza $H_0$; si no, no se rechaza.
\end{definicion}

La hipótesis nula representa la situación por defecto («no hay efecto», «la moneda es equilibrada», «los dos modelos son igual de buenos») y el contraste busca evidencia contra ella. Las alternativas más comunes para un parámetro real son la \emph{bilateral} $H_1:\theta\neq\theta_0$ y las \emph{unilaterales} $H_1:\theta>\theta_0$ o $H_1:\theta<\theta_0$.

\begin{definicion}[Errores, nivel y potencia]\label{def:t5-errores}
Se comete un \emph{error de tipo~I} si se rechaza $H_0$ siendo cierta, y un \emph{error de tipo~II} si no se rechaza siendo falsa:
\begin{center}
\small
\begin{tabular}{@{}lcc@{}}
\toprule
 & $H_0$ cierta & $H_0$ falsa \\
\midrule
No se rechaza $H_0$ & decisión correcta & error de tipo~II \\
Se rechaza $H_0$ & error de tipo~I & decisión correcta \\
\bottomrule
\end{tabular}
\end{center}
El contraste tiene \emph{nivel de significación} $\alpha$ si $\PP_\theta(T\in R)\le\alpha$ para todo $\theta\in\Theta_0$. La \emph{función de potencia} es $\pi(\theta)=\PP_\theta(T\in R)$; para $\theta\in\Theta\setminus\Theta_0$, $\pi(\theta)$ es la \emph{potencia} contra esa alternativa y $\beta(\theta)=1-\pi(\theta)$ es la probabilidad de error de tipo~II~\cite[cap.~10]{wasserman}.
\end{definicion}

Los dos errores no se pueden reducir a la vez con un tamaño muestral fijo: agrandar la región de rechazo reduce los errores de tipo~II y aumenta los de tipo~I. El convenio es fijar $\alpha$ (habitualmente $0{,}05$ o $0{,}01$) y, entre los contrastes de nivel $\alpha$, preferir los más potentes. Las dos hipótesis no son simétricas: no rechazar $H_0$ no prueba que sea cierta, solo que los datos no aportan evidencia suficiente contra ella.

\section{El \texorpdfstring{$p$}{p}-valor}

\begin{definicion}[$p$-valor]\label{def:t5-pvalor}
Supóngase que se rechaza $H_0$ para valores grandes de un estadístico $T$ y que $H_0$ es simple. Si el valor observado es $t_{\mathrm{obs}}$, el \emph{$p$-valor} es
\[
p=\PP_{H_0}(T\ge t_{\mathrm{obs}}),
\]
la probabilidad, calculada suponiendo $H_0$ cierta, de obtener un resultado al menos tan extremo como el observado. (Si se rechaza para valores pequeños, $p=\PP_{H_0}(T\le t_{\mathrm{obs}})$; en un contraste bilateral con estadístico de distribución simétrica, $p=\PP_{H_0}(\abs T\ge\abs{t_{\mathrm{obs}}})$.)
\end{definicion}

\begin{proposicion}[Decisión con el $p$-valor]\label{prop:t5-pvalor}
Sea un contraste que rechaza $H_0$ si $T\ge c_\alpha$, donde $c_\alpha$ cumple $\PP_{H_0}(T\ge c_\alpha)=\alpha$ y $T$ tiene distribución continua bajo $H_0$. Entonces se rechaza $H_0$ al nivel $\alpha$ si y solo si $p\le\alpha$.
\end{proposicion}

\begin{proof}
Sea $G(t)=\PP_{H_0}(T\ge t)$, que es continua y no creciente, de modo que $p=G(t_{\mathrm{obs}})$ y $G(c_\alpha)=\alpha$. Si $t_{\mathrm{obs}}\ge c_\alpha$, entonces $p=G(t_{\mathrm{obs}})\le G(c_\alpha)=\alpha$. Recíprocamente, si $t_{\mathrm{obs}}<c_\alpha$, entonces $p=\PP_{H_0}(T\ge t_{\mathrm{obs}})=\PP_{H_0}(t_{\mathrm{obs}}\le T<c_\alpha)+\alpha$, y el primer sumando es positivo si la densidad de $T$ no se anula entre $t_{\mathrm{obs}}$ y $c_\alpha$, de modo que $p>\alpha$. (Si la densidad se anulara en ese tramo, se toma como $c_\alpha$ el menor valor con $G(c_\alpha)=\alpha$ y el argumento es el mismo.)
\end{proof}

El $p$-valor es, por tanto, el menor nivel al que los datos permiten rechazar $H_0$, y resume en un número la fuerza de la evidencia contra $H_0$: cuanto menor, más incompatibles son los datos con $H_0$.

\begin{observacion}[Lo que el $p$-valor no es]\label{obs:t5-pvalor}
El $p$-valor \emph{no} es la probabilidad de que $H_0$ sea cierta, ni la probabilidad de que los resultados se deban al azar, ni mide el tamaño del efecto. Un $p$-valor muy pequeño puede acompañar a un efecto irrelevante si $n$ es enorme, y un efecto grande puede dar un $p$-valor alto si $n$ es pequeño. Por eso conviene informar siempre del efecto estimado con su intervalo de confianza, además del $p$-valor~\cite[lecturas~17b y~20]{mit1805}.
\end{observacion}

\section{Contrastes \texorpdfstring{$z$}{z} y \texorpdfstring{$t$}{t}}

Los contrastes de esta sección usan como estadístico una estimación tipificada con su error estándar calculado bajo $H_0$. En todos ellos, las regiones de rechazo de nivel $\alpha$ son $T>q_{1-\alpha}$ (alternativa $>$), $T<-q_{1-\alpha}$ (alternativa $<$) y $\abs T>q_{1-\alpha/2}$ (bilateral), donde $q$ es el cuantil de la distribución de referencia.

\begin{center}
\small
\begin{tabular}{@{}llll@{}}
\toprule
Contraste & $H_0$ & Estadístico & Distribución bajo $H_0$ \\
\midrule
$z$ para la media ($\sigma$ conocida) & $\mu=\mu_0$ & $Z=\dfrac{\Xbar-\mu_0}{\sigma/\sqrt n}$ & $\Nor(0,1)$ (datos normales) \\[2ex]
$t$ para la media & $\mu=\mu_0$ & $T=\dfrac{\Xbar-\mu_0}{S/\sqrt n}$ & $t_{n-1}$ (datos normales) \\[2ex]
$z$ para una proporción & $p=p_0$ & $Z=\dfrac{\widehat p-p_0}{\sqrt{p_0(1-p_0)/n}}$ & $\approx\Nor(0,1)$ ($n$ grande) \\
\bottomrule
\end{tabular}
\end{center}

La distribución bajo $H_0$ se sigue del \cref{teo:t3-suma-normales}, del \cref{cor:t3-t} y del \cref{cor:t3-binomial}, respectivamente. En el contraste de una proporción, el error estándar se calcula con $p_0$, no con $\widehat p$, porque se razona bajo $H_0$.

\begin{ejemplo}[Contraste $z$ unilateral]\label{ej:t5-z}
Se afirma que un modelo genera respuestas de $\mu_0=200$ tokens de media, con $\sigma=80$. En $n=64$ respuestas se observa $\xbar=221$. Se contrasta $H_0:\mu=200$ frente a $H_1:\mu>200$ al nivel $0{,}05$. El estadístico vale $z=(221-200)/(80/8)=2{,}1>z_{0{,}95}=1{,}645$, y $p=1-\Phi(2{,}1)=0{,}0179<0{,}05$: se rechaza $H_0$.
\end{ejemplo}

\section{Potencia y tamaño muestral}

\begin{proposicion}[Potencia del contraste $z$ unilateral]\label{prop:t5-potencia}
Para el contraste $z$ de $H_0:\mu=\mu_0$ frente a $H_1:\mu>\mu_0$ al nivel $\alpha$ con datos $\Nor(\mu,\sigma^{2})$, la función de potencia es
\[
\pi(\mu)=1-\Phi\Bigl(z_{1-\alpha}-\frac{(\mu-\mu_0)\sqrt n}{\sigma}\Bigr).
\]
\end{proposicion}

\begin{proof}
Se rechaza si $\Xbar>\mu_0+z_{1-\alpha}\sigma/\sqrt n$. Si la media verdadera es $\mu$, entonces $(\Xbar-\mu)/(\sigma/\sqrt n)\sim\Nor(0,1)$, y
\[
\pi(\mu)=\PP_\mu\Bigl(\frac{\Xbar-\mu}{\sigma/\sqrt n}>\frac{\mu_0-\mu}{\sigma/\sqrt n}+z_{1-\alpha}\Bigr)=1-\Phi\Bigl(z_{1-\alpha}-\frac{(\mu-\mu_0)\sqrt n}{\sigma}\Bigr).\qedhere
\]
\end{proof}

En el \cref{ej:t5-z}, si la media real fuera $220$, la potencia sería $1-\Phi(1{,}645-2)=\Phi(0{,}355)\approx0{,}64$: con $64$ respuestas, el contraste detectaría esa diferencia solo en dos de cada tres ocasiones. El \cref{eje:t5-potencia} deduce el tamaño muestral necesario para alcanzar una potencia dada.

\section{Contraste \texorpdfstring{$\chi^{2}$}{chi cuadrado} de bondad de ajuste}

Sea una variable categórica con $k$ categorías y probabilidades $p_1,\dots,p_k$. En $n$ observaciones independientes se obtienen los recuentos $O_1,\dots,O_k$, con $\sum_jO_j=n$. Se quiere contrastar $H_0:p_j=p_{0j}$ para todo $j$, con $p_{0j}$ dados.

\begin{definicion}[Estadístico $\chi^{2}$ de Pearson]\label{def:t5-chi2}
Sean $E_j=np_{0j}$ los recuentos \emph{esperados} bajo $H_0$. El estadístico es
\[
X^{2}=\sum_{j=1}^{k}\frac{(O_j-E_j)^{2}}{E_j}.
\]
Bajo $H_0$, y si todos los $E_j$ son al menos~$5$, $X^{2}$ sigue aproximadamente una distribución $\chi^{2}_{k-1}$. Se rechaza $H_0$ al nivel $\alpha$ si $X^{2}>\chi^{2}_{k-1;1-\alpha}$~\cite[§11.2]{openstax}.
\end{definicion}

Los $k-1$ grados de libertad se deben a la restricción $\sum_jO_j=n$: conocidos $k-1$ recuentos, el último queda determinado. Si se estiman $m$ parámetros del modelo a partir de los mismos datos, se restan también: $k-1-m$ grados de libertad. Si algún $E_j$ es menor que $5$, se agrupan categorías contiguas.

\begin{ejemplo}[¿Es equilibrado un dado?]\label{ej:t5-dado}
En $120$ lanzamientos de un dado se obtienen los recuentos $25,17,15,23,24,16$ para las caras $1$ a $6$. Bajo $H_0$ (dado equilibrado), $E_j=20$ y
\[
X^{2}=\frac{25+9+25+9+16+16}{20}=\frac{100}{20}=5.
\]
Con $5$ grados de libertad, $\chi^{2}_{5;0{,}95}=11{,}07$; como $5<11{,}07$, no se rechaza $H_0$ al $5\,\%$ (el $p$-valor es aproximadamente $0{,}42$).
\end{ejemplo}

\section{Comparaciones múltiples}

Si se realizan muchos contrastes, cada uno al nivel $\alpha$, la probabilidad de cometer \emph{algún} error de tipo~I crece rápidamente.

\begin{proposicion}\label{prop:t5-multiples}
Si se realizan $m$ contrastes independientes al nivel $\alpha$ y todas las hipótesis nulas son ciertas, la probabilidad de rechazar al menos una es $1-(1-\alpha)^{m}$, y el número esperado de rechazos (todos falsos) es $m\alpha$.
\end{proposicion}

\begin{proof}
Cada contraste no rechaza con probabilidad $1-\alpha$; por independencia, ninguno rechaza con probabilidad $(1-\alpha)^{m}$. El número de rechazos es una suma de $m$ indicadores con esperanza $\alpha$ cada uno.
\end{proof}

Con $m=20$ y $\alpha=0{,}05$, la probabilidad de al menos un «descubrimiento» falso es $1-0{,}95^{20}\approx0{,}64$.

\begin{definicion}[Tasa de error por familia]\label{def:t5-fwer}
La \emph{tasa de error por familia} (FWER, \emph{family-wise error rate}) de un procedimiento que realiza $m$ contrastes es la probabilidad de cometer al menos un error de tipo~I~\cite[clase~11]{stats200}.
\end{definicion}

\begin{teorema}[Corrección de Bonferroni]\label{teo:t5-bonferroni}
Si cada uno de los $m$ contrastes se realiza al nivel $\alpha/m$ (es decir, se rechaza la hipótesis $i$ si $p_i\le\alpha/m$), la tasa de error por familia es a lo sumo $\alpha$, cualquiera que sea la dependencia entre los contrastes y cualquiera que sea el conjunto de hipótesis nulas ciertas.
\end{teorema}

\begin{proof}
Sea $I_0$ el conjunto de índices de las hipótesis nulas ciertas, con $m_0=\abs{I_0}\le m$, y sea $A_i$ el suceso «se rechaza $H_{0,i}$». Para $i\in I_0$, $\PP(A_i)\le\alpha/m$ porque el contraste $i$ tiene nivel $\alpha/m$. Por la desigualdad de la unión (subaditividad de la probabilidad),
\[
\text{FWER}=\PP\Bigl(\bigcup_{i\in I_0}A_i\Bigr)\le\sum_{i\in I_0}\PP(A_i)\le m_0\frac{\alpha}{m}\le\alpha.\qedhere
\]
\end{proof}

La corrección de Bonferroni es sencilla y muy general, pero conservadora: con $m$ grande reduce mucho la potencia. Existen procedimientos más potentes, como el de Holm o el de Benjamini y Hochberg, que controla una cantidad distinta, la tasa de falsos descubrimientos~\cite[clase~11]{stats200}; quedan fuera de este módulo.

\begin{observacion}[Comparar modelos en muchos benchmarks]\label{obs:t5-benchmarks}
La evaluación de modelos de lenguaje es un terreno abonado para las comparaciones múltiples: se comparan dos modelos en decenas de benchmarks, se prueban muchas variantes de \emph{prompt} y se informa de la mejor, o se ordena una tabla con cientos de modelos. Si en $20$ benchmarks independientes no hay ninguna diferencia real, se esperan $20\cdot0{,}05=1$ diferencias «significativas» al $5\,\%$. Para comparar la exactitud de dos modelos se puede usar un contraste $z$ de diferencia de proporciones (\cref{eje:t5-dos-modelos}); cuando ambos modelos responden a las \emph{mismas} preguntas, es preferible un contraste pareado que use las diferencias pregunta a pregunta, porque elimina la variabilidad debida a la dificultad de cada pregunta (ampliación).
\end{observacion}

\begin{observacion}[Errores frecuentes]\label{obs:t5-errores}
(a) Formular la hipótesis nula a partir de los datos ya vistos, o elegir el contraste unilateral después de ver en qué sentido va el efecto.
(b) Concluir «$H_0$ es cierta» porque no se rechaza.
(c) Interpretar el $p$-valor como $\PP(H_0\mid\text{datos})$ (\cref{obs:t5-pvalor}).
(d) En el contraste $\chi^{2}$, usar proporciones o porcentajes en lugar de recuentos, o no restar los parámetros estimados de los grados de libertad.
(e) Realizar muchos contrastes y destacar solo los significativos sin corregir por multiplicidad.
\end{observacion}

\begin{observacion}[Conexión con el curso]
Los módulos~22 y~23 comparan modelos con exactitudes medidas sobre conjuntos finitos de preguntas y discuten la significación de las diferencias, la contaminación de los benchmarks y su saturación; las ideas de este tema (error estándar, potencia, comparaciones múltiples) son las que permiten decidir si una mejora de un punto en una tabla de resultados es real. En el módulo~21, las preferencias humanas entre pares de respuestas se agregan con modelos probabilísticos cuyas diferencias también deben valorarse frente al ruido de muestreo.
\end{observacion}

\begin{lecturas}
\item Wasserman~\cite{wasserman}, cap.~10 «Hypothesis Testing and p-values» (incluye el contraste de Wald, el $\chi^{2}$ de Pearson y los contrastes múltiples).
\item Orloff y Bloom~\cite{mit1805}, lecturas~17b «Null Hypothesis Significance Testing I», 18 «NHST II», 19 «NHST III» y 20 «Comparison of frequentist and Bayesian inference».
\item Stats~200~\cite{stats200}, clases~5 «Testing a simple null hypothesis», 7 «Composite hypotheses and the t-test» y 11 «Testing multiple hypotheses».
\item OpenStax~\cite{openstax}, cap.~9: §9.1 «Null and Alternative Hypotheses», §9.2 «Outcomes and the Type I and Type II Errors» y §9.5 «Additional Information and Full Hypothesis Test Examples»; §11.2 «Goodness-of-Fit Test».
\item Rice~\cite{rice}, cap.~9 «Testing Hypotheses and Assessing Goodness of Fit».
\end{lecturas}

\section*{Ejercicios del tema 5}
\addcontentsline{toc}{section}{Ejercicios del tema 5}

\begin{ejercicio}{1}{eje:t5-vf}
Un contraste da $p=0{,}03$. Indique si cada afirmación es verdadera o falsa. (a) La probabilidad de que $H_0$ sea cierta es $0{,}03$. (b) Al nivel $0{,}05$ se rechaza $H_0$. (c) Al nivel $0{,}01$ se rechaza $H_0$. (d) Si $H_0$ fuera cierta, un resultado al menos tan extremo como el observado ocurriría con probabilidad $0{,}03$. (e) El efecto observado es grande. (f) Si no se hubiera rechazado $H_0$, quedaría demostrado que es cierta.
\end{ejercicio}

\begin{ejercicio}{2}{eje:t5-moneda}
Una moneda se hace girar sobre su canto $250$ veces y sale cara $140$ veces. Sea $\theta$ la probabilidad de cara y $H_0:\theta=0{,}5$ frente a $H_1:\theta\neq0{,}5$. (a) Calcule el $p$-valor con la aproximación normal y corrección por continuidad. (b) ¿Se rechaza $H_0$ al nivel $0{,}10$? ¿Y al $0{,}05$? (c) ¿Cuántas caras harían falta para rechazar al nivel $0{,}01$? (d) Con la región de rechazo de nivel $0{,}05$, estime la potencia contra $\theta=0{,}6$.
\fuente{adaptado de Orloff y Bloom, MIT 18.05~\cite{mit1805}, \emph{problem set}~8, problema~1.}
\end{ejercicio}

\begin{ejercicio}{2}{eje:t5-poligrafo}
En un estudio sobre detectores de mentiras, $140$ personas dicen la verdad y $140$ mienten. De las sinceras, $131$ son clasificadas como sinceras y $9$ como mentirosas; de las que mienten, $15$ son clasificadas como sinceras y $125$ como mentirosas. Con $H_0$: «la persona dice la verdad», describa los errores de tipo~I y~II y estime sus probabilidades, el nivel y la potencia del detector.
\fuente{adaptado de Orloff y Bloom, MIT 18.05~\cite{mit1805}, \emph{problem set}~8, problema~2.}
\end{ejercicio}

\begin{ejercicio}{2}{eje:t5-radar}
Tres radares miden la velocidad de cada coche con errores independientes $\Nor(0,5^{2})$ (en km/h) y se usa la media $\Xbar$ de las tres lecturas. El límite es $40$ y se quiere multar a lo sumo al $4\,\%$ de los conductores que circulan exactamente a $40$. (a) Plantee el contraste y halle el umbral de multa. (b) Si nadie superase el límite, ¿qué fracción de las multas sería errónea? (c) Calcule la potencia contra $\mu=45$. (d) ¿Cuántos radares harían falta para que esa potencia fuera $0{,}9$?
\fuente{adaptado de Orloff y Bloom, MIT 18.05~\cite{mit1805}, \emph{problem set}~8, problema~3.}
\end{ejercicio}

\begin{ejercicio}{2}{eje:t5-uniforme}
Se observa un único valor $x$ de una distribución $\Unif(0,\theta)$ para contrastar $H_0:\theta=2$ frente a $H_1:\theta\neq2$, y se rechaza $H_0$ si $x\le0{,}1$ o $x\ge1{,}9$. (a) Calcule la probabilidad de error de tipo~I. (b) Calcule la probabilidad de error de tipo~II si $\theta=2{,}5$.
\fuente{adaptado de Orloff y Bloom, MIT 18.05~\cite{mit1805}, \emph{problem set}~8, problema~4.}
\end{ejercicio}

\begin{ejercicio}{2}{eje:t5-t-test}
Una muestra normal de tamaño $n=16$ tiene $\xbar=11$ y $s^{2}=4$. Contraste $H_0:\mu=10$ al nivel $0{,}05$ (a) frente a $H_1:\mu\neq10$ y (b) frente a $H_1:\mu>10$. Dé en cada caso el valor crítico y acote el $p$-valor con la \cref{tab:t}.
\fuente{adaptado de Orloff y Bloom, MIT 18.05~\cite{mit1805}, \emph{problem set}~9, problema~1.}
\end{ejercicio}

\begin{ejercicio}{2}{eje:t5-openstax-z}
Un fabricante afirma que sus envases contienen de media $16{,}43$ onzas, con $\sigma=0{,}8$. Una muestra de $15$ envases da $\xbar=16$. Contraste al nivel $0{,}05$ si la media es menor que la anunciada, calculando el estadístico y el $p$-valor.
\fuente{adaptado de OpenStax, \emph{Introductory Statistics}~2e~\cite{openstax}, §9.5, ejemplo~9.14.}
\end{ejercicio}

\begin{ejercicio}{2}{eje:t5-openstax-prop}
Un tratamiento es eficaz, según su fabricante, en el $25\,\%$ de los casos o menos. En una muestra de $42$ casos es eficaz en $17$. Al nivel $0{,}01$, ¿hay evidencia de que la proporción de éxitos es mayor que $0{,}25$?
\fuente{adaptado de OpenStax, \emph{Introductory Statistics}~2e~\cite{openstax}, §9.5, ejemplo~9.19.}
\end{ejercicio}

\begin{ejercicio}{2}{eje:t5-chi-dias}
Las ausencias de $60$ empleados se reparten entre los días laborables así: lunes $15$, martes $12$, miércoles $9$, jueves $9$, viernes $15$. Contraste al nivel $0{,}05$ si las ausencias se reparten uniformemente entre los cinco días.
\fuente{adaptado de OpenStax, \emph{Introductory Statistics}~2e~\cite{openstax}, §11.2, ejemplo~11.2.}
\end{ejercicio}

\begin{ejercicio}{2}{eje:t5-benford}
Según la ley de Benford, el primer dígito $k$ de muchas cantidades contables tiene probabilidades $0{,}301;\allowbreak\ 0{,}176;\allowbreak\ 0{,}125;\allowbreak\ 0{,}097;\allowbreak\ 0{,}079;\allowbreak\ 0{,}067;\allowbreak\ 0{,}058;\allowbreak\ 0{,}051;\allowbreak\ 0{,}046$ para $k=1,\dots,9$. En $100$ partidas de una nota de gastos, los primeros dígitos aparecen $7,13,12,9,9,13,11,10,16$ veces. Contraste al nivel $0{,}001$ si los datos siguen la ley de Benford.
\fuente{adaptado de Orloff y Bloom, MIT 18.05~\cite{mit1805}, \emph{problem set}~9, problema~4.}
\end{ejercicio}

\begin{ejercicio}{2}{eje:t5-dos-modelos}
Dos modelos se evalúan sobre muestras independientes de $1000$ preguntas cada una: A acierta $780$ y B, $750$. (a) Plantee el contraste $H_0:p_A=p_B$ frente a $H_1:p_A\neq p_B$. (b) Bajo $H_0$, estime la varianza de $\widehat p_A-\widehat p_B$ con la proporción combinada y calcule el estadístico $z$ y el $p$-valor. (c) ¿Es significativa la diferencia al $5\,\%$?
\fuente{adaptado de Stats~200, Stanford~\cite{stats200}, \emph{homework}~3, problema~3.}
\end{ejercicio}

\begin{ejercicio}{2}{eje:t5-multiples}
Se compara un modelo con otro en $20$ benchmarks independientes y, en realidad, ambos son igual de buenos en todos. (a) Si cada comparación se contrasta al $5\,\%$, ¿cuál es la probabilidad de encontrar al menos una diferencia significativa? ¿Cuántas se esperan? (b) ¿A qué nivel debe hacerse cada contraste con la corrección de Bonferroni? ¿Cuál es entonces exactamente la probabilidad de al menos un falso positivo?
\end{ejercicio}

\begin{ejercicio}{3}{eje:t5-sidak}
Se realizan $n$ contrastes independientes con $p$-valores $P_1,\dots,P_n$. (a) Si todas las nulas son ciertas y los $P_i$ son \iid\ $\Unif(0,1)$, demuestre que $\PP(\min_iP_i\le t)=1-(1-t)^{n}$ para $t\in(0,1)$. (b) Deduzca que realizar cada contraste al nivel $1-(1-\alpha)^{1/n}$ da una probabilidad de algún falso positivo exactamente igual a $\alpha$. (c) Demuestre que $\alpha/n\le1-(1-\alpha)^{1/n}$, de modo que este procedimiento es al menos tan potente como el de Bonferroni. (\emph{Indicación:} desigualdad de Bernoulli, $(1+x)^{n}\ge1+nx$ para $x\ge-1$.)
\fuente{adaptado de Stats~200, Stanford~\cite{stats200}, \emph{homework}~4, problema~4.}
\end{ejercicio}

\begin{ejercicio}{3}{eje:t5-uniforme-p}
Sea $T$ un estadístico con función de distribución continua y estrictamente creciente $F_0$ bajo $H_0$, y sea $p=1-F_0(T)$ el $p$-valor del contraste que rechaza para valores grandes de $T$. Demuestre que, bajo $H_0$, $p\sim\Unif(0,1)$, y deduzca que $\PP_{H_0}(p\le\alpha)=\alpha$.
\end{ejercicio}

\begin{ejercicio}{3}{eje:t5-dualidad}
Sean datos $\Nor(\mu,\sigma^{2})$ con $\sigma$ conocida. Demuestre que, para cada $\mu_0$, el contraste $z$ bilateral de nivel $\alpha$ \emph{no} rechaza $H_0:\mu=\mu_0$ si y solo si $\mu_0$ pertenece al intervalo de confianza de nivel $1-\alpha$ del \cref{teo:t4-ic-z}. Interprete el resultado.
\fuente{adaptado de Orloff y Bloom, MIT 18.05~\cite{mit1805}, problemas de la clase~23, problema~2.}
\end{ejercicio}

\begin{ejercicio}{3}{eje:t5-potencia}
Con la función de potencia de la \cref{prop:t5-potencia}: (a) demuestre que $\pi$ es creciente en $\mu$ y que $\pi(\mu_0)=\alpha$; (b) demuestre que, para $\mu>\mu_0$ fijo, $\pi(\mu)\to1$ cuando $n\to\infty$; (c) demuestre que la potencia contra $\mu_1=\mu_0+\delta$ ($\delta>0$) es al menos $1-\beta$ si y solo si
\[
n\ge\Bigl(\frac{(z_{1-\alpha}+z_{1-\beta})\,\sigma}{\delta}\Bigr)^{2};
\]
(d) aplíquelo al \cref{ej:t5-z} con $\delta=20$, $\alpha=0{,}05$ y potencia $0{,}9$.
\end{ejercicio}

% =====================================================================
\chapter{Regresión lineal simple y correlación}
\label{tema:6}
% =====================================================================

La regresión estudia cómo cambia en promedio una variable respuesta $y$ cuando cambia una variable explicativa $x$. Este tema deduce la recta de mínimos cuadrados, estudia los residuos y la descomposición de la variabilidad que da lugar al coeficiente de determinación $R^{2}$, muestra cómo las transformaciones logarítmicas permiten ajustar leyes de potencia y termina con la distinción entre correlación y causalidad. Se sigue el capítulo~12 de OpenStax~\cite{openstax} (§12.3 «The Regression Equation»), el capítulo~14 de Rice~\cite{rice}, «Linear Least Squares», el capítulo~13 de Wasserman~\cite{wasserman}, «Linear and Logistic Regression», y la lectura~26 del curso 18.05 del MIT~\cite{mit1805}; para la causalidad, el capítulo~16 de Wasserman, «Causal Inference».

\section{La recta de mínimos cuadrados}

Se dispone de $n\ge2$ pares $(x_1,y_1),\dots,(x_n,y_n)$ con los $x_i$ no todos iguales. Se usan las sumas centradas
\[
S_{xx}=\sum_{i=1}^{n}(x_i-\xbar)^{2},\qquad S_{yy}=\sum_{i=1}^{n}(y_i-\ybar)^{2},\qquad S_{xy}=\sum_{i=1}^{n}(x_i-\xbar)(y_i-\ybar),
\]
de modo que $s_x^{2}=S_{xx}/(n-1)$, $s_{xy}=S_{xy}/(n-1)$ y $r=S_{xy}/\sqrt{S_{xx}S_{yy}}$ (\cref{def:t1-cov}). Por la \cref{prop:t1-abreviada}, $S_{xx}=\sum x_i^{2}-n\xbar^{2}$ y, análogamente, $S_{xy}=\sum x_iy_i-n\xbar\ybar$.

\begin{definicion}[Recta de mínimos cuadrados]\label{def:t6-mc}
Para cada recta $y=b_0+b_1x$, el \emph{valor ajustado} del dato $i$ es $\widehat y_i=b_0+b_1x_i$ y el \emph{residuo} es $e_i=y_i-\widehat y_i$. La \emph{recta de mínimos cuadrados} (o de regresión de $y$ sobre $x$) es la que minimiza la suma de cuadrados de los residuos
\[
\operatorname{SC}(b_0,b_1)=\sum_{i=1}^{n}(y_i-b_0-b_1x_i)^{2}.
\]
\end{definicion}

\begin{teorema}[Coeficientes de mínimos cuadrados]\label{teo:t6-mc}
Si $S_{xx}>0$, la función $\operatorname{SC}$ tiene un único mínimo, en
\[
\widehat b_1=\frac{S_{xy}}{S_{xx}}=r\,\frac{s_y}{s_x},\qquad\widehat b_0=\ybar-\widehat b_1\xbar,
\]
y el valor mínimo es $S_{yy}-S_{xy}^{2}/S_{xx}=S_{yy}(1-r^{2})$.
\end{teorema}

\begin{proof}
\emph{Paso 1 (separar la ordenada).} Sea $c=\ybar-b_0-b_1\xbar$. Entonces $y_i-b_0-b_1x_i=\bigl[(y_i-\ybar)-b_1(x_i-\xbar)\bigr]+c$. Al desarrollar el cuadrado y sumar, el término cruzado es $2c\sum_i\bigl[(y_i-\ybar)-b_1(x_i-\xbar)\bigr]=0$, porque las desviaciones respecto de la media suman cero (\cref{prop:t1-extremal}). Por tanto,
\[
\operatorname{SC}(b_0,b_1)=\sum_{i}\bigl[(y_i-\ybar)-b_1(x_i-\xbar)\bigr]^{2}+n\,c^{2}.
\]
\emph{Paso 2 (completar el cuadrado en $b_1$).} Desarrollando el primer sumando,
\[
\sum_i\bigl[(y_i-\ybar)-b_1(x_i-\xbar)\bigr]^{2}=S_{yy}-2b_1S_{xy}+b_1^{2}S_{xx}=S_{xx}\Bigl(b_1-\frac{S_{xy}}{S_{xx}}\Bigr)^{2}+S_{yy}-\frac{S_{xy}^{2}}{S_{xx}}.
\]
\emph{Paso 3 (minimizar).} Como $S_{xx}>0$, $\operatorname{SC}(b_0,b_1)\ge S_{yy}-S_{xy}^{2}/S_{xx}$, con igualdad si y solo si $b_1=S_{xy}/S_{xx}$ y $c=0$, es decir, $b_0=\ybar-b_1\xbar$. Finalmente, $S_{xy}/S_{xx}=r\sqrt{S_{yy}/S_{xx}}=r\,s_y/s_x$ y $S_{xy}^{2}/S_{xx}=r^{2}S_{yy}$.
\end{proof}

La condición $\widehat b_0=\ybar-\widehat b_1\xbar$ dice que la recta pasa por el punto de medias $(\xbar,\ybar)$. La pendiente tiene el signo de $r$: es positiva cuando la asociación es creciente.

\begin{ejemplo}\label{ej:t6-recta}
Para los datos del \cref{ej:t1-r}, $(1,2),(2,4),(3,5),(4,4),(5,5)$: $S_{xx}=10$, $S_{yy}=6$ y $S_{xy}=6$. Entonces $\widehat b_1=0{,}6$, $\widehat b_0=4-0{,}6\cdot3=2{,}2$ y la recta es $\widehat y=2{,}2+0{,}6x$ (\cref{fig:t6-recta}). Los residuos son $-0{,}8;\allowbreak\ 0{,}6;\allowbreak\ 1{,}0;\allowbreak\ -0{,}6;\allowbreak\ -0{,}2$, que suman cero, y su suma de cuadrados es $2{,}4=6\,(1-0{,}6)$.
\end{ejemplo}

\begin{figure}[ht]
\centering
\begin{tikzpicture}[x=1.1cm,y=0.75cm]
  \draw[->] (0,0) -- (6,0) node[right]{\small $x$};
  \draw[->] (0,0) -- (0,6.2) node[above]{\small $y$};
  \foreach \i in {1,...,5}{\draw (\i,0) -- (\i,-0.12) node[below]{\scriptsize\i};}
  \foreach \j in {1,...,6}{\draw (0,\j) -- (-0.1,\j) node[left]{\scriptsize\j};}
  \draw[thick] (0.5,2.5) -- (5.6,5.56) node[right]{\small $\widehat y=2{,}2+0{,}6x$};
  \foreach \x/\y/\yh in {1/2/2.8,2/4/3.4,3/5/4.0,4/4/4.6,5/5/5.2}{
    \draw[dashed,gray] (\x,\y) -- (\x,\yh);
    \fill (\x,\y) circle (2pt);
  }
  \draw (3,4) circle (3pt);
  \node[below right] at (3,4) {\scriptsize $(\xbar,\ybar)$};
\end{tikzpicture}
\caption{Recta de mínimos cuadrados del \cref{ej:t6-recta}. Los segmentos discontinuos son los residuos; la recta pasa por el punto de medias $(3,4)$.}\label{fig:t6-recta}
\end{figure}

\section{Residuos y coeficiente de determinación}

\begin{proposicion}[Propiedades de los residuos]\label{prop:t6-residuos}
Para la recta de mínimos cuadrados, $\sum_ie_i=0$, $\sum_ix_ie_i=0$ y $\sum_i\widehat y_ie_i=0$.
\end{proposicion}

La demostración es el \cref{eje:t6-residuos}. Las dos primeras igualdades son las \emph{ecuaciones normales} (\cref{eje:t6-normales}): los residuos son ortogonales, como vectores de $\R^{n}$, al vector de unos y al vector de las $x_i$. Este es el punto de vista geométrico que el módulo~8 generaliza a varias variables explicativas.

\begin{teorema}[Descomposición de la variabilidad]\label{teo:t6-anova}
Sean $\operatorname{SC}_T=\sum_i(y_i-\ybar)^{2}=S_{yy}$ (suma de cuadrados total), $\operatorname{SC}_R=\sum_i(\widehat y_i-\ybar)^{2}$ (explicada por la regresión) y $\operatorname{SC}_E=\sum_ie_i^{2}$ (residual o de error). Entonces
\[
\operatorname{SC}_T=\operatorname{SC}_R+\operatorname{SC}_E.
\]
\end{teorema}

\begin{proof}
Se escribe $y_i-\ybar=(\widehat y_i-\ybar)+e_i$ y se eleva al cuadrado:
$\sum_i(y_i-\ybar)^{2}=\sum_i(\widehat y_i-\ybar)^{2}+\sum_ie_i^{2}+2\sum_i(\widehat y_i-\ybar)e_i$.
El último término es $2\sum_i\widehat y_ie_i-2\ybar\sum_ie_i=0-0$ por la \cref{prop:t6-residuos}.
\end{proof}

\begin{definicion}[Coeficiente de determinación]\label{def:t6-r2}
El \emph{coeficiente de determinación} es $R^{2}=\operatorname{SC}_R/\operatorname{SC}_T=1-\operatorname{SC}_E/\operatorname{SC}_T$, la fracción de la variabilidad de $y$ explicada por la recta. Por el \cref{teo:t6-anova}, $0\le R^{2}\le1$.
\end{definicion}

\begin{corolario}\label{cor:t6-r2}
En la regresión lineal simple, $R^{2}=r^{2}$.
\end{corolario}

\begin{proof}
Por el \cref{teo:t6-mc}, $\operatorname{SC}_E=S_{yy}(1-r^{2})$, luego $R^{2}=1-\operatorname{SC}_E/S_{yy}=r^{2}$.
\end{proof}

En el \cref{ej:t6-recta}, $R^{2}=0{,}6$: la recta explica el $60\,\%$ de la variabilidad de $y$. La dispersión de los puntos alrededor de la recta se resume en el \emph{error típico residual} $s_e=\sqrt{\operatorname{SC}_E/(n-2)}$; se divide por $n-2$ porque se han estimado dos parámetros (compárese con la \cref{obs:t1-n-1}).

\section{Diagnóstico con residuos}

Un $R^{2}$ alto no garantiza que el modelo lineal sea adecuado. La herramienta básica de diagnóstico es el \emph{gráfico de residuos}, que representa $e_i$ frente a $x_i$ (o frente a $\widehat y_i$). Si el modelo es adecuado, los residuos se distribuyen sin estructura alrededor de cero. Algunas señales de alarma~\cite[§12.3]{openstax}:
\begin{enumerate}
  \item una curvatura sistemática (residuos positivos en los extremos y negativos en el centro, o al revés) indica que la relación no es lineal;
  \item una dispersión que crece con $x$ (forma de embudo) indica varianza no constante;
  \item un punto con un residuo muy grande es un posible atípico, y un punto con $x_i$ muy alejado de $\xbar$ puede ser muy \emph{influyente}: puede cambiar mucho la recta por sí solo.
\end{enumerate}
Además, la recta solo describe la relación en el rango de los $x_i$ observados: \emph{extrapolar} fuera de ese rango (\cref{eje:t6-buceo}) es arriesgado.

\begin{observacion}[El modelo estadístico]\label{obs:t6-modelo}
Si se supone que $y_i=\beta_0+\beta_1x_i+\varepsilon_i$ con errores $\varepsilon_i\sim\Nor(0,\sigma^{2})$ \iid, la log-verosimilitud de $(\beta_0,\beta_1)$ es una constante menos $\frac{1}{2\sigma^{2}}\operatorname{SC}(\beta_0,\beta_1)$; por tanto, los estimadores de mínimos cuadrados son los estimadores de máxima verosimilitud~\cite[cap.~13]{wasserman}, \cite[\emph{problem set}~6]{mit1805}. Con este modelo pueden construirse intervalos de confianza y contrastes para la pendiente (ampliación), que se estudian en~\cite[cap.~14]{rice}.
\end{observacion}

\section{Transformaciones: ajuste de leyes de potencia}

Muchas relaciones no lineales se vuelven lineales tras una transformación. Una \emph{ley de potencia} $y=a\,x^{b}$ con $a,x,y>0$ cumple
\[
\log y=\log a+b\log x,
\]
de modo que en escala log-log es una recta de pendiente $b$ y ordenada $\log a$. Para ajustarla se aplica mínimos cuadrados a los pares $(\log x_i,\log y_i)$ y se deshace la transformación: $\widehat a=10^{\widehat b_0}$ (o $e^{\widehat b_0}$ si se usan logaritmos naturales) y $\widehat b=\widehat b_1$. Análogamente, una relación exponencial $y=a\,e^{bx}$ se linealiza con $\ln y=\ln a+bx$. Obsérvese que así se minimizan los errores \emph{relativos} (en escala logarítmica), no los absolutos.

\begin{ejemplo}\label{ej:t6-potencia}
Si la pérdida de un modelo cumpliera exactamente $L=20\,N^{-0{,}1}$, al multiplicar el número de parámetros $N$ por $10$ la pérdida se multiplicaría por $10^{-0{,}1}\approx0{,}794$, y en escala log-log los puntos estarían sobre una recta de pendiente $-0{,}1$. El \cref{eje:t6-loglog} ajusta una ley así a partir de cuatro mediciones.
\end{ejemplo}

\section{Correlación y causalidad}

Una correlación entre $x$ e $y$ puede deberse a que $x$ influye en $y$, a que $y$ influye en $x$ (causalidad inversa), a que una tercera variable $z$ influye en ambas (\emph{variable de confusión}) o simplemente al azar. Los datos observacionales, por sí solos, no permiten distinguir estos casos~\cite[cap.~16]{wasserman}.

\begin{ejemplo}[Variable de confusión]\label{ej:t6-confusion}
En un conjunto de modelos de lenguaje, la longitud media de las respuestas y la puntuación en un benchmark de razonamiento pueden estar fuertemente correlacionadas. De ahí no se sigue que alargar artificialmente las respuestas mejore la puntuación: ambos rasgos pueden deberse a una tercera variable, como el tamaño del modelo o el entrenamiento con cadenas de razonamiento.
\end{ejemplo}

Un caso extremo de confusión es la \emph{paradoja de Simpson}: una asociación que se da en cada uno de los grupos de una población puede invertirse al agregar los grupos, cuando los tamaños de los grupos difieren entre las condiciones comparadas (\cref{eje:t6-simpson}).

La forma más fiable de establecer una relación causal es el \emph{experimento aleatorizado}: si el tratamiento (por ejemplo, el \emph{prompt} A o el B) se asigna al azar a las unidades (las preguntas), las variables de confusión se reparten en promedio por igual entre los grupos, y una diferencia sistemática en la respuesta solo puede atribuirse al tratamiento o al azar, cuya magnitud se cuantifica con los contrastes del \cref{tema:5}.

\begin{observacion}[Errores frecuentes]\label{obs:t6-errores}
(a) Interpretar la pendiente como efecto causal con datos observacionales.
(b) Extrapolar la recta fuera del rango de los datos.
(c) Concluir que no hay relación porque $r\approx0$ (la relación puede ser no lineal; \cref{eje:t6-cuadratica}).
(d) Creer que la recta de regresión de $x$ sobre $y$ es la inversa de la de $y$ sobre $x$: solo lo es si $\abs r=1$ (\cref{eje:t6-resumen}).
(e) Juzgar un ajuste solo por $R^{2}$, sin mirar el gráfico de residuos.
\end{observacion}

\begin{observacion}[Conexión con el curso]
El módulo~8 generaliza la recta de mínimos cuadrados a varias variables mediante las ecuaciones normales $X^{\top}X\widehat{\bm\beta}=X^{\top}\mathbf y$ y su interpretación como proyección, de la que la \cref{prop:t6-residuos} es el caso más sencillo, y demuestra que mínimos cuadrados es máxima verosimilitud con errores normales (\cref{obs:t6-modelo}). El módulo~16 ajusta en escala log-log las leyes de escala que relacionan la pérdida con el número de parámetros, de datos y de operaciones (\cref{ej:t6-potencia}). En el módulo~23, las sondas lineales detectan \emph{correlaciones} entre activaciones y conceptos, mientras que las intervenciones causales (modificar una activación y observar el efecto) son el análogo del experimento aleatorizado.
\end{observacion}

\begin{lecturas}
\item OpenStax~\cite{openstax}, §12.3 «The Regression Equation» y §12.4 «Testing the Significance of the Correlation Coefficient».
\item Rice~\cite{rice}, cap.~14 «Linear Least Squares».
\item Wasserman~\cite{wasserman}, cap.~13 «Linear and Logistic Regression» y cap.~16 «Causal Inference».
\item Orloff y Bloom~\cite{mit1805}, lectura~26 «Linear regression».
\item Stats~200~\cite{stats200}, clase~25 «The linear model».
\end{lecturas}

\section*{Ejercicios del tema 6}
\addcontentsline{toc}{section}{Ejercicios del tema 6}

\begin{ejercicio}{1}{eje:t6-vf}
Indique si cada afirmación es verdadera o falsa. (a) La recta de mínimos cuadrados pasa siempre por $(\xbar,\ybar)$. (b) Si $r=-0{,}9$, entonces $R^{2}=-0{,}81$. (c) Un $R^{2}$ cercano a~$1$ prueba que $x$ causa $y$. (d) La suma de los residuos de mínimos cuadrados es cero. (e) Si se intercambian los papeles de $x$ e $y$, la pendiente pasa a ser $1/\widehat b_1$.
\end{ejercicio}

\begin{ejercicio}{2}{eje:t6-basico}
Para los datos $(0,1),(1,3),(2,2),(3,5),(4,4)$: (a) calcule $S_{xx}$, $S_{yy}$, $S_{xy}$ y $r$; (b) obtenga la recta de mínimos cuadrados; (c) calcule los residuos, $\operatorname{SC}_E$, $R^{2}$ y $s_e$; (d) prediga $y$ para $x=2{,}5$.
\end{ejercicio}

\begin{ejercicio}{2}{eje:t6-gauss}
Se ajusta el modelo $y_i=ax_i+b+\varepsilon_i$ con errores $\Nor(0,\sigma^{2})$ \iid\ a los datos $(1,8),(3,2),(5,1)$. (a) Escriba la log-verosimilitud de un dato $(x_1,y_1)$. (b) Obtenga los EMV de $a$ y $b$ y compruebe que no dependen de $\sigma$.
\fuente{adaptado de Orloff y Bloom, MIT 18.05~\cite{mit1805}, \emph{problem set}~6, problema~2.}
\end{ejercicio}

\begin{ejercicio}{2}{eje:t6-examenes}
Las notas de once estudiantes en un tercer examen ($x$, sobre $80$) y en el final ($y$, sobre $200$) cumplen $n=11$, $\sum x_i=761$, $\sum y_i=1765$, $\sum x_i^{2}=52\,729$, $\sum y_i^{2}=287\,529$ y $\sum x_iy_i=122\,500$. (a) Calcule la recta de regresión y $r$. (b) Prediga la nota final de un estudiante con $73$ en el tercer examen. (c) Un estudiante tiene $(65,175)$; calcule su residuo.
\fuente{adaptado de OpenStax, \emph{Introductory Statistics}~2e~\cite{openstax}, §12.3, ejemplo~12.6.}
\end{ejercicio}

\begin{ejercicio}{2}{eje:t6-buceo}
La profundidad de una inmersión ($x$, en pies) y el tiempo máximo de inmersión ($y$, en minutos) son $(50,80),(60,55),(70,45),(80,35),(90,25),(100,22)$. (a) Calcule la recta de regresión y $r$. (b) Prediga $y$ para $x=85$ y para $x=110$. ¿Cuál de las dos predicciones es fiable? (c) Observe el signo de los residuos en los extremos y en el centro: ¿qué sugiere?
\fuente{adaptado de OpenStax, \emph{Introductory Statistics}~2e~\cite{openstax}, §12.3, \emph{Try It}~12.6.}
\end{ejercicio}

\begin{ejercicio}{2}{eje:t6-helados}
La temperatura máxima ($x$, en grados Fahrenheit) y el número de helados vendidos ($y$) en nueve días son $(70,105)$, $(85,240)$, $(65,49)$, $(72,147)$, $(80,231)$, $(61,38)$, $(75,193)$, $(78,196)$, $(68,89)$. Se dan $\sum x_i=654$, $\sum y_i=1288$, $\sum x_i^{2}=47\,988$, $\sum y_i^{2}=231\,026$ y $\sum x_iy_i=98\,132$. Calcule la recta de regresión, $r$ y $R^{2}$, y prediga las ventas de un día a $74$ grados.
\fuente{adaptado de OpenStax, \emph{Introductory Statistics}~2e~\cite{openstax}, §12.3, ejemplo «Ice Cream Cones».}
\end{ejercicio}

\begin{ejercicio}{2}{eje:t6-resumen}
De un conjunto de pares se conocen $\xbar=10$, $\ybar=50$, $s_x=2$, $s_y=8$ y $r=0{,}75$. (a) Obtenga la recta de regresión de $y$ sobre $x$ y prediga $y$ para $x=12$. (b) Obtenga la recta de regresión de $x$ sobre $y$. (c) Compruebe que el producto de las dos pendientes es $r^{2}$ y que las dos rectas no coinciden.
\end{ejercicio}

\begin{ejercicio}{2}{eje:t6-loglog}
Se mide la pérdida $L$ de cuatro modelos con $N=10^{6},10^{7},10^{8},10^{9}$ parámetros y se obtiene $L=5{,}0;\allowbreak\ 4{,}1;\allowbreak\ 3{,}3;\allowbreak\ 2{,}7$, cuyos logaritmos decimales son $0{,}6990;\allowbreak\ 0{,}6128;\allowbreak\ 0{,}5185;\allowbreak\ 0{,}4314$. (a) Ajuste una ley de potencia $L=aN^{-\alpha}$ por mínimos cuadrados en escala log-log. (b) Prediga la pérdida para $N=10^{10}$. (c) ¿Por qué esta predicción es una extrapolación?
\end{ejercicio}

\begin{ejercicio}{2}{eje:t6-simpson}
Dos modelos se evalúan en un banco con preguntas fáciles y difíciles. El modelo A acierta $45$ de $50$ fáciles y $75$ de $250$ difíciles; el modelo B acierta $200$ de $250$ fáciles y $10$ de $50$ difíciles. (a) Calcule las exactitudes de cada modelo en cada tipo de pregunta y en total. (b) ¿Qué modelo es mejor? Explique la aparente contradicción.
\end{ejercicio}

\begin{ejercicio}{2}{eje:t6-cuadratica}
Para los puntos $(-2,4),(-1,1),(0,0),(1,1),(2,4)$, en los que $y=x^{2}$: (a) calcule $r$ y la recta de mínimos cuadrados; (b) calcule los residuos y describa el patrón del gráfico de residuos; (c) ¿qué se concluye?
\end{ejercicio}

\begin{ejercicio}{3}{eje:t6-normales}
Demuestre el \cref{teo:t6-mc} por otra vía: (a) obtenga las ecuaciones normales igualando a cero las derivadas parciales de $\operatorname{SC}(b_0,b_1)$; (b) resuélvalas; (c) compruebe que la matriz hessiana es definida positiva, de modo que el punto crítico es un mínimo (y, por ser $\operatorname{SC}$ un polinomio de grado~2, global).
\end{ejercicio}

\begin{ejercicio}{3}{eje:t6-residuos}
Demuestre la \cref{prop:t6-residuos}: $\sum_ie_i=0$, $\sum_ix_ie_i=0$ y $\sum_i\widehat y_ie_i=0$. Deduzca que la media de los valores ajustados es $\ybar$.
\end{ejercicio}

\begin{ejercicio}{3}{eje:t6-origen}
(a) Demuestre que la recta que pasa por el origen y minimiza $\sum_i(y_i-bx_i)^{2}$ tiene pendiente $\widehat b=\sum_ix_iy_i\big/\sum_ix_i^{2}$ (suponga que algún $x_i\ne0$). (b) Demuestre que sus residuos cumplen $\sum_ix_ie_i=0$ y dé un ejemplo en el que $\sum_ie_i\neq0$.
\end{ejercicio}

\begin{ejercicio}{3}{eje:t6-transformaciones}
Sean $x_i'=\alpha x_i+\beta$ e $y_i'=\gamma y_i+\delta$ con $\alpha,\gamma\neq0$. Demuestre que la pendiente de mínimos cuadrados de $y'$ sobre $x'$ es $\widehat b_1'=(\gamma/\alpha)\,\widehat b_1$, que $R^{2}$ no cambia y que los residuos se multiplican por $\gamma$. ¿Cómo cambia la pendiente al pasar $x$ de metros a centímetros?
\end{ejercicio}

\begin{ejercicio}{3}{eje:t6-regresion-media}
Sean $z_{x,i}=(x_i-\xbar)/s_x$ y $z_{y,i}=(y_i-\ybar)/s_y$. (a) Demuestre que la recta de mínimos cuadrados, escrita en unidades tipificadas, es $\widehat z_y=r\,z_x$. (b) Deduzca que, si $\abs r<1$, la predicción de $y$ está a menos desviaciones típicas de su media que $x$ de la suya («regresión a la media»). (c) Un modelo obtiene en una primera evaluación una puntuación $2$ desviaciones típicas por encima de la media de los modelos; si la correlación entre dos evaluaciones es $0{,}5$, ¿qué puntuación tipificada se predice para la segunda?
\end{ejercicio}


\appendix
% =====================================================================
\chapter{Soluciones de los ejercicios}
\label{ap:soluciones}
% =====================================================================

Cada solución se presenta en cinco partes: datos, planteamiento, desarrollo, resultado y comprobación. En los pasos intermedios se redondea a cuatro cifras significativas salvo que se indique otra cosa; los valores de $\Phi$ y los cuantiles se toman de las tablas del \cref{ap:tablas} o se dan con cuatro decimales. En los ejercicios de demostración, el apartado «Desarrollo» contiene la demostración redactada tal como debería entregarse.

% ---------------------------------------------------------------------
\section{Tema 1}

\begin{solucion}{eje:t1-vf}
\planteamiento contrastar cada afirmación con las definiciones del \cref{tema:1}.
\pasos
\begin{enumerate}[label=(\alph*)]
  \item \textbf{Verdadera} (si $n\ge 3$). El mayor dato sigue siendo el mayor después de multiplicarlo por $10$, de modo que el orden de los demás no cambia y la mediana, que depende solo de los valores centrales, tampoco. (Con $n=2$ la mediana es la media de los dos datos y sí cambia.)
  \item \textbf{Falsa.} $s$ es la raíz cuadrada no negativa de una suma de cuadrados dividida por $n-1>0$.
  \item \textbf{Verdadera.} Si $x_i=c$ para todo $i$, entonces $\xbar=c$, todas las desviaciones son nulas y $s=0$; además $Q_1=Q_3=c$.
  \item \textbf{Falsa.} Por el \cref{teo:t1-r}, los puntos están sobre una recta si y solo si $\abs r=1$.
  \item \textbf{Verdadera.} Por la \cref{prop:t1-lineal} con $a=1$, $b=5$: $s_y^{2}=s_x^{2}$ y $\RIC_y=\RIC_x$.
\end{enumerate}
\resultado \textbf{(a) V (para $n\ge3$); (b) F; (c) V; (d) F; (e) V.}
\end{solucion}

\begin{solucion}{eje:t1-basico}
\datos $4,7,2,9,5,6,3,25,6$; $n=9$.
\planteamiento ordenar, aplicar las \cref{def:t1-centro,def:t1-cuartiles,def:t1-caja} y la fórmula de cálculo (\cref{prop:t1-abreviada}).
\pasos
\begin{enumerate}[label=\arabic*.]
  \item Ordenados: $2,3,4,5,6,6,7,9,25$. Suma $67$; $\xbar=67/9\approx 7{,}444$.
  \item Mediana: el dato quinto, $6$. Moda: $6$ (aparece dos veces).
  \item Mitad inferior $2,3,4,5$: $Q_1=(3+4)/2=3{,}5$. Mitad superior $6,7,9,25$: $Q_3=(7+9)/2=8$. $\RIC=4{,}5$.
  \item Vallas: $3{,}5-1{,}5\cdot4{,}5=-3{,}25$ y $8+6{,}75=14{,}75$. El dato $25$ es atípico.
  \item $\sum x_i^{2}=16+49+4+81+25+36+9+625+36=881$; $n\xbar^{2}=67^{2}/9=4489/9\approx 498{,}78$; suma de cuadrados $881-498{,}78=382{,}22$; $s^{2}=382{,}22/8\approx 47{,}78$; $s\approx 6{,}912$.
  \item Sin el $25$ ($n=8$): suma $42$, $\xbar=5{,}25$; mediana $(5+6)/2=5{,}5$; $\sum x_i^{2}=256$, suma de cuadrados $256-8\cdot5{,}25^{2}=256-220{,}5=35{,}5$, $s^{2}=35{,}5/7\approx5{,}071$, $s\approx 2{,}252$.
\end{enumerate}
\resultado \textbf{$\xbar\approx7{,}44$; mediana $6$; moda $6$; $Q_1=3{,}5$; $Q_3=8$; $\RIC=4{,}5$; vallas $-3{,}25$ y $14{,}75$; $s^{2}\approx47{,}78$; $s\approx6{,}91$; el $25$ es atípico.} Sin él, $\xbar=5{,}25$, mediana $5{,}5$, $s\approx2{,}25$: la media baja unas dos unidades y $s$ se reduce a menos de la tercera parte, mientras que la mediana apenas cambia.
\comprobacion la suma de desviaciones con todos los datos es $67-9\cdot\frac{67}{9}=0$, como exige la \cref{prop:t1-extremal}.
\end{solucion}

\begin{solucion}{eje:t1-edades}
\datos $n=20$ edades; $\sum x_i=210{,}5$; $\sum x_i^{2}=2225{,}25$.
\planteamiento aplicar las definiciones; los datos ya están ordenados.
\pasos
\begin{enumerate}[label=\arabic*.]
  \item $\xbar=210{,}5/20=10{,}525$.
  \item Mediana: media de los datos décimo y undécimo, ambos $10{,}5$; mediana $10{,}5$.
  \item Mitad inferior (diez primeros): sus valores quinto y sexto son $10$ y $10$, luego $Q_1=10$. Mitad superior (diez últimos): sus valores quinto y sexto son $11$ y $11$, luego $Q_3=11$.
  \item Suma de cuadrados: $2225{,}25-20\cdot10{,}525^{2}=2225{,}25-2215{,}5125=9{,}7375$; $s^{2}=9{,}7375/19\approx0{,}5125$; $s\approx0{,}7159$.
  \item Intervalo $\xbar\pm s=(9{,}809;\allowbreak\ 11{,}241)$. Contiene los valores $10$ (cuatro veces), $10{,}5$ (cuatro) y $11$ (seis): $14$ datos.
\end{enumerate}
\resultado \textbf{$\xbar=10{,}525$; mediana $10{,}5$; $Q_1=10$; $Q_3=11$; $s\approx0{,}716$; el $70\,\%$ de los datos está a menos de una desviación típica de la media.}
\comprobacion la fuente da $s\approx0{,}72$. El $70\,\%$ es parecido al $68\,\%$ que se esperaría con datos de forma aproximadamente normal (\cref{tema:3}).
\end{solucion}

\begin{solucion}{eje:t1-cajas}
\datos mañana $(32;56;74{,}5;82{,}5;99)$; tarde $(25{,}5;78;81;89;98)$.
\planteamiento rango $=$ máximo $-$ mínimo; $\RIC=Q_3-Q_1$; vallas de la \cref{def:t1-caja}.
\pasos
\begin{enumerate}[label=\arabic*.]
  \item Mañana: rango $99-32=67$; $\RIC=82{,}5-56=26{,}5$.
  \item Tarde: rango $98-25{,}5=72{,}5$; $\RIC=89-78=11$.
  \item La mitad central más dispersa es la de la mañana ($26{,}5>11$), aunque el rango de la tarde sea mayor.
  \item Vallas de la mañana: $56-1{,}5\cdot26{,}5=16{,}25$ y $82{,}5+39{,}75=122{,}25$; el mínimo $32$ no es atípico.
  \item Vallas de la tarde: $78-16{,}5=61{,}5$ y $89+16{,}5=105{,}5$; el mínimo $25{,}5<61{,}5$ es atípico, y el bigote inferior acaba en el menor dato no atípico, que no se conoce con el resumen (está entre $61{,}5$ y $78$).
  \item Esbozo: la caja de la tarde es estrecha y está desplazada a la derecha (de $78$ a $89$, mediana $81$), con el $25{,}5$ marcado aparte; la de la mañana es ancha (de $56$ a $82{,}5$, mediana $74{,}5$) con bigotes hasta $32$ y $99$.
\end{enumerate}
\resultado \textbf{Rangos $67$ y $72{,}5$; $\RIC$ $26{,}5$ y $11$; la mañana tiene la mitad central más dispersa; solo el mínimo de la tarde es atípico.}
\comprobacion el rango de la tarde es mayor solo por un dato: el $\RIC$, robusto, revela que el grupo de la tarde es más homogéneo.
\end{solucion}

\begin{solucion}{eje:t1-agrupados}
\datos marcas $m_j=1,4,7,10,13,16$; frecuencias $f_j=1,6,10,7,0,2$; $n=26$.
\planteamiento $\xbar\approx\frac1n\sum_j f_jm_j$ y $s^{2}\approx\frac{1}{n-1}\bigl(\sum_j f_jm_j^{2}-n\xbar^{2}\bigr)$.
\pasos
\begin{enumerate}[label=\arabic*.]
  \item $\sum f_jm_j=1+24+70+70+0+32=197$; $\xbar=197/26\approx7{,}577$.
  \item $\sum f_jm_j^{2}=1+96+490+700+0+512=1799$.
  \item $n\xbar^{2}=197^{2}/26=38\,809/26\approx1492{,}65$; suma de cuadrados $1799-1492{,}65=306{,}35$.
  \item $s^{2}\approx306{,}35/25\approx12{,}25$; $s\approx3{,}50$.
\end{enumerate}
\resultado \textbf{$\xbar\approx7{,}58$ y $s\approx3{,}50$.}
\comprobacion la fuente obtiene los mismos valores. Son aproximaciones: se ha supuesto que todos los datos de cada clase valen la marca de clase.
\end{solucion}

\begin{solucion}{eje:t1-abreviada}
\datos $n=10$, $\sum x_i=50$, $\sum x_i^{2}=290$; nuevo dato $x_{11}=5$.
\planteamiento fórmula de cálculo (\cref{prop:t1-abreviada}); las sumas se actualizan sumando el nuevo dato.
\pasos
\begin{enumerate}[label=\arabic*.]
  \item $\xbar=50/10=5$; suma de cuadrados $290-10\cdot25=40$; $s^{2}=40/9\approx4{,}444$.
  \item Con el nuevo dato: $\sum x_i=55$, $\sum x_i^{2}=290+25=315$, $n=11$, $\xbar=5$.
  \item Suma de cuadrados $315-11\cdot25=40$; $s^{2}=40/10=4$.
\end{enumerate}
\resultado \textbf{(a) $\xbar=5$, $s^{2}=40/9\approx4{,}44$. (b) $\xbar=5$, $s^{2}=4$.}
\comprobacion el dato añadido coincide con la media, de modo que no altera la suma de cuadrados ($40$); solo aumenta el denominador, y la varianza disminuye.
\end{solucion}

\begin{solucion}{eje:t1-celsius}
\datos $18,21,25,19,22,30,24$ ($n=7$); $F=1{,}8C+32$.
\planteamiento calcular en Celsius y transformar con la \cref{prop:t1-lineal} ($a=1{,}8>0$).
\pasos
\begin{enumerate}[label=\arabic*.]
  \item Ordenados: $18,19,21,22,24,25,30$; suma $159$; $\xbar=159/7\approx22{,}71$; mediana $22$.
  \item Mitad inferior $18,19,21$: $Q_1=19$; superior $24,25,30$: $Q_3=25$; $\RIC=6$.
  \item $\sum x_i^{2}=3711$; suma de cuadrados $3711-159^{2}/7=3711-3611{,}57=99{,}43$; $s^{2}=99{,}43/6\approx16{,}57$; $s\approx4{,}071$.
  \item Fahrenheit: $\xbar_F=1{,}8\cdot22{,}714+32\approx72{,}89$; mediana $1{,}8\cdot22+32=71{,}6$; $s_F=1{,}8\cdot4{,}071\approx7{,}327$; $\RIC_F=1{,}8\cdot6=10{,}8$.
  \item Las puntuaciones $z$ no cambian: $z_i^{F}=\dfrac{1{,}8x_i+32-(1{,}8\xbar+32)}{1{,}8\,s}=\dfrac{x_i-\xbar}{s}$.
\end{enumerate}
\resultado \textbf{Celsius: $\xbar\approx22{,}71$, mediana $22$, $s\approx4{,}07$, $\RIC=6$. Fahrenheit: $\xbar\approx72{,}89$, mediana $71{,}6$, $s\approx7{,}33$, $\RIC=10{,}8$. Las puntuaciones $z$ son las mismas.}
\comprobacion el desplazamiento $+32$ afecta a la media y a la mediana, pero no a $s$ ni al $\RIC$.
\end{solucion}

\begin{solucion}{eje:t1-zscore}
\datos Juan: $x=2{,}85$, media $3{,}0$, $s=0{,}7$. Alí: $x=77$, media $80$, $s=10$.
\planteamiento comparar puntuaciones tipificadas (\cref{def:t1-z}).
\pasos
\begin{enumerate}[label=\arabic*.]
  \item $z_{\text{Juan}}=(2{,}85-3{,}0)/0{,}7=-0{,}15/0{,}7\approx-0{,}214$.
  \item $z_{\text{Alí}}=(77-80)/10=-0{,}3$.
\end{enumerate}
\resultado \textbf{Juan}, porque su nota está a $0{,}21$ desviaciones típicas por debajo de la media de su centro, frente a $0{,}30$ en el caso de Alí.
\comprobacion la comparación directa de $2{,}85$ con $77$ carece de sentido porque las escalas son distintas; la tipificación las hace comparables.
\end{solucion}

\begin{solucion}{eje:t1-cov-tabla}
\datos $X,Y\in\{-1,1\}$ con $\PP(X=1)=\PP(Y=1)=1/2$; $c=\PP(X=1,Y=1)$.
\planteamiento completar la tabla con las marginales y usar $\Cov(X,Y)=\E[XY]-\E[X]\E[Y]$.
\pasos
\begin{enumerate}[label=\arabic*.]
  \item Como $\PP(X=1)=1/2$, $\PP(X=1,Y=-1)=1/2-c$; análogamente $\PP(X=-1,Y=1)=1/2-c$ y $\PP(X=-1,Y=-1)=1-c-2(1/2-c)=c$. Debe ser $0\le c\le1/2$.
  \item $\E[X]=\E[Y]=0$ y $\Var X=\Var Y=\E[X^{2}]=1$.
  \item $XY=1$ con probabilidad $2c$ y $XY=-1$ con probabilidad $1-2c$, luego $\E[XY]=2c-(1-2c)=4c-1$.
  \item $\Cov(X,Y)=4c-1$ y $\Cor(X,Y)=(4c-1)/\sqrt{1\cdot1}=4c-1$.
  \item Independencia: debe cumplirse $\PP(X=1,Y=1)=\frac12\cdot\frac12$, es decir, $c=1/4$; entonces las cuatro casillas valen $1/4$ y la tabla es producto de las marginales.
  \item $\Cor=1\iff c=1/2$ (entonces $Y=X$); $\Cor=-1\iff c=0$ (entonces $Y=-X$).
\end{enumerate}
\resultado \textbf{$\Cov(X,Y)=\Cor(X,Y)=4c-1$; independientes si y solo si $c=1/4$; correlación $1$ si $c=1/2$ y $-1$ si $c=0$.}
\comprobacion $c\in[0,1/2]$ da $\Cor\in[-1,1]$, de acuerdo con la acotación de la correlación.
\end{solucion}

\begin{solucion}{eje:t1-agrupar}
\datos $n_1=10$, $\xbar_1=70$, $s_1=8$; $n_2=15$, $\xbar_2=80$, $s_2=6$.
\planteamiento la suma de cuadrados respecto de la media global $\xbar$ se descompone por la \cref{prop:t1-extremal}(ii): $\sum_{\text{grupo }g}(x_i-\xbar)^{2}=(n_g-1)s_g^{2}+n_g(\xbar_g-\xbar)^{2}$.
\pasos
\begin{enumerate}[label=\arabic*.]
  \item $\xbar=(10\cdot70+15\cdot80)/25=(700+1200)/25=76$.
  \item Grupo 1: $9\cdot64+10\cdot(70-76)^{2}=576+360=936$.
  \item Grupo 2: $14\cdot36+15\cdot(80-76)^{2}=504+240=744$.
  \item Total: $936+744=1680$; $s^{2}=1680/24=70$; $s=\sqrt{70}\approx8{,}367$.
\end{enumerate}
\resultado \textbf{$\xbar=76$ y $s\approx8{,}37$.}
\comprobacion $s$ es mayor que las dos desviaciones de grupo porque se añade la variabilidad entre las medias de los grupos.
\end{solucion}

\begin{solucion}{eje:t1-mediana-min}
\planteamiento estudiar la pendiente de $g$ en cada intervalo entre datos ordenados consecutivos.
\pasos
\begin{enumerate}[label=\arabic*.]
  \item $g$ es continua y lineal a trozos. Sea $a$ un punto que no coincide con ningún dato y sean $L(a)=\#\{i:x_i<a\}$ y $U(a)=\#\{i:x_i>a\}$, de modo que $L(a)+U(a)=n$. Cerca de $a$, cada término con $x_i<a$ vale $a-x_i$ (pendiente $+1$) y cada término con $x_i>a$ vale $x_i-a$ (pendiente $-1$). Por tanto, $g'(a)=L(a)-U(a)$.
  \item Si $L(a)<n/2$, entonces $U(a)>n/2$ y $g'(a)<0$: $g$ decrece. Si $L(a)>n/2$, $g$ crece.
  \item Caso $n=2m+1$ impar. Para $a<x_{(m+1)}$ (y fuera de los datos) hay a lo sumo $m<n/2$ datos menores que $a$, y $g$ decrece; para $a>x_{(m+1)}$ hay al menos $m+1>n/2$, y $g$ crece. Por continuidad, el mínimo está en $a=x_{(m+1)}$, la mediana.
  \item Caso $n=2m$ par. Para $a<x_{(m)}$, $L(a)\le m-1<m$ y $g$ decrece; para $a>x_{(m+1)}$, $L(a)\ge m+1$ y $g$ crece; para $x_{(m)}<a<x_{(m+1)}$, $L(a)=U(a)=m$ y $g$ es constante. Luego $g$ es mínima en todo el intervalo $[x_{(m)},x_{(m+1)}]$, que contiene la mediana $\tfrac12(x_{(m)}+x_{(m+1)})$.
\end{enumerate}
\resultado \textbf{$g$ alcanza su mínimo en la mediana} (si $n$ es par, en todo el intervalo entre los dos datos centrales). $\blacksquare$
\comprobacion con los datos $1,2,10$: $g(2)=1+0+8=9$, $g(3)=2+1+7=10$ y $g(\xbar)=g(13/3)\approx3{,}33+2{,}33+5{,}67=11{,}33$; la media no minimiza la suma de valores absolutos.
\end{solucion}

\begin{solucion}{eje:t1-r-invariante}
\planteamiento por la \cref{prop:t1-lineal}, $x_i'-\xbar'=a(x_i-\xbar)$ e $y_i'-\ybar'=c(y_i-\ybar)$.
\pasos
\begin{enumerate}[label=\arabic*.]
  \item $s_{x'y'}=\frac{1}{n-1}\sum_i a(x_i-\xbar)\,c(y_i-\ybar)=ac\,s_{xy}$.
  \item $s_{x'}=\abs a\,s_x$ y $s_{y'}=\abs c\,s_y$.
  \item $r_{x'y'}=\dfrac{ac\,s_{xy}}{\abs a\abs c\,s_xs_y}=\dfrac{ac}{\abs{ac}}\,r_{xy}=\operatorname{signo}(ac)\,r_{xy}$.
  \item Un cambio de unidades es una transformación con $a>0$ (o $c>0$), luego no altera $r$.
\end{enumerate}
\resultado \textbf{$s_{x'y'}=ac\,s_{xy}$ y $r_{x'y'}=\operatorname{signo}(ac)\,r_{xy}$}; la correlación es adimensional. $\blacksquare$
\comprobacion con los datos del \cref{ej:t1-r} y $x'=10x$: la covarianza pasa de $1{,}5$ a $15$ y $r$ sigue valiendo $0{,}775$.
\end{solucion}

\begin{solucion}{eje:t1-chebyshev}
\planteamiento acotar inferiormente la suma de cuadrados por la contribución de los datos alejados.
\pasos
\begin{enumerate}[label=\arabic*.]
  \item Sea $A=\{i:\abs{x_i-\xbar}\ge ks\}$ y $N_A=\abs A$.
  \item Como todos los sumandos son no negativos,
  \[
  (n-1)s^{2}=\sum_{i=1}^{n}(x_i-\xbar)^{2}\ge\sum_{i\in A}(x_i-\xbar)^{2}\ge N_A\,k^{2}s^{2}.
  \]
  \item Dividiendo por $nk^{2}s^{2}>0$: $N_A/n\le(n-1)/(nk^{2})<1/k^{2}$.
  \item Con $k=2$: menos del $25\,\%$ de los datos dista $2s$ o más de la media; al menos el $75\,\%$ está en $(\xbar-2s,\xbar+2s)$.
\end{enumerate}
\resultado \textbf{La fracción es a lo sumo $(n-1)/(nk^{2})<1/k^{2}$}; para $k=2$, al menos el $75\,\%$ de los datos está a menos de dos desviaciones típicas de la media, sea cual sea la forma de los datos. $\blacksquare$
\comprobacion en el \cref{ej:t1-tokens}, $\xbar\pm2s=(-1{,}45;\allowbreak\ 36{,}2)$ deja fuera solo el $40$: $1/8=12{,}5\,\%\le7/32\approx21{,}9\,\%$.
\end{solucion}

\begin{solucion}{eje:t1-welford}
\planteamiento escribir $m\xbar_m=(m-1)\xbar_{m-1}+x_m$ y aplicar la \cref{prop:t1-extremal}(ii) a los $m-1$ primeros datos con $a=\xbar_m$.
\pasos
\begin{enumerate}[label=\arabic*.]
  \item Media: $\xbar_m=\dfrac{(m-1)\xbar_{m-1}+x_m}{m}=\xbar_{m-1}+\dfrac{x_m-\xbar_{m-1}}{m}$.
  \item Por la \cref{prop:t1-extremal}(ii), $\sum_{i=1}^{m-1}(x_i-\xbar_m)^{2}=M_{m-1}+(m-1)(\xbar_{m-1}-\xbar_m)^{2}$. Sumando el término $i=m$: $M_m=M_{m-1}+(m-1)(\xbar_{m-1}-\xbar_m)^{2}+(x_m-\xbar_m)^{2}$.
  \item Sea $d=x_m-\xbar_{m-1}$. Por el paso 1, $\xbar_m-\xbar_{m-1}=d/m$ y $x_m-\xbar_m=d-d/m=d(m-1)/m$.
  \item Sustituyendo: $(m-1)\dfrac{d^{2}}{m^{2}}+\dfrac{(m-1)^{2}d^{2}}{m^{2}}=\dfrac{(m-1)d^{2}(1+m-1)}{m^{2}}=\dfrac{(m-1)d^{2}}{m}=d\cdot\dfrac{(m-1)d}{m}=(x_m-\xbar_{m-1})(x_m-\xbar_m)$.
  \item Datos $2,4,9$: $\xbar_1=2$, $M_1=0$; $\xbar_2=2+2/2=3$, $M_2=0+(4-2)(4-3)=2$; $\xbar_3=3+6/3=5$, $M_3=2+(9-3)(9-5)=26$.
\end{enumerate}
\resultado \textbf{Las dos fórmulas son ciertas}; con $2,4,9$ se obtiene $\xbar_3=5$, $M_3=26$ y $s^{2}=26/2=13$. $\blacksquare$
\comprobacion directamente: $(2-5)^{2}+(4-5)^{2}+(9-5)^{2}=9+1+16=26$.
\end{solucion}

\begin{solucion}{eje:t1-varianza-pares}
\planteamiento desarrollar la suma sobre todos los pares ordenados y usar $\sum_i(x_i-\xbar)=0$.
\pasos
\begin{enumerate}[label=\arabic*.]
  \item Sean $u_i=x_i-\xbar$; como $x_i-x_j=u_i-u_j$, se trabaja con los $u_i$, que suman cero.
  \item Sumando sobre todos los pares ordenados $(i,j)$, incluidos $i=j$ (que aportan $0$):
  \[
  \sum_{i=1}^{n}\sum_{j=1}^{n}(u_i-u_j)^{2}=\sum_{i,j}u_i^{2}+\sum_{i,j}u_j^{2}-2\Bigl(\sum_iu_i\Bigr)\Bigl(\sum_ju_j\Bigr)=2n\sum_iu_i^{2}-0.
  \]
  \item Cada par no ordenado $\{i,j\}$ con $i\neq j$ aparece dos veces, luego $\sum_{i<j}(x_i-x_j)^{2}=n\sum_iu_i^{2}=n(n-1)s^{2}$.
  \item Dividiendo por $n(n-1)$ se obtiene la identidad.
\end{enumerate}
\resultado \textbf{$s^{2}=\frac{1}{n(n-1)}\sum_{i<j}(x_i-x_j)^{2}$.} $\blacksquare$
\comprobacion con $2,4,9$: $(2-4)^{2}+(2-9)^{2}+(4-9)^{2}=4+49+25=78$ y $78/(3\cdot2)=13$, que coincide con el \cref{eje:t1-welford}.
\end{solucion}

% ---------------------------------------------------------------------
\section{Tema 2}

\begin{solucion}{eje:t2-conceptos}
\planteamiento aplicar las \cref{def:t2-mas,def:t2-parametro,def:t2-estimador}.
\pasos
\begin{enumerate}[label=\arabic*.]
  \item Población: el banco de un millón de preguntas (o, mejor, la variable «acierto» en cada una de ellas).
  \item Muestra: las $500$ preguntas elegidas al azar, con variables $X_1,\dots,X_{500}$, donde $X_i=1$ si el modelo acierta la pregunta $i$.
  \item Parámetro: la exactitud real $p$, fracción de aciertos en todo el banco; es un número fijo y desconocido.
  \item Estimador: $\widehat p=\Xbar=\frac{1}{500}\sum_iX_i$, una variable aleatoria. Estimación: su valor observado $410/500=0{,}82$.
\end{enumerate}
\resultado \textbf{El número $0{,}82$ no es aleatorio (es un valor observado); $\widehat p$ sí lo es}, porque con otra muestra de $500$ preguntas tomaría otro valor.
\comprobacion la muestra ($500$) es mucho menor que la población ($10^{6}$), de modo que el muestreo sin reemplazamiento se puede tratar como \iid\ (\cref{def:t2-mas}).
\end{solucion}

\begin{solucion}{eje:t2-enumeracion}
\datos población $\{1,3,5\}$ equiprobable; $n=2$ con reemplazamiento.
\planteamiento enumerar las $3^{2}=9$ muestras equiprobables y calcular $\xbar$ y $s^{2}=(x_1-x_2)^{2}/2$ en cada una.
\pasos
\begin{enumerate}[label=\arabic*.]
  \item $\mu=(1+3+5)/3=3$; $\sigma^{2}=\bigl((1-3)^{2}+0+(5-3)^{2}\bigr)/3=8/3$.
  \item Muestras y medias: $(1,1)\to1$; $(1,3),(3,1)\to2$; $(1,5),(3,3),(5,1)\to3$; $(3,5),(5,3)\to4$; $(5,5)\to5$. Distribución de $\Xbar$: valores $1,2,3,4,5$ con probabilidades $\frac19,\frac29,\frac39,\frac29,\frac19$.
  \item $\E[\Xbar]=(1+4+9+8+5)/9=27/9=3=\mu$.
  \item $\Var(\Xbar)=\bigl(4\cdot1+1\cdot2+0\cdot3+1\cdot2+4\cdot1\bigr)/9=12/9=4/3=\sigma^{2}/2$.
  \item Para $n=2$, $s^{2}=\bigl((x_1-\xbar)^{2}+(x_2-\xbar)^{2}\bigr)/1=(x_1-x_2)^{2}/2$. Vale $0$ en las tres muestras con $x_1=x_2$, $2$ en las cuatro con $\abs{x_1-x_2}=2$ y $8$ en las dos con $\abs{x_1-x_2}=4$.
  \item $\E[S^{2}]=(0\cdot3+2\cdot4+8\cdot2)/9=24/9=8/3=\sigma^{2}$.
  \item $\widehat\sigma^{2}=S^{2}/2$, luego $\E[\widehat\sigma^{2}]=4/3$.
\end{enumerate}
\resultado \textbf{$\mu=3$, $\sigma^{2}=8/3$; $\E[\Xbar]=3$, $\Var(\Xbar)=4/3$; $\E[S^{2}]=8/3$; $\E[\widehat\sigma^{2}]=4/3$}, la mitad de $\sigma^{2}$.
\comprobacion el \cref{cor:t2-sesgo} predice $\E[\widehat\sigma^{2}]=\frac{n-1}{n}\sigma^{2}=\frac12\cdot\frac83=\frac43$.
\end{solucion}

\begin{solucion}{eje:t2-tres-estimadores}
\datos $X_1,X_2,X_3$ \iid, media $\mu$, varianza $\sigma^{2}$.
\planteamiento para $T=\sum_ia_iX_i$: $\E[T]=\mu\sum_ia_i$ y $\Var(T)=\sigma^{2}\sum_ia_i^{2}$ (\cref{prop:t2-props}); después, \cref{prop:t2-ecm}.
\pasos
\begin{enumerate}[label=\arabic*.]
  \item $T_1$: $a=(\frac13,\frac13,\frac13)$; $\sum a_i=1$, sesgo $0$; $\sum a_i^{2}=\frac13$; $\ECM=\Var=\sigma^{2}/3$.
  \item $T_2$: $a=(\frac14,\frac12,\frac14)$; $\sum a_i=1$, sesgo $0$; $\sum a_i^{2}=\frac1{16}+\frac14+\frac1{16}=\frac38$; $\ECM=3\sigma^{2}/8$.
  \item $T_3$: $a=(1,1,-1)$; $\sum a_i=1$, sesgo $0$; $\sum a_i^{2}=3$; $\ECM=3\sigma^{2}$.
  \item $T_4$: $a=(\frac14,\frac14,\frac14)$; $\E[T_4]=\frac34\mu$, sesgo $-\mu/4$; $\Var=\frac3{16}\sigma^{2}$; $\ECM=\frac{3\sigma^{2}}{16}+\frac{\mu^{2}}{16}$.
  \item Entre los insesgados, $\frac13<\frac38<3$: el mejor es $T_1$.
  \item $\ECM(T_4)<\ECM(T_1)\iff\frac{3\sigma^{2}+\mu^{2}}{16}<\frac{\sigma^{2}}{3}\iff\mu^{2}<\frac{16\sigma^{2}}{3}-3\sigma^{2}=\frac{7\sigma^{2}}{3}$.
\end{enumerate}
\resultado \textbf{Sesgos $0,0,0,-\mu/4$; varianzas $\sigma^{2}/3$, $3\sigma^{2}/8$, $3\sigma^{2}$, $3\sigma^{2}/16$; el mejor insesgado es $\Xbar$; $T_4$ es mejor que $\Xbar$ si y solo si $\abs\mu<\sigma\sqrt{7/3}\approx1{,}528\,\sigma$.}
\comprobacion para $\mu=0$, $T_4$ es insesgado y tiene varianza $3\sigma^{2}/16<\sigma^{2}/3$: encoger hacia el cero ayuda cuando el parámetro está cerca de cero.
\end{solucion}

\begin{solucion}{eje:t2-poisson}
\datos $X_i\sim\Pois(\lambda)$, $\E X_i=\Var X_i=\lambda$; $n=50$, $\xbar=3{,}2$.
\planteamiento \cref{teo:t2-media} con $\mu=\sigma^{2}=\lambda$.
\pasos
\begin{enumerate}[label=\arabic*.]
  \item $\E[\widehat\lambda]=\lambda$: insesgado.
  \item $\Var(\widehat\lambda)=\lambda/n$; $\ee(\widehat\lambda)=\sqrt{\lambda/n}$; $\ECM=\Var=\lambda/n$.
  \item Estimación: $\widehat\lambda=3{,}2$. Error estándar estimado: $\sqrt{3{,}2/50}=\sqrt{0{,}064}\approx0{,}2530$.
\end{enumerate}
\resultado \textbf{$\widehat\lambda$ es insesgado, con $\Var=\ECM=\lambda/n$ y $\ee=\sqrt{\lambda/n}$; con los datos, $\widehat\lambda=3{,}2$ y $\widehat{\ee}\approx0{,}253$.}
\comprobacion el error estándar es pequeño frente a $3{,}2$ (un $8\,\%$), coherente con $n=50$ bloques.
\end{solucion}

\begin{solucion}{eje:t2-expo-sesgo}
\datos $\E[1/\Xbar]=n\lambda/(n-1)$; $n=5$; $\sum x_i=10$.
\planteamiento sesgo $=\E[\widehat\lambda]-\lambda$; elegir $c$ para que $\E[c/\Xbar]=\lambda$.
\pasos
\begin{enumerate}[label=\arabic*.]
  \item $\sesgo(1/\Xbar)=\dfrac{n\lambda}{n-1}-\lambda=\dfrac{\lambda}{n-1}>0$: sobrestima $\lambda$.
  \item $\E[c/\Xbar]=c\,n\lambda/(n-1)=\lambda\iff c=(n-1)/n$. El estimador insesgado es $\dfrac{n-1}{n\Xbar}=\dfrac{n-1}{\sum_iX_i}$.
  \item Datos: $\xbar=2$; $1/\xbar=0{,}5$; insesgado: $4/10=0{,}4$.
\end{enumerate}
\resultado \textbf{Sesgo $\lambda/(n-1)$; estimador insesgado $(n-1)/\sum_iX_i$; estimaciones $0{,}5$ y $0{,}4$.}
\comprobacion el sesgo tiende a cero cuando $n\to\infty$ y la corrección $(n-1)/n\to1$; con $n=5$ la diferencia es del $20\,\%$.
\end{solucion}

\begin{solucion}{eje:t2-sumas-solapadas}
\datos $\Var(X_i)=1/4$; $X=\sum_{i=1}^{n}X_i$; $Y=\sum_{i=n-7}^{2n-8}X_i$.
\planteamiento bilinealidad de la covarianza (\cref{prop:t2-props}(v)): $\Cov(X_i,X_j)=0$ si $i\ne j$ y $\Var(X_i)$ si $i=j$; solo cuentan los índices comunes.
\pasos
\begin{enumerate}[label=\arabic*.]
  \item $Y$ tiene $(2n-8)-(n-7)+1=n$ sumandos, luego $\Var(Y)=n/4=\Var(X)$.
  \item Índices comunes: de $n-7$ a $n$, que son $8$ (como $n\ge8$, $n-7\ge1$).
  \item $\Cov(X,Y)=\sum_{i=n-7}^{n}\Var(X_i)=8\cdot\frac14=2$.
  \item $\Cor(X,Y)=\dfrac{2}{\sqrt{(n/4)(n/4)}}=\dfrac{2}{n/4}=\dfrac8n$.
\end{enumerate}
\resultado \textbf{$\Cov(X,Y)=2$ y $\Cor(X,Y)=8/n$.}
\comprobacion para $n=8$ se obtiene $\Cor=1$: entonces $Y=X_1+\dots+X_8=X$.
\end{solucion}

\begin{solucion}{eje:t2-benchmark}
\datos $n=400$; $300$ aciertos.
\planteamiento \cref{ej:t2-proporcion}: $\widehat{\ee}=\sqrt{\widehat p(1-\widehat p)/n}$.
\pasos
\begin{enumerate}[label=\arabic*.]
  \item $\widehat p=300/400=0{,}75$.
  \item $\widehat{\ee}=\sqrt{0{,}75\cdot0{,}25/400}=\sqrt{0{,}000\,468\,75}\approx0{,}02165$.
  \item El error estándar es proporcional a $1/\sqrt n$; para dividirlo por $2$ se necesita $4n=1600$ preguntas.
\end{enumerate}
\resultado \textbf{$\widehat p=0{,}75$, $\widehat{\ee}\approx0{,}0217$ (unos $2{,}2$ puntos porcentuales); $1600$ preguntas.}
\comprobacion $\sqrt{0{,}1875/1600}\approx0{,}01083$, la mitad de $0{,}02165$.
\end{solucion}

\begin{solucion}{eje:t2-combinar}
\datos $\E A=\E B=\theta$; $A$, $B$ independientes; $\Var A=4$, $\Var B=1$.
\planteamiento linealidad para el sesgo; varianza de una combinación de independientes; minimizar un polinomio de grado~2.
\pasos
\begin{enumerate}[label=\arabic*.]
  \item $\E[T_w]=w\theta+(1-w)\theta=\theta$ para todo $w$.
  \item $V(w)=\Var(T_w)=4w^{2}+(1-w)^{2}$. $V'(w)=8w-2(1-w)=10w-2=0\iff w=0{,}2$; $V''=10>0$, mínimo.
  \item $V(0{,}2)=4\cdot0{,}04+0{,}64=0{,}8$.
  \item En general, $V(w)=w^{2}\sigma_A^{2}+(1-w)^{2}\sigma_B^{2}$ y $V'(w)=0$ da $w^{*}=\dfrac{\sigma_B^{2}}{\sigma_A^{2}+\sigma_B^{2}}$, con $V(w^{*})=\dfrac{\sigma_A^{2}\sigma_B^{2}}{\sigma_A^{2}+\sigma_B^{2}}$.
\end{enumerate}
\resultado \textbf{$w^{*}=0{,}2$ y $\Var_{\min}=0{,}8$}; en general, cada estimador se pondera de forma inversamente proporcional a su varianza y $1/\Var_{\min}=1/\sigma_A^{2}+1/\sigma_B^{2}$.
\comprobacion $1/4+1/1=1{,}25=1/0{,}8$. La combinación es mejor que el mejor de los dos ($0{,}8<1$).
\end{solucion}

\begin{solucion}{eje:t2-chebyshev-n}
\datos $\varepsilon=0{,}05$; probabilidad máxima $0{,}05$; $\Var(\widehat p)=p(1-p)/n\le1/(4n)$.
\planteamiento \cref{lem:t2-chebyshev} con $W=\widehat p$.
\pasos
\begin{enumerate}[label=\arabic*.]
  \item $\PP(\abs{\widehat p-p}\ge0{,}05)\le\dfrac{p(1-p)}{n\cdot0{,}0025}\le\dfrac{1}{4n\cdot0{,}0025}=\dfrac{100}{n}$.
  \item $100/n\le0{,}05\iff n\ge2000$.
\end{enumerate}
\resultado \textbf{$n=2000$ preguntas bastan}, sea cual sea $p$.
\comprobacion la aproximación normal del \cref{tema:3} dará $n=385$ (\cref{eje:t3-tamano}); Chebyshev es válida para cualquier distribución, pero es muy conservadora.
\end{solucion}

\begin{solucion}{eje:t2-ecm-normal}
\datos $\E[S^{2}]=\sigma^{2}$; $\Var(S^{2})=2\sigma^{4}/(n-1)$; $\widehat\sigma^{2}=\frac{n-1}{n}S^{2}$.
\planteamiento \cref{prop:t2-ecm} para cada estimador y comparación de las dos expresiones.
\pasos
\begin{enumerate}[label=\arabic*.]
  \item $S^{2}$ es insesgado, luego $\ECM(S^{2})=\dfrac{2\sigma^{4}}{n-1}$.
  \item $\sesgo(\widehat\sigma^{2})=-\sigma^{2}/n$ (\cref{cor:t2-sesgo}) y $\Var(\widehat\sigma^{2})=\dfrac{(n-1)^{2}}{n^{2}}\cdot\dfrac{2\sigma^{4}}{n-1}=\dfrac{2(n-1)\sigma^{4}}{n^{2}}$.
  \item $\ECM(\widehat\sigma^{2})=\dfrac{2(n-1)\sigma^{4}}{n^{2}}+\dfrac{\sigma^{4}}{n^{2}}=\dfrac{(2n-1)\sigma^{4}}{n^{2}}$.
  \item Comparación: $\dfrac{2n-1}{n^{2}}<\dfrac{2}{n-1}\iff(2n-1)(n-1)<2n^{2}\iff2n^{2}-3n+1<2n^{2}\iff1<3n$, que se cumple para todo $n\ge1$.
\end{enumerate}
\resultado \textbf{$\ECM(S^{2})=2\sigma^{4}/(n-1)$ y $\ECM(\widehat\sigma^{2})=(2n-1)\sigma^{4}/n^{2}$, que es menor para todo $n\ge2$.} $\blacksquare$
\comprobacion con $n=2$: $\ECM(S^{2})=2\sigma^{4}$ y $\ECM(\widehat\sigma^{2})=\frac34\sigma^{4}$. La insesgadez no garantiza el menor error cuadrático medio.
\end{solucion}

\begin{solucion}{eje:t2-uniforme}
\datos $X_i\sim\Unif(0,\theta)$: densidad $1/\theta$ en $[0,\theta]$, media $\theta/2$, varianza $\theta^{2}/12$.
\planteamiento (a) \cref{teo:t2-media}; (b) la función de distribución del máximo por independencia; (c) comparar.
\pasos
\begin{enumerate}[label=\arabic*.]
  \item $\E[2\Xbar]=2\cdot\theta/2=\theta$; $\Var(2\Xbar)=4\cdot\dfrac{\theta^{2}}{12n}=\dfrac{\theta^{2}}{3n}$.
  \item $M\le x$ si y solo si todos los $X_i\le x$; por independencia, $\PP(M\le x)=\prod_i\PP(X_i\le x)=(x/\theta)^{n}$ para $0\le x\le\theta$. Derivando, la densidad de $M$ es $f_M(x)=nx^{n-1}/\theta^{n}$ en $[0,\theta]$.
  \item $\E[M]=\displaystyle\int_0^{\theta}x\,\frac{nx^{n-1}}{\theta^{n}}\,dx=\frac{n}{\theta^{n}}\cdot\frac{\theta^{n+1}}{n+1}=\frac{n\theta}{n+1}$; análogamente $\E[M^{2}]=\dfrac{n}{\theta^{n}}\cdot\dfrac{\theta^{n+2}}{n+2}=\dfrac{n\theta^{2}}{n+2}$.
  \item $\E[T_2]=\frac{n+1}{n}\cdot\frac{n\theta}{n+1}=\theta$: insesgado.
  \item $\Var(M)=\dfrac{n\theta^{2}}{n+2}-\dfrac{n^{2}\theta^{2}}{(n+1)^{2}}=\dfrac{n\theta^{2}\bigl[(n+1)^{2}-n(n+2)\bigr]}{(n+2)(n+1)^{2}}=\dfrac{n\theta^{2}}{(n+2)(n+1)^{2}}$.
  \item $\Var(T_2)=\dfrac{(n+1)^{2}}{n^{2}}\cdot\dfrac{n\theta^{2}}{(n+2)(n+1)^{2}}=\dfrac{\theta^{2}}{n(n+2)}$.
  \item $\dfrac{\theta^{2}}{n(n+2)}\le\dfrac{\theta^{2}}{3n}\iff3\le n+2\iff n\ge1$, con igualdad solo si $n=1$.
\end{enumerate}
\resultado \textbf{Ambos son insesgados; $\Var(T_1)=\theta^{2}/(3n)$ y $\Var(T_2)=\theta^{2}/(n(n+2))$.} $T_2$ es mejor para $n\ge2$ y su varianza decrece como $1/n^{2}$, no como $1/n$. $\blacksquare$
\comprobacion con $n=1$, $M=X_1$ y $T_2=2X_1=T_1$: deben coincidir, y en efecto $\theta^{2}/3=\theta^{2}/(1\cdot3)$.
\end{solucion}

\begin{solucion}{eje:t2-max-consistente}
\datos $\PP(M\le x)=(x/\theta)^{n}$; $\ECM(M)$ a partir del \cref{eje:t2-uniforme}.
\planteamiento (a) como $M\le\theta$, el suceso $\abs{M-\theta}>\varepsilon$ es $M<\theta-\varepsilon$; (b) calcular el $\ECM$.
\pasos
\begin{enumerate}[label=\arabic*.]
  \item Si $\varepsilon\ge\theta$, la probabilidad es $0$. Si $0<\varepsilon<\theta$: $\PP(\abs{M-\theta}>\varepsilon)=\PP(M<\theta-\varepsilon)=\bigl(1-\varepsilon/\theta\bigr)^{n}\to0$, porque $0<1-\varepsilon/\theta<1$.
  \item Sesgo: $\E[M]-\theta=-\theta/(n+1)\to0$. Varianza: $\dfrac{n\theta^{2}}{(n+2)(n+1)^{2}}\le\dfrac{\theta^{2}}{(n+1)^{2}}\to0$. Luego $\ECM(M)\to0$ y la \cref{prop:t2-ecm-consistencia} da la consistencia.
\end{enumerate}
\resultado \textbf{$M\convp\theta$}: el máximo es un estimador consistente (aunque sesgado). $\blacksquare$
\comprobacion con $\theta=1$, $\varepsilon=0{,}05$ y $n=100$: $0{,}95^{100}\approx0{,}0059$.
\end{solucion}

\begin{solucion}{eje:t2-cov-desv}
\planteamiento bilinealidad de la covarianza y $\Cov(X_j,X_i)=\sigma^{2}$ si $j=i$ y $0$ si $j\neq i$.
\pasos
\begin{enumerate}[label=\arabic*.]
  \item $\Cov(\Xbar,X_i)=\frac1n\sum_{j=1}^{n}\Cov(X_j,X_i)=\frac1n\sigma^{2}$.
  \item $\Cov(\Xbar,\Xbar)=\Var(\Xbar)=\sigma^{2}/n$.
  \item $\Cov(\Xbar,X_i-\Xbar)=\Cov(\Xbar,X_i)-\Cov(\Xbar,\Xbar)=\sigma^{2}/n-\sigma^{2}/n=0$.
\end{enumerate}
\resultado \textbf{$\Cov(\Xbar,X_i-\Xbar)=0$ para todo $i$.} $\blacksquare$
\comprobacion la incorrelación no implica en general independencia; para datos normales, el vector $(\Xbar,X_1-\Xbar,\dots,X_n-\Xbar)$ es normal multivariante y en ese caso sí la implica, lo que lleva a la independencia de $\Xbar$ y $S^{2}$ (\cref{teo:t3-fisher}).
\end{solucion}

\begin{solucion}{eje:t2-s-sesgo}
\planteamiento $\Var(S)=\E[S^{2}]-(\E S)^{2}\ge0$ y $\E[S^{2}]=\sigma^{2}$ (\cref{teo:t2-s2}).
\pasos
\begin{enumerate}[label=\arabic*.]
  \item $0\le\Var(S)=\sigma^{2}-(\E S)^{2}$, luego $(\E S)^{2}\le\sigma^{2}$.
  \item Como $S\ge0$, $\E S\ge0$, y tomando raíces: $\E S\le\sigma$.
  \item Igualdad $\iff\Var(S)=0\iff S$ es constante con probabilidad~1.
\end{enumerate}
\resultado \textbf{$\E[S]\le\sigma$}, con igualdad solo en el caso degenerado; $S$ subestima $\sigma$ en promedio. $\blacksquare$
\comprobacion con la población del \cref{eje:t2-enumeracion}: $S=\sqrt{S^{2}}$ vale $0$, $\sqrt2$ y $\sqrt8$ con probabilidades $3/9$, $4/9$ y $2/9$, luego $\E S=(4\sqrt2+2\sqrt8)/9=8\sqrt2/9\approx1{,}257<\sigma=\sqrt{8/3}\approx1{,}633$.
\end{solucion}

\begin{solucion}{eje:t2-passk}
\datos $n$ soluciones independientes, cada una correcta con probabilidad $p$; $C$ correctas.
\planteamiento introducir un subconjunto aleatorio $K$ de $k$ índices, elegido uniformemente entre los $\binom nk$ posibles e independiente de las soluciones, y aplicar la ley de la probabilidad total de dos formas.
\pasos
\begin{enumerate}[label=\arabic*.]
  \item Sea $A$ el suceso «todas las soluciones de índice en $K$ son incorrectas».
  \item Condicionando a los resultados de las $n$ soluciones (y, por tanto, a $C$), $K$ es un subconjunto uniforme, y el número de subconjuntos de $k$ índices formados solo por incorrectas es $\binom{n-C}{k}$ (es $0$ si $n-C<k$). Luego $\PP(A\mid\text{resultados})=\binom{n-C}{k}/\binom nk$.
  \item Tomando esperanzas (probabilidad total): $\PP(A)=\E\Bigl[\binom{n-C}{k}\big/\binom nk\Bigr]$.
  \item Condicionando ahora a $K=J$ para un subconjunto fijo $J$: las $k$ soluciones de $J$ son independientes y cada una es incorrecta con probabilidad $1-p$, luego $\PP(A\mid K=J)=(1-p)^{k}$ para todo $J$, y $\PP(A)=(1-p)^{k}$.
  \item Igualando 3 y 4: $\E\bigl[\binom{n-C}{k}/\binom nk\bigr]=(1-p)^{k}$, y por linealidad $\E[\widehat\pi_k]=1-(1-p)^{k}$.
\end{enumerate}
\resultado \textbf{$\widehat\pi_k$ es un estimador insesgado de pass@$k$.} $\blacksquare$
\comprobacion para $k=1$, $\widehat\pi_1=1-(n-C)/n=C/n$, la proporción muestral, que es insesgada para $p$. Numéricamente, con $n=10$, $k=3$ y $p=0{,}3$, la esperanza calculada sumando sobre $C=0,\dots,10$ vale $0{,}657=1-0{,}7^{3}$.
\end{solucion}

% ---------------------------------------------------------------------
\section{Tema 3}

\begin{solucion}{eje:t3-vf}
\planteamiento contrastar con el \cref{teo:t3-tcl}, el \cref{teo:t3-suma-normales}, el \cref{cor:t3-binomial} y la \cref{def:t3-chi-t}.
\pasos
\begin{enumerate}[label=(\alph*)]
  \item \textbf{Falsa.} Los datos siguen siendo exponenciales (asimétricos); lo que es aproximadamente normal es la distribución de $\Xbar$.
  \item \textbf{Verdadera}, por el \cref{teo:t3-suma-normales}.
  \item \textbf{Verdadera}: $\ee(\Xbar)=\sigma/\sqrt{100}=\sigma/10$.
  \item \textbf{Verdadera}: por ejemplo, $t_{10;0{,}975}=2{,}228>1{,}960=z_{0{,}975}$, de modo que $\PP(T_{10}>1{,}96)>0{,}025$.
  \item \textbf{Falsa}: $np=0{,}4$, muy lejos de $np\ge10$; la distribución es muy asimétrica (de hecho, $\PP(X=0)=0{,}98^{20}\approx0{,}668$).
\end{enumerate}
\resultado \textbf{(a) F; (b) V; (c) V; (d) V; (e) F.}
\end{solucion}

\begin{solucion}{eje:t3-tabla}
\planteamiento leer la \cref{tab:phi} y usar $\Phi(-z)=1-\Phi(z)$.
\pasos
\begin{enumerate}[label=(\alph*)]
  \item Fila $1{,}2$, columna $0{,}03$: $\Phi(1{,}23)=0{,}8907$.
  \item $\PP(Z>-0{,}5)=1-\Phi(-0{,}5)=\Phi(0{,}5)=0{,}6915$.
  \item $\Phi(2)-\Phi(-1)=0{,}9772-(1-0{,}8413)=0{,}8185$.
  \item $z_{0{,}95}=1{,}6449$ y $z_{0{,}99}=2{,}3263$ (\cref{tab:zq}).
  \item $\PP(\abs Z\le c)=2\Phi(c)-1=0{,}90\iff\Phi(c)=0{,}95\iff c=z_{0{,}95}=1{,}6449$.
\end{enumerate}
\resultado \textbf{(a) $0{,}8907$; (b) $0{,}6915$; (c) $0{,}8185$; (d) $1{,}6449$ y $2{,}3263$; (e) $c\approx1{,}645$.}
\comprobacion en (e), cada cola tiene probabilidad $0{,}05$, y por eso aparece el cuantil $0{,}95$ y no el $0{,}90$.
\end{solucion}

\begin{solucion}{eje:t3-normal-basico}
\datos $X\sim\Nor(100,15^{2})$.
\planteamiento tipificar (\cref{prop:t3-tipificacion}).
\pasos
\begin{enumerate}[label=\arabic*.]
  \item $\PP(X>130)=1-\Phi\bigl(\frac{130-100}{15}\bigr)=1-\Phi(2)=1-0{,}9772=0{,}0228$.
  \item $\PP(85<X<115)=\Phi(1)-\Phi(-1)=2\Phi(1)-1=2\cdot0{,}8413-1=0{,}6826$.
  \item Percentil $90$: $x=100+15\,z_{0{,}90}=100+15\cdot1{,}2816=119{,}22$.
\end{enumerate}
\resultado \textbf{$0{,}0228$; $0{,}683$; $119{,}2$.}
\comprobacion los dos primeros resultados reproducen la regla empírica ($68\,\%$ a una desviación típica; $2{,}3\,\%$ por encima de dos).
\end{solucion}

\begin{solucion}{eje:t3-openstax-media}
\datos $\mu=90$, $\sigma=15$, $n=25$.
\planteamiento por el teorema central del límite, $\Xbar\approx\Nor(90,\,3^{2})$, pues $\ee=15/\sqrt{25}=3$.
\pasos
\begin{enumerate}[label=\arabic*.]
  \item $z_1=(85-90)/3\approx-1{,}67$; $z_2=(92-90)/3\approx0{,}67$.
  \item $\PP(85<\Xbar<92)\approx\Phi(0{,}67)-\Phi(-1{,}67)=0{,}7486-0{,}0475=0{,}7011$.
  \item Dos errores estándar por encima: $90+2\cdot3=96$.
\end{enumerate}
\resultado \textbf{(a) $\approx0{,}70$; (b) $96$.}
\comprobacion sin redondear $z$ a dos decimales se obtiene $0{,}6997$, que es el valor de la fuente.
\end{solucion}

\begin{solucion}{eje:t3-openstax-suma}
\datos $\mu=90$, $\sigma=15$, $n=80$.
\planteamiento $S_n\approx\Nor(n\mu,n\sigma^{2})$ (\cref{obs:t3-uso}).
\pasos
\begin{enumerate}[label=\arabic*.]
  \item $\E S_n=80\cdot90=7200$; $\sqrt{\Var S_n}=\sqrt{80}\cdot15\approx134{,}16$.
  \item $z=(7500-7200)/134{,}16\approx2{,}24$; $\PP(S_n>7500)\approx1-\Phi(2{,}24)=1-0{,}9875=0{,}0125$.
  \item $7200+1{,}5\cdot134{,}16\approx7401{,}2$.
\end{enumerate}
\resultado \textbf{(a) $\approx0{,}0125$; (b) $\approx7401$.}
\comprobacion con $z=2{,}236$ sin redondear se obtiene $0{,}0127$, el valor de la fuente.
\end{solucion}

\begin{solucion}{eje:t3-elecciones}
\datos $n=400$; $p_A=0{,}5$; $p_{\text{menores}}=0{,}3$.
\planteamiento $\widehat p\approx\Nor\bigl(p,p(1-p)/n\bigr)$ (\cref{cor:t3-binomial}).
\pasos
\begin{enumerate}[label=\arabic*.]
  \item (a) $\ee=\sqrt{0{,}25/400}=0{,}025$; $z=(0{,}525-0{,}5)/0{,}025=1$; $\PP(\widehat p\ge0{,}525)\approx1-\Phi(1)=0{,}1587$.
  \item (b) $\ee=\sqrt{0{,}21/400}\approx0{,}02291$; $z=(0{,}31-0{,}3)/0{,}02291\approx0{,}44$; $\PP(\widehat p<0{,}31)\approx\Phi(0{,}44)=0{,}6700$.
\end{enumerate}
\resultado \textbf{(a) $\approx0{,}16$; (b) $\approx0{,}67$.}
\comprobacion en (a), con corrección por continuidad, $\PP(X\ge210)\approx1-\Phi\bigl((209{,}5-200)/10\bigr)=1-\Phi(0{,}95)=0{,}1711$, que coincide con el valor binomial exacto ($0{,}1711$).
\end{solucion}

\begin{solucion}{eje:t3-redondeo}
\datos errores $E_i$ uniformes en $\{-2,-1,0,1,2\}$, independientes; $n=1000$.
\planteamiento calcular media y varianza de cada error y aplicar el teorema central del límite a la suma $S$.
\pasos
\begin{enumerate}[label=\arabic*.]
  \item $\E E_i=0$; $\Var E_i=\E E_i^{2}=(4+1+0+1+4)/5=2$.
  \item $\E S=0$; $\Var S=1000\cdot2=2000$; desviación típica $\sqrt{2000}\approx44{,}72$.
  \item $\PP(\abs S>100)\approx2\bigl(1-\Phi(100/44{,}72)\bigr)=2\bigl(1-\Phi(2{,}24)\bigr)=2\cdot0{,}0125=0{,}025$.
\end{enumerate}
\resultado \textbf{$p\approx0{,}025$.}
\comprobacion como $S$ es entera, $\abs S>100$ equivale a $\abs S\ge101$, y con corrección por continuidad se obtiene $2\bigl(1-\Phi(100{,}5/44{,}72)\bigr)\approx0{,}0246$; la diferencia es pequeña.
\end{solucion}

\begin{solucion}{eje:t3-binomial}
\datos $np=50$, $\sqrt{np(1-p)}=5$.
\planteamiento \cref{cor:t3-binomial} y \cref{def:t3-cc}.
\pasos
\begin{enumerate}[label=\arabic*.]
  \item $\PP(X=50)=\PP(49{,}5\le Y\le50{,}5)\approx\Phi(0{,}1)-\Phi(-0{,}1)=2\cdot0{,}5398-1=0{,}0796$.
  \item Con corrección: $\PP(39{,}5\le Y\le60{,}5)\approx\Phi(2{,}1)-\Phi(-2{,}1)=2\cdot0{,}9821-1=0{,}9642$.
  \item Sin corrección: $\PP(40\le Y\le60)\approx2\Phi(2)-1=2\cdot0{,}9772-1=0{,}9544$.
\end{enumerate}
\resultado \textbf{(a) $0{,}0796$; (b) $0{,}9642$ con corrección y $0{,}9544$ sin ella.}
\comprobacion los valores exactos son $0{,}0796$ y $0{,}9648$: con corrección el error es de $0{,}0006$ y sin ella de $0{,}0104$. Sin corrección, $\PP(X=50)$ se aproximaría por $0$.
\end{solucion}

\begin{solucion}{eje:t3-benchmark}
\datos $p=0{,}8$; $n=1000$.
\planteamiento $\widehat p\approx\Nor\bigl(0{,}8,\ 0{,}16/1000\bigr)$.
\pasos
\begin{enumerate}[label=\arabic*.]
  \item $\ee=\sqrt{0{,}8\cdot0{,}2/1000}=\sqrt{0{,}000\,16}\approx0{,}01265$.
  \item $z=(0{,}82-0{,}8)/0{,}01265\approx1{,}58$; $\PP(\widehat p\ge0{,}82)\approx1-\Phi(1{,}58)=1-0{,}9429=0{,}0571$.
  \item Intervalo: $0{,}8\pm z_{0{,}975}\,\ee=0{,}8\pm1{,}96\cdot0{,}01265=0{,}8\pm0{,}0248$, es decir, $[0{,}775;\allowbreak\ 0{,}825]$.
\end{enumerate}
\resultado \textbf{(a) $0{,}0126$; (b) $\approx0{,}057$; (c) $[0{,}775;\allowbreak\,0{,}825]$.}
\comprobacion un modelo con exactitud real del $80\,\%$ obtiene un $82\,\%$ o más en casi un $6\,\%$ de las evaluaciones: una ventaja de dos puntos en un benchmark de $1000$ preguntas no es concluyente por sí sola.
\end{solucion}

\begin{solucion}{eje:t3-tamano}
\datos probabilidad $0{,}95$; tolerancia $0{,}1\sigma$; en (b), tolerancia $0{,}05$ y $\sigma=\sqrt{p(1-p)}\le1/2$.
\planteamiento $\PP(\abs{\Xbar-\mu}<c)\approx2\Phi\bigl(c\sqrt n/\sigma\bigr)-1\ge0{,}95\iff c\sqrt n/\sigma\ge z_{0{,}975}=1{,}96$.
\pasos
\begin{enumerate}[label=\arabic*.]
  \item (a) $c=0{,}1\sigma$: $0{,}1\sqrt n\ge1{,}96\iff n\ge(19{,}6)^{2}=384{,}16$, luego $n=385$.
  \item (b) $c=0{,}05$ y $\sigma\le0{,}5$: basta $0{,}05\sqrt n/0{,}5\ge1{,}96$, que es la misma condición que en (a): $n=385$.
\end{enumerate}
\resultado \textbf{$n=385$ en ambos casos}, frente a $n=2000$ con Chebyshev (\cref{eje:t2-chebyshev-n}).
\comprobacion con $n=385$ y $p=1/2$: $\ee=0{,}5/\sqrt{385}\approx0{,}02548$ y $1{,}96\cdot0{,}02548\approx0{,}0499<0{,}05$.
\end{solucion}

\begin{solucion}{eje:t3-t-cuantiles}
\planteamiento leer las \cref{tab:t,tab:chi2} y usar la simetría de la $t$.
\pasos
\begin{enumerate}[label=\arabic*.]
  \item $t_{9;0{,}975}=2{,}262$, $t_{29;0{,}975}=2{,}045$, $t_{120;0{,}975}=1{,}980$, frente a $z_{0{,}975}=1{,}960$.
  \item $t_{15;0{,}975}=2{,}131$, luego $\PP(T>2{,}131)=0{,}025$.
  \item $t_{15;0{,}95}=1{,}753$, luego $\PP(\abs T\le1{,}753)=1-2\cdot0{,}05=0{,}90$.
  \item $\PP(V>c)=0{,}05\iff c=\chi^{2}_{10;0{,}95}=18{,}307$.
\end{enumerate}
\resultado \textbf{(a) $2{,}262$; $2{,}045$; $1{,}980$ (decrecen hacia $1{,}960$); (b) $0{,}025$ y $0{,}90$; (c) $c=18{,}307$.}
\comprobacion la $\chi^{2}_{10}$ tiene media $10$ y desviación típica $\sqrt{20}\approx4{,}47$; $18{,}3$ está unas $1{,}9$ desviaciones por encima, razonable para un cuantil $0{,}95$ de una distribución asimétrica a la derecha.
\end{solucion}

\begin{solucion}{eje:t3-lineal-normal}
\planteamiento calcular la función de distribución de $Y$, derivarla y reconocer la densidad normal.
\pasos
\begin{enumerate}[label=\arabic*.]
  \item Caso $a>0$: $F_Y(y)=\PP(aX+b\le y)=\PP\bigl(X\le\frac{y-b}{a}\bigr)=F_X\bigl(\frac{y-b}{a}\bigr)$. Derivando (regla de la cadena), $f_Y(y)=\frac1af_X\bigl(\frac{y-b}{a}\bigr)$.
  \item Caso $a<0$: $F_Y(y)=\PP\bigl(X\ge\frac{y-b}{a}\bigr)=1-F_X\bigl(\frac{y-b}{a}\bigr)$, y $f_Y(y)=-\frac1af_X\bigl(\frac{y-b}{a}\bigr)$. En ambos casos, $f_Y(y)=\frac{1}{\abs a}f_X\bigl(\frac{y-b}{a}\bigr)$.
  \item Sustituyendo la densidad normal y usando $\frac{y-b}{a}-\mu=\frac{y-(a\mu+b)}{a}$:
  \[
  f_Y(y)=\frac{1}{\abs a\sigma\sqrt{2\pi}}\exp\Bigl(-\frac{(y-(a\mu+b))^{2}}{2a^{2}\sigma^{2}}\Bigr),
  \]
  que es la densidad $\Nor(a\mu+b,a^{2}\sigma^{2})$, con desviación típica $\abs a\sigma$.
  \item Con $a=1/\sigma$ y $b=-\mu/\sigma$ se obtiene $Z=(X-\mu)/\sigma\sim\Nor(0,1)$ y $\PP(X\le x)=\PP\bigl(Z\le\frac{x-\mu}{\sigma}\bigr)=\Phi\bigl(\frac{x-\mu}{\sigma}\bigr)$.
\end{enumerate}
\resultado \textbf{$aX+b\sim\Nor(a\mu+b,a^{2}\sigma^{2})$}, y la tipificación lleva cualquier normal a la estándar. $\blacksquare$
\comprobacion con $a=-1$, $b=0$, $\mu=0$, $\sigma=1$: $-Z\sim\Nor(0,1)$, que expresa la simetría de la normal estándar.
\end{solucion}

\begin{solucion}{eje:t3-chi1}
\planteamiento para $v>0$, $Z^{2}\le v\iff-\sqrt v\le Z\le\sqrt v$.
\pasos
\begin{enumerate}[label=\arabic*.]
  \item $\PP(V\le v)=\Phi(\sqrt v)-\Phi(-\sqrt v)=\Phi(\sqrt v)-\bigl(1-\Phi(\sqrt v)\bigr)=2\Phi(\sqrt v)-1$.
  \item Sea $\varphi(z)=\frac{1}{\sqrt{2\pi}}e^{-z^{2}/2}=\Phi'(z)$. Por la regla de la cadena,
  \[
  f_V(v)=2\varphi(\sqrt v)\cdot\frac{1}{2\sqrt v}=\frac{\varphi(\sqrt v)}{\sqrt v}=\frac{1}{\sqrt{2\pi v}}e^{-v/2},\qquad v>0.
  \]
  \item Para $v\le0$, $\PP(V\le v)=\PP(Z^{2}\le v)\le\PP(Z=0)=0$, y la densidad es nula.
\end{enumerate}
\resultado \textbf{$f_V(v)=(2\pi v)^{-1/2}e^{-v/2}$ para $v>0$.} $\blacksquare$
\comprobacion $\PP(V\le3{,}841)=2\Phi(\sqrt{3{,}841})-1=2\Phi(1{,}96)-1=0{,}95$, de acuerdo con $\chi^{2}_{1;0{,}95}=3{,}841$ en la \cref{tab:chi2}.
\end{solucion}

\begin{solucion}{eje:t3-chi-momentos}
\planteamiento $\varphi'(z)=-z\varphi(z)$; integrar por partes $\int z^{m}\cdot z\varphi(z)\,dz$; después, linealidad e independencia.
\pasos
\begin{enumerate}[label=\arabic*.]
  \item Para $m\ge0$ entero, con $u=z^{m+1}$ y $dv=z\varphi(z)\,dz$, $v=-\varphi(z)$:
  \[
  \begin{aligned}
  \E[Z^{m+2}]&=\int_{-\infty}^{\infty}z^{m+1}\,z\varphi(z)\,dz=\Bigl[-z^{m+1}\varphi(z)\Bigr]_{-\infty}^{\infty}+(m+1)\int_{-\infty}^{\infty}z^{m}\varphi(z)\,dz\\
  &=(m+1)\E[Z^{m}],
  \end{aligned}
  \]
  porque $z^{m+1}e^{-z^{2}/2}\to0$ cuando $\abs z\to\infty$.
  \item $m=0$: $\E[Z^{2}]=1\cdot\E[1]=1$. $m=2$: $\E[Z^{4}]=3\E[Z^{2}]=3$.
  \item $\Var(Z^{2})=\E[Z^{4}]-(\E Z^{2})^{2}=3-1=2$.
  \item $V=\sum_{i=1}^{k}Z_i^{2}$ con $Z_i$ \iid: $\E V=k\cdot1=k$ y, por independencia, $\Var V=k\cdot2=2k$.
\end{enumerate}
\resultado \textbf{$\E[Z^{4}]=3$; $\E[\chi^{2}_k]=k$ y $\Var(\chi^{2}_k)=2k$.} $\blacksquare$
\comprobacion para $k=10$, media $10$ y varianza $20$, coherente con el \cref{eje:t3-t-cuantiles}.
\end{solucion}

\begin{solucion}{eje:t3-t-simetrica}
\planteamiento escribir $T=Z/\sqrt{V/k}$ con $Z\sim\Nor(0,1)$ y $V\sim\chi^{2}_k$ independientes, y observar que $-T$ es la misma función aplicada a $(-Z,V)$.
\pasos
\begin{enumerate}[label=\arabic*.]
  \item $-Z\sim\Nor(0,1)$ (\cref{eje:t3-lineal-normal} con $a=-1$) y $-Z$ es independiente de $V$. Luego el par $(-Z,V)$ tiene la misma distribución conjunta que $(Z,V)$.
  \item Como $-T=(-Z)/\sqrt{V/k}$ es la misma función del par $(-Z,V)$ que $T$ del par $(Z,V)$, $-T$ y $T$ tienen la misma distribución.
  \item Sea $F$ la función de distribución de $T$, continua. Para todo $t$: $F(-t)=\PP(T\le-t)=\PP(-T\ge t)=\PP(T\ge t)=1-F(t)$.
  \item Con $t=t_{k;q}$: $F(-t_{k;q})=1-q$, es decir, $-t_{k;q}=t_{k;1-q}$.
  \item Con $c=t_{k;1-\alpha/2}$: $\PP(\abs T\le c)=F(c)-F(-c)=2F(c)-1=2(1-\alpha/2)-1=1-\alpha$.
\end{enumerate}
\resultado \textbf{$T\stackrel{d}{=}-T$, $t_{k;1-q}=-t_{k;q}$ y $\PP(\abs T\le t_{k;1-\alpha/2})=1-\alpha$.} $\blacksquare$
\comprobacion el mismo argumento con $V=k$ constante da la simetría de la normal estándar, cuyo caso límite es la $t_k$ cuando $k\to\infty$.
\end{solucion}

% ---------------------------------------------------------------------
\section{Tema 4}

\begin{solucion}{eje:t4-interpretacion}
\planteamiento aplicar la \cref{obs:t4-interpretacion}: la confianza se refiere al procedimiento, no a la probabilidad de un valor concreto del parámetro.
\pasos
\begin{enumerate}[label=(\roman*)]
  \item \textbf{No se deduce.} En el enfoque frecuentista la mejora real es fija; el $95\,\%$ es la cobertura del procedimiento.
  \item \textbf{No se deduce}, por la misma razón: el intervalo no asigna probabilidades a hipótesis.
  \item \textbf{Se deduce}: el centro, $15$, es la estimación puntual y el intervalo indica su precisión.
  \item \textbf{No se deduce}: «probablemente falsa» es una afirmación de probabilidad sobre una hipótesis.
  \item \textbf{Se deduce}: es la definición de cobertura (\cref{def:t4-ic}).
  \item \textbf{Se deduce}: el valor $0$ queda fuera del intervalo, y un intervalo de nivel $95\,\%$ contiene exactamente los valores no rechazados por el contraste bilateral al $5\,\%$ (\cref{eje:t5-dualidad}).
\end{enumerate}
\resultado \textbf{Se deducen (iii), (v) y (vi); no se deducen (i), (ii) y (iv).}
\end{solucion}

\begin{solucion}{eje:t4-momentos-basico}
\planteamiento \cref{ej:t4-momentos}: igualar media (y segundo momento) muestrales y poblacionales.
\pasos
\begin{enumerate}[label=\arabic*.]
  \item (a) $\xbar=(0{,}5+1{,}2+0{,}3+2{,}0)/4=4{,}0/4=1{,}0$; $\widehat\lambda=1/\xbar=1{,}0$ peticiones por segundo.
  \item (b) $\widehat\mu=\xbar=22/5=4{,}4$. Desviaciones: $-2{,}4;\allowbreak\,-0{,}4;\allowbreak\,-0{,}4;\allowbreak\,0{,}6;\allowbreak\,2{,}6$; cuadrados: $5{,}76;\allowbreak\,0{,}16;\allowbreak\,0{,}16;\allowbreak\,0{,}36;\allowbreak\,6{,}76$; suma $13{,}2$.
  \item $\widehat\sigma^{2}=13{,}2/5=2{,}64$, mientras que $s^{2}=13{,}2/4=3{,}3$.
\end{enumerate}
\resultado \textbf{(a) $\widehat\lambda=1$; (b) $\widehat\mu=4{,}4$ y $\widehat\sigma^{2}=2{,}64$, menor que $s^{2}=3{,}3$.}
\comprobacion $\frac1n\sum x_i^{2}-\xbar^{2}=(4+16+16+25+49)/5-19{,}36=22-19{,}36=2{,}64$.
\end{solucion}

\begin{solucion}{eje:t4-geometrica}
\datos $f(x\mid p)=p(1-p)^{x-1}$, $\E X=1/p$; datos $2,1,4,3,1,1$.
\planteamiento (a) $1/p=\xbar$; (b) derivar la log-verosimilitud.
\pasos
\begin{enumerate}[label=\arabic*.]
  \item Momentos: $1/\widehat p=\xbar$, luego $\widehat p=1/\xbar$.
  \item Verosimilitud: $\Lik(p)=\prod_ip(1-p)^{x_i-1}=p^{n}(1-p)^{\sum x_i-n}$; $\ell(p)=n\ln p+\bigl(\sum x_i-n\bigr)\ln(1-p)$.
  \item $\ell'(p)=\dfrac np-\dfrac{\sum x_i-n}{1-p}=0\iff n(1-p)=\bigl(\sum x_i-n\bigr)p\iff n=p\sum x_i\iff p=\dfrac{n}{\sum x_i}=\dfrac1{\xbar}$.
  \item $\ell''(p)=-n/p^{2}-\bigl(\sum x_i-n\bigr)/(1-p)^{2}<0$: máximo global.
  \item Datos: $\sum x_i=12$, $n=6$, $\xbar=2$, $\widehat p=0{,}5$.
\end{enumerate}
\resultado \textbf{Ambos métodos dan $\widehat p=1/\xbar$; con los datos, $\widehat p=0{,}5$.}
\comprobacion un agente que resuelve cada intento con probabilidad $0{,}5$ necesita en media $2$ intentos, que es la media observada.
\end{solucion}

\begin{solucion}{eje:t4-gamma}
\datos $f(x\mid a)=a^{5}x^{4}e^{-ax}/24$; datos $2{,}0;\allowbreak\,3{,}0;\allowbreak\,2{,}5;\allowbreak\,4{,}5$.
\planteamiento escribir $\ell(a)$ con $S=\sum x_i$ y $P=\prod x_i$ y resolver $\ell'(a)=0$.
\pasos
\begin{enumerate}[label=\arabic*.]
  \item $\Lik(a)=\dfrac{a^{5n}P^{4}e^{-aS}}{24^{n}}$; $\ell(a)=5n\ln a+4\ln P-aS-n\ln24$.
  \item $\ell'(a)=5n/a-S=0\iff a=5n/S$.
  \item $\ell''(a)=-5n/a^{2}<0$: máximo global en $(0,\infty)$.
  \item Datos: $n=4$, $S=12$; $\widehat a=20/12=5/3\approx1{,}667$.
\end{enumerate}
\resultado \textbf{$\widehat a=5n/\sum x_i=5/\xbar$; con los datos, $\widehat a=5/3$.}
\comprobacion la media de esta distribución es $5/a$ (es una gamma de forma $5$), luego el EMV coincide con el estimador por momentos: $5/\widehat a=\xbar=3$.
\end{solucion}

\begin{solucion}{eje:t4-uniforme-ab}
\datos $1{,}2;\allowbreak\,2{,}1;\allowbreak\,1{,}3;\allowbreak\,10{,}5;\allowbreak\,5$; $n=5$.
\planteamiento (a) $\Lik(a,b)=(b-a)^{-n}$ si $a\le\min x_i$ y $b\ge\max x_i$, y $0$ en otro caso; (b) igualar media y varianza (con denominador $n$).
\pasos
\begin{enumerate}[label=\arabic*.]
  \item (a) $\Lik$ es positiva solo si el intervalo $[a,b]$ contiene todos los datos y, en ese caso, es mayor cuanto menor es $b-a$. El intervalo más corto que contiene los datos es $[\min x_i,\max x_i]$: $\widehat a=1{,}2$ y $\widehat b=10{,}5$.
  \item (b) $\xbar=20{,}1/5=4{,}02$; $\frac1n\sum x_i^{2}=(1{,}44+4{,}41+1{,}69+110{,}25+25)/5=142{,}79/5=28{,}558$; $\widehat\sigma^{2}=28{,}558-4{,}02^{2}=12{,}3976$; $\widehat\sigma\approx3{,}521$.
  \item De $(b-a)/2=\sqrt3\,\widehat\sigma\approx6{,}099$ y $(a+b)/2=4{,}02$: $\widehat a\approx-2{,}08$ y $\widehat b\approx10{,}12$.
  \item (c) $\widehat b\approx10{,}12<10{,}5$: según el modelo estimado, el dato $10{,}5$ sería imposible.
\end{enumerate}
\resultado \textbf{EMV: $\widehat a=1{,}2$, $\widehat b=10{,}5$. Momentos: $\widehat a\approx-2{,}08$, $\widehat b\approx10{,}12$, incompatibles con los datos.}
\comprobacion la media del modelo con los estimadores de momentos es $(-2{,}08+10{,}12)/2=4{,}02=\xbar$, como debe ser.
\end{solucion}

\begin{solucion}{eje:t4-poisson-datos}
\datos $2,0,3,1,4,2,1,3$; $n=8$.
\planteamiento $\widehat\lambda=\xbar$ (\cref{eje:t4-emv-poisson}); intervalo $\widehat\lambda\pm1{,}96\sqrt{\widehat\lambda/n}$.
\pasos
\begin{enumerate}[label=\arabic*.]
  \item $\sum x_i=16$; $\widehat\lambda=2$.
  \item $\widehat{\ee}=\sqrt{2/8}=0{,}5$; margen $1{,}96\cdot0{,}5=0{,}98$.
\end{enumerate}
\resultado \textbf{$\widehat\lambda=2$; intervalo aproximado $[1{,}02;\allowbreak\ 2{,}98]$.}
\comprobacion con $n=8$ la aproximación normal es tosca; el intervalo es solo orientativo, pero su amplitud ya muestra que $\lambda$ está poco determinado.
\end{solucion}

\begin{solucion}{eje:t4-soda}
\datos $n=9$; $\sum x_i=3204$; $\xbar=356$.
\planteamiento (a) \cref{teo:t4-ic-z}; (b) \cref{teo:t4-ic-t} con $8$ grados de libertad.
\pasos
\begin{enumerate}[label=\arabic*.]
  \item $\sigma/\sqrt n=3/3=1$. Al $95\,\%$: $356\pm1{,}96$, es decir, $[354{,}04;\allowbreak\ 357{,}96]$. Al $98\,\%$ ($z_{0{,}99}=2{,}326$): $[353{,}67;\allowbreak\ 358{,}33]$.
  \item Desviaciones respecto de $356$: $-4,-5,5,-3,-4,2,4,2,3$; suma de cuadrados $16+25+25+9+16+4+16+4+9=124$; $s^{2}=124/8=15{,}5$; $s\approx3{,}937$; $s/\sqrt n\approx1{,}312$.
  \item Al $95\,\%$: $t_{8;0{,}975}=2{,}306$, margen $3{,}026$, intervalo $[352{,}97;\allowbreak\ 359{,}03]$.
  \item Al $98\,\%$: $t_{8;0{,}99}=2{,}896$, margen $3{,}801$, intervalo $[352{,}20;\allowbreak\ 359{,}80]$.
\end{enumerate}
\resultado \textbf{$\sigma$ conocida: $[354{,}04;\allowbreak\,357{,}96]$ y $[353{,}67;\allowbreak\,358{,}33]$. $\sigma$ desconocida: $[352{,}97;\allowbreak\,359{,}03]$ y $[352{,}20;\allowbreak\,359{,}80]$.}
\comprobacion los intervalos $t$ son más anchos por dos motivos: $t_{8}>z$ y, en esta muestra, $s\approx3{,}94>\sigma=3$. Al aumentar el nivel, todos los intervalos se ensanchan.
\end{solucion}

\begin{solucion}{eje:t4-alturas}
\datos $n=20$; $\xbar=69{,}55$; $s^{2}=14{,}26$; en (b), $\sigma=3{,}77$.
\planteamiento \cref{teo:t4-ic-t,teo:t4-ic-z} y \cref{cor:t4-tamano} con $E=0{,}5$ (la mitad de la amplitud).
\pasos
\begin{enumerate}[label=\arabic*.]
  \item (a) $s/\sqrt n=\sqrt{14{,}26/20}=\sqrt{0{,}713}\approx0{,}8444$; $t_{19;0{,}95}=1{,}729$; margen $1{,}460$; intervalo $[68{,}09;\allowbreak\ 71{,}01]$.
  \item (b) $\sigma/\sqrt n=3{,}77/\sqrt{20}\approx0{,}8430$; $z_{0{,}95}=1{,}645$; margen $1{,}387$; intervalo $[68{,}16;\allowbreak\ 70{,}94]$.
  \item (c) $n\ge(1{,}645\cdot3{,}77/0{,}5)^{2}=(12{,}40)^{2}\approx153{,}8$, luego $n=154$.
\end{enumerate}
\resultado \textbf{(a) $[68{,}09;\allowbreak\,71{,}01]$; (b) $[68{,}16;\allowbreak\,70{,}94]$; (c) $154$ personas.}
\comprobacion con $n=154$ la amplitud es $2\cdot1{,}645\cdot3{,}77/\sqrt{154}\approx0{,}9995<1$.
\end{solucion}

\begin{solucion}{eje:t4-openstax-z}
\datos (a) $\sigma=3$, $n=36$, $\xbar=68$, nivel $90\,\%$; (b) $\sigma=15$, $E=2$, nivel $95\,\%$.
\planteamiento \cref{teo:t4-ic-z} y \cref{cor:t4-tamano}.
\pasos
\begin{enumerate}[label=\arabic*.]
  \item (a) Margen $1{,}645\cdot3/6=0{,}8225$; intervalo $[67{,}18;\allowbreak\ 68{,}82]$.
  \item (b) $n\ge(1{,}96\cdot15/2)^{2}=14{,}7^{2}=216{,}09$, luego $n=217$.
\end{enumerate}
\resultado \textbf{(a) $[67{,}18;\allowbreak\,68{,}82]$; (b) $n=217$.}
\comprobacion se redondea hacia arriba: con $n=216$ el margen sería $1{,}96\cdot15/\sqrt{216}\approx2{,}0004>2$.
\end{solucion}

\begin{solucion}{eje:t4-proporcion}
\datos (a) $n=500$, $x=421$; (b) $E=0{,}03$, nivel $90\,\%$.
\planteamiento \cref{def:t4-wald} y \cref{cor:t4-tamano-p}.
\pasos
\begin{enumerate}[label=\arabic*.]
  \item (a) $\widehat p=421/500=0{,}842$; $\widehat{\ee}=\sqrt{0{,}842\cdot0{,}158/500}\approx0{,}01631$; margen $1{,}96\cdot0{,}01631\approx0{,}0320$; intervalo $[0{,}810;\allowbreak\ 0{,}874]$.
  \item (b) $n\ge z_{0{,}95}^{2}/(4E^{2})=1{,}645^{2}/(4\cdot0{,}0009)=2{,}706/0{,}0036\approx751{,}7$, luego $n=752$.
\end{enumerate}
\resultado \textbf{(a) $[0{,}810;\allowbreak\,0{,}874]$; (b) $752$ personas.}
\comprobacion en (a), $n\widehat p=421$ y $n(1-\widehat p)=79$ son grandes, y la aproximación normal es adecuada.
\end{solucion}

\begin{solucion}{eje:t4-margen}
\planteamiento margen conservador $z/(2\sqrt n)=E\iff n=\bigl(z/(2E)\bigr)^{2}$.
\pasos
\begin{enumerate}[label=\arabic*.]
  \item (a) $z=1{,}96$: $n=(1{,}96/0{,}02)^{2}=98^{2}=9604$. Con $z\approx2$: $n=100^{2}=10\,000$.
  \item (b) $z_{0{,}90}=1{,}2816$: $n=(1{,}2816/0{,}02)^{2}=64{,}08^{2}\approx4106$.
  \item (c) Con $z\approx2$: $2/(2\sqrt n)=0{,}04\iff\sqrt n=25\iff n=625$ (con $1{,}96$ saldría unas $600$).
\end{enumerate}
\resultado \textbf{(a) $9604$ (o $10\,000$ con $z\approx2$); (b) unas $4106$; (c) unas $600$ personas.}
\comprobacion la regla práctica «margen $\approx1/\sqrt n$ al $95\,\%$» da $1/\sqrt{625}=0{,}04$: un benchmark de $10\,000$ preguntas tiene un margen de un punto porcentual.
\end{solucion}

\begin{solucion}{eje:t4-emv-poisson}
\planteamiento escribir $\ell(\lambda)$, resolver $\ell'(\lambda)=0$ y estudiar la concavidad.
\pasos
\begin{enumerate}[label=\arabic*.]
  \item $f(x\mid\lambda)=e^{-\lambda}\lambda^{x}/x!$, luego $\ell(\lambda)=\Bigl(\sum_ix_i\Bigr)\ln\lambda-n\lambda-\sum_i\ln(x_i!)$ para $\lambda>0$.
  \item $\ell'(\lambda)=\dfrac{\sum_ix_i}{\lambda}-n=0\iff\lambda=\dfrac{\sum_ix_i}{n}=\xbar$.
  \item $\ell''(\lambda)=-\sum_ix_i/\lambda^{2}<0$ para todo $\lambda>0$ (pues $\sum x_i>0$): $\ell$ es estrictamente cóncava en $(0,\infty)$ y su único punto crítico es el máximo global.
  \item Por el \cref{ej:t4-momentos}, el estimador por momentos también es $\Xbar$.
\end{enumerate}
\resultado \textbf{El EMV es $\widehat\lambda=\xbar$ y coincide con el de momentos.} $\blacksquare$
\comprobacion con los datos del \cref{eje:t4-poisson-datos}, $\ell'(2)=16/2-8=0$.
\end{solucion}

\begin{solucion}{eje:t4-emv-expo}
\planteamiento derivar la log-verosimilitud en cada parametrización.
\pasos
\begin{enumerate}[label=\arabic*.]
  \item (a) $\ell(\lambda)=n\ln\lambda-\lambda\sum_ix_i$; $\ell'(\lambda)=n/\lambda-\sum_ix_i=0\iff\lambda=n/\sum_ix_i=1/\xbar$; $\ell''(\lambda)=-n/\lambda^{2}<0$: máximo global.
  \item (b) Con $\beta=1/\lambda$, $f(x\mid\beta)=\beta^{-1}e^{-x/\beta}$ y $\ell(\beta)=-n\ln\beta-\sum_ix_i/\beta$.
  \item $\ell'(\beta)=-\dfrac n\beta+\dfrac{\sum_ix_i}{\beta^{2}}=\dfrac{\sum_ix_i-n\beta}{\beta^{2}}$, positiva para $\beta<\xbar$ y negativa para $\beta>\xbar$: máximo global en $\widehat\beta=\xbar$.
  \item $\widehat\beta=\xbar=1/\widehat\lambda$, como predice la invariancia con $g(\lambda)=1/\lambda$.
\end{enumerate}
\resultado \textbf{$\widehat\lambda=1/\xbar$ y $\widehat\beta=\xbar$.} $\blacksquare$
\comprobacion con los datos del \cref{eje:t4-momentos-basico}(a), $\widehat\lambda=1$ y $\widehat\beta=1$ segundo.
\end{solucion}

\begin{solucion}{eje:t4-laplace}
\planteamiento la log-verosimilitud es, salvo una constante, $-\sum_i\abs{x_i-\mu}$.
\pasos
\begin{enumerate}[label=\arabic*.]
  \item $\ell(\mu)=\sum_i\ln\bigl(\tfrac12e^{-\abs{x_i-\mu}}\bigr)=-n\ln2-\sum_i\abs{x_i-\mu}$.
  \item Maximizar $\ell$ equivale a minimizar $g(\mu)=\sum_i\abs{x_i-\mu}$, que alcanza su mínimo en la mediana (\cref{eje:t1-mediana-min}).
  \item En el modelo normal con $\sigma$ fija, $\ell(\mu)=c-\frac{1}{2\sigma^{2}}\sum_i(x_i-\mu)^{2}$, que se maximiza en $\xbar$ (\cref{prop:t1-extremal}).
\end{enumerate}
\resultado \textbf{El EMV de la posición de Laplace es la mediana muestral} (cualquier punto entre los dos datos centrales si $n$ es par). $\blacksquare$
\comprobacion la densidad de Laplace tiene colas más pesadas que la normal, y su EMV es el estimador robusto frente a atípicos: el modelo elegido determina el estimador.
\end{solucion}

\begin{solucion}{eje:t4-regla3}
\planteamiento $\PP_p(\text{cero fallos})=(1-p)^{n}$, decreciente en $p$; usar $e^{x}\ge1+x$.
\pasos
\begin{enumerate}[label=\arabic*.]
  \item (a) $(1-p)^{n}\ge\alpha\iff1-p\ge\alpha^{1/n}\iff p\le1-\alpha^{1/n}$, pues $t\mapsto t^{n}$ es creciente en $[0,1]$. El conjunto es $[0,p_{\max}]$ con $p_{\max}=1-\alpha^{1/n}$.
  \item (b) $\alpha^{1/n}=e^{(\ln\alpha)/n}\ge1+\dfrac{\ln\alpha}{n}$, luego $p_{\max}=1-\alpha^{1/n}\le-\dfrac{\ln\alpha}{n}$. Para $\alpha=0{,}05$, $-\ln0{,}05\approx2{,}996\approx3$.
  \item (c) $n=300$: $p_{\max}=1-0{,}05^{1/300}\approx0{,}009\,936$; $3/n=0{,}01$.
\end{enumerate}
\resultado \textbf{$p_{\max}=1-\alpha^{1/n}\le-\ln(\alpha)/n\approx3/n$; con $n=300$, $p_{\max}\approx0{,}0099$}: con $95\,\%$ de confianza, la tasa de fallos graves es menor que un $1\,\%$. $\blacksquare$
\comprobacion $(1-0{,}009\,936)^{300}=e^{300\ln(0{,}990\,064)}\approx e^{-2{,}996}\approx0{,}05$.
\end{solucion}

\begin{solucion}{eje:t4-conservador}
\planteamiento calcular la cobertura aproximada con la normal y usar $\sqrt{p(1-p)}\le\frac12$.
\pasos
\begin{enumerate}[label=\arabic*.]
  \item Sea $c=z_{1-\alpha/2}$ y $\tau=\sqrt{p(1-p)/n}$. El intervalo cubre $p$ si y solo si $\abs{\widehat p-p}\le c/(2\sqrt n)$.
  \item Con $\widehat p\approx\Nor(p,\tau^{2})$: $\PP\bigl(\abs{\widehat p-p}\le\tfrac{c}{2\sqrt n}\bigr)\approx2\Phi\Bigl(\dfrac{c}{2\sqrt n\,\tau}\Bigr)-1=2\Phi\Bigl(\dfrac{c}{2\sqrt{p(1-p)}}\Bigr)-1$.
  \item Como $p(1-p)\le1/4$ con igualdad solo si $p=1/2$, el argumento es $\ge c$, y como $\Phi$ es creciente, la cobertura aproximada es $\ge2\Phi(c)-1=1-\alpha$, con igualdad solo si $p=1/2$.
  \item $p=0{,}1$: $\sqrt{p(1-p)}=0{,}3$ y el argumento vale $1{,}96/0{,}6\approx3{,}27$; cobertura aproximada $2\Phi(3{,}27)-1\approx0{,}9989$, independiente de $n$.
\end{enumerate}
\resultado \textbf{La cobertura aproximada es al menos $1-\alpha$}, y el intervalo es tanto más conservador cuanto más lejos está $p$ de $1/2$; con $p=0{,}1$ cubre el $99{,}9\,\%$ de las veces en lugar del $95\,\%$. $\blacksquare$
\comprobacion el precio es la amplitud: con $p=0{,}1$ el intervalo de Wald tendría un margen de $1{,}96\cdot0{,}3/\sqrt n$, un $40\,\%$ menor que el conservador $1{,}96\cdot0{,}5/\sqrt n$.
\end{solucion}

% ---------------------------------------------------------------------
\section{Tema 5}

\begin{solucion}{eje:t5-vf}
\planteamiento \cref{def:t5-pvalor}, \cref{prop:t5-pvalor} y \cref{obs:t5-pvalor}.
\pasos
\begin{enumerate}[label=(\alph*)]
  \item \textbf{Falsa.} El $p$-valor se calcula suponiendo $H_0$ cierta; no es $\PP(H_0\mid\text{datos})$.
  \item \textbf{Verdadera}: $0{,}03\le0{,}05$.
  \item \textbf{Falsa}: $0{,}03>0{,}01$.
  \item \textbf{Verdadera}: es la definición.
  \item \textbf{Falsa}: el $p$-valor no mide el tamaño del efecto; depende también de $n$.
  \item \textbf{Falsa}: no rechazar no demuestra $H_0$.
\end{enumerate}
\resultado \textbf{(a) F; (b) V; (c) F; (d) V; (e) F; (f) F.}
\end{solucion}

\begin{solucion}{eje:t5-moneda}
\datos $n=250$, $x=140$; bajo $H_0$, $X\sim\Bin(250;\allowbreak\,0{,}5)$, con media $125$ y desviación típica $\sqrt{62{,}5}\approx7{,}906$.
\planteamiento aproximación normal con corrección por continuidad (\cref{def:t3-cc}); contraste bilateral.
\pasos
\begin{enumerate}[label=\arabic*.]
  \item (a) $p=2\PP(X\ge140)\approx2\Bigl(1-\Phi\bigl(\tfrac{139{,}5-125}{7{,}906}\bigr)\Bigr)=2\bigl(1-\Phi(1{,}83)\bigr)=2\cdot0{,}0336=0{,}0672$.
  \item (b) $0{,}0672\le0{,}10$: se rechaza al $10\,\%$; $0{,}0672>0{,}05$: no se rechaza al $5\,\%$.
  \item (c) Al $1\,\%$ se rechaza si $\abs{x-125}-0{,}5\ge z_{0{,}995}\cdot7{,}906=2{,}576\cdot7{,}906\approx20{,}37$, es decir, $\abs{x-125}\ge20{,}87$: hacen falta al menos $146$ caras (o a lo sumo $104$).
  \item (d) Al $5\,\%$ se rechaza si $\abs{x-125}-0{,}5\ge1{,}96\cdot7{,}906\approx15{,}50$, es decir, $x\ge141$ o $x\le109$. Si $\theta=0{,}6$: media $150$, desviación típica $\sqrt{60}\approx7{,}746$, y $\PP(X\ge141)\approx1-\Phi\bigl(\tfrac{140{,}5-150}{7{,}746}\bigr)=\Phi(1{,}23)=0{,}8907$; $\PP(X\le109)$ es despreciable.
\end{enumerate}
\resultado \textbf{(a) $p\approx0{,}067$; (b) se rechaza al $10\,\%$, no al $5\,\%$; (c) $146$ caras; (d) potencia $\approx0{,}89$.}
\comprobacion los valores binomiales exactos son $p=0{,}0664$ (el «$7\,\%$» de la fuente), $2\PP(X\ge146)=0{,}0094\le0{,}01$ y potencia $0{,}8896$.
\end{solucion}

\begin{solucion}{eje:t5-poligrafo}
\datos sinceras: $131$ bien, $9$ mal; mentirosas: $15$ mal, $125$ bien.
\planteamiento con $H_0$: «dice la verdad», rechazar $H_0$ es clasificar como mentirosa.
\pasos
\begin{enumerate}[label=\arabic*.]
  \item Error de tipo~I: clasificar como mentirosa a una persona sincera. Estimación: $9/140\approx0{,}064$, que es el nivel estimado del detector.
  \item Error de tipo~II: clasificar como sincera a una persona que miente. Estimación: $15/140\approx0{,}107$.
  \item Potencia: $1-0{,}107=125/140\approx0{,}893$.
\end{enumerate}
\resultado \textbf{$\PP(\text{tipo I})\approx0{,}064$; $\PP(\text{tipo II})\approx0{,}107$; nivel $\approx0{,}064$; potencia $\approx0{,}893$.}
\comprobacion nivel $=\PP(\text{tipo I})$ y potencia $=1-\PP(\text{tipo II})$, como en la \cref{def:t5-errores}.
\end{solucion}

\begin{solucion}{eje:t5-radar}
\datos $n=3$; $\sigma=5$; $\alpha=0{,}04$; $\mu_0=40$.
\planteamiento contraste $z$ unilateral (\cref{prop:t5-potencia}); $\sigma/\sqrt3\approx2{,}887$; $z_{0{,}96}=1{,}7507$.
\pasos
\begin{enumerate}[label=\arabic*.]
  \item (a) $H_0:\mu=40$ (el coche no supera el límite; también puede tomarse la compuesta $\mu\le40$) frente a $H_1:\mu>40$. Se multa si $\xbar>40+1{,}7507\cdot2{,}887\approx45{,}05$.
  \item (b) Si nadie supera el límite, toda multa es un error de tipo~I: el $100\,\%$ de las multas serían erróneas (aunque solo se multaría a un $4\,\%$ o menos de los conductores).
  \item (c) $\pi(45)=1-\Phi\bigl(1{,}7507-5\sqrt3/5\bigr)=1-\Phi(1{,}7507-1{,}7321)=1-\Phi(0{,}019)\approx0{,}49$.
  \item (d) $n\ge\bigl((z_{0{,}96}+z_{0{,}90})\sigma/\delta\bigr)^{2}=\bigl((1{,}7507+1{,}2816)\cdot5/5\bigr)^{2}\approx9{,}19$, luego $10$ radares.
\end{enumerate}
\resultado \textbf{(a) umbral $\approx45{,}05$ km/h; (b) todas; (c) potencia $\approx0{,}49$; (d) $10$ radares.}
\comprobacion la potencia contra $45$ es casi $1/2$ porque el umbral ($45{,}05$) es casi igual a $45$: un coche a $45$ supera el umbral la mitad de las veces.
\end{solucion}

\begin{solucion}{eje:t5-uniforme}
\datos $X\sim\Unif(0,\theta)$; región de rechazo $\{x\le0{,}1\}\cup\{x\ge1{,}9\}$.
\planteamiento calcular probabilidades de intervalos con la densidad uniforme.
\pasos
\begin{enumerate}[label=\arabic*.]
  \item (a) Con $\theta=2$: $\PP(X\le0{,}1)+\PP(X\ge1{,}9)=0{,}1/2+0{,}1/2=0{,}1$.
  \item (b) Con $\theta=2{,}5$, no se rechaza si $0{,}1<X<1{,}9$: probabilidad $1{,}8/2{,}5=0{,}72$.
\end{enumerate}
\resultado \textbf{(a) $0{,}10$; (b) $0{,}72$.}
\comprobacion la potencia contra $\theta=2{,}5$ es $0{,}28$, mayor que el nivel $0{,}10$, pero pequeña: con una sola observación es difícil distinguir $\theta=2$ de $\theta=2{,}5$.
\end{solucion}

\begin{solucion}{eje:t5-t-test}
\datos $n=16$; $\xbar=11$; $s=2$; $\mu_0=10$; $\alpha=0{,}05$; $15$ grados de libertad.
\planteamiento contraste $t$; $t_{15;0{,}975}=2{,}131$ y $t_{15;0{,}95}=1{,}753$.
\pasos
\begin{enumerate}[label=\arabic*.]
  \item $t=\dfrac{11-10}{2/\sqrt{16}}=\dfrac{1}{0{,}5}=2$.
  \item (a) Bilateral: $\abs t=2<2{,}131$, no se rechaza. Como $1{,}753<2<2{,}131$, el $p$-valor unilateral está entre $0{,}025$ y $0{,}05$ y el bilateral entre $0{,}05$ y $0{,}10$.
  \item (b) Unilateral: $t=2>1{,}753$, se rechaza; $0{,}025<p<0{,}05$.
\end{enumerate}
\resultado \textbf{(a) No se rechaza ($0{,}05<p<0{,}10$); (b) se rechaza ($0{,}025<p<0{,}05$).}
\comprobacion los $p$-valores exactos son $0{,}0639$ y $0{,}0320$. El ejemplo muestra por qué la alternativa debe fijarse antes de ver los datos (\cref{obs:t5-errores}(a)).
\end{solucion}

\begin{solucion}{eje:t5-openstax-z}
\datos $\mu_0=16{,}43$; $\sigma=0{,}8$; $n=15$; $\xbar=16$; $\alpha=0{,}05$.
\planteamiento contraste $z$ unilateral: $H_0:\mu=16{,}43$ frente a $H_1:\mu<16{,}43$.
\pasos
\begin{enumerate}[label=\arabic*.]
  \item $\sigma/\sqrt n=0{,}8/\sqrt{15}\approx0{,}2066$.
  \item $z=(16-16{,}43)/0{,}2066\approx-2{,}08$.
  \item $p=\Phi(-2{,}08)=1-0{,}9812=0{,}0188<0{,}05$.
\end{enumerate}
\resultado \textbf{$z\approx-2{,}08$, $p\approx0{,}019$: se rechaza $H_0$}; hay evidencia de que la media es menor que la anunciada.
\comprobacion equivalentemente, $z=-2{,}08<-z_{0{,}95}=-1{,}645$.
\end{solucion}

\begin{solucion}{eje:t5-openstax-prop}
\datos $p_0=0{,}25$; $n=42$; $x=17$; $\alpha=0{,}01$.
\planteamiento contraste $z$ unilateral para una proporción, con el error estándar calculado con $p_0$.
\pasos
\begin{enumerate}[label=\arabic*.]
  \item $\widehat p=17/42\approx0{,}4048$.
  \item $\sqrt{p_0(1-p_0)/n}=\sqrt{0{,}1875/42}\approx0{,}06682$.
  \item $z=(0{,}4048-0{,}25)/0{,}06682\approx2{,}32$; $p=1-\Phi(2{,}32)=1-0{,}9898=0{,}0102$.
  \item $0{,}0102>0{,}01$: no se rechaza.
\end{enumerate}
\resultado \textbf{$z\approx2{,}32$, $p\approx0{,}010$: no se rechaza $H_0$ al $1\,\%$} (sí se rechazaría al $5\,\%$).
\comprobacion $np_0=10{,}5$ y $n(1-p_0)=31{,}5$, de modo que la aproximación normal es aceptable; el resultado está justo en el límite, y conviene informar del $p$-valor en lugar de solo «no significativo».
\end{solucion}

\begin{solucion}{eje:t5-chi-dias}
\datos $O=(15,12,9,9,15)$; $n=60$; $k=5$.
\planteamiento \cref{def:t5-chi2} con $p_{0j}=1/5$.
\pasos
\begin{enumerate}[label=\arabic*.]
  \item $E_j=60/5=12$ para todo $j$.
  \item $X^{2}=\dfrac{3^{2}+0^{2}+3^{2}+3^{2}+3^{2}}{12}=\dfrac{36}{12}=3$.
  \item Grados de libertad: $k-1=4$; $\chi^{2}_{4;0{,}95}=9{,}488$.
  \item $3<9{,}488$: no se rechaza.
\end{enumerate}
\resultado \textbf{$X^{2}=3$ con $4$ grados de libertad; no se rechaza la uniformidad ($p\approx0{,}56$).}
\comprobacion $\E[\chi^{2}_4]=4$: un valor de $3$ es completamente típico bajo $H_0$.
\end{solucion}

\begin{solucion}{eje:t5-benford}
\datos $n=100$; $E_j=100p_j$; $O=(7,13,12,9,9,13,11,10,16)$.
\planteamiento \cref{def:t5-chi2} con $k-1=8$ grados de libertad.
\pasos
\begin{enumerate}[label=\arabic*.]
  \item $E=(30{,}1;\allowbreak\,17{,}6;\allowbreak\,12{,}5;\allowbreak\,9{,}7;\allowbreak\,7{,}9;\allowbreak\,6{,}7;\allowbreak\,5{,}8;\allowbreak\,5{,}1;\allowbreak\,4{,}6)$.
  \item Contribuciones $(O_j-E_j)^{2}/E_j$: $17{,}73$; $1{,}20$; $0{,}02$; $0{,}05$; $0{,}15$; $5{,}92$; $4{,}66$; $4{,}71$; $28{,}25$.
  \item $X^{2}\approx62{,}70$; $\chi^{2}_{8;0{,}999}=26{,}12$.
  \item $62{,}70>26{,}12$: se rechaza $H_0$ al nivel $0{,}001$.
\end{enumerate}
\resultado \textbf{$X^{2}\approx62{,}7$ con $8$ grados de libertad: se rechaza con mucha claridad} (hay demasiados nueves y muy pocos unos); se recomienda revisar la nota de gastos.
\comprobacion $E_9=4{,}6<5$. Agrupando los dígitos $8$ y $9$ ($E=9{,}7$, $O=26$) se obtiene $X^{2}\approx57{,}1$ con $7$ grados de libertad, frente a $\chi^{2}_{7;0{,}999}=24{,}32$: la conclusión no cambia.
\end{solucion}

\begin{solucion}{eje:t5-dos-modelos}
\datos $n_A=n_B=1000$; $\widehat p_A=0{,}78$; $\widehat p_B=0{,}75$.
\planteamiento bajo $H_0$, $p_A=p_B=p$, $\Var(\widehat p_A-\widehat p_B)=p(1-p)(1/n_A+1/n_B)$, y $p$ se estima con la proporción combinada.
\pasos
\begin{enumerate}[label=\arabic*.]
  \item (a) $H_0:p_A=p_B$ (compuesta: no fija el valor común) frente a $H_1:p_A\ne p_B$.
  \item (b) $\widehat p=(780+750)/2000=0{,}765$; $\widehat{\Var}=0{,}765\cdot0{,}235\cdot(2/1000)\approx0{,}000\,359\,6$; $\widehat{\ee}\approx0{,}01896$.
  \item $z=(0{,}78-0{,}75)/0{,}01896\approx1{,}58$; $p=2\bigl(1-\Phi(1{,}58)\bigr)=2\cdot0{,}0571=0{,}1142$.
  \item (c) $0{,}114>0{,}05$: no es significativa.
\end{enumerate}
\resultado \textbf{$z\approx1{,}58$, $p\approx0{,}11$: la diferencia de tres puntos no es significativa al $5\,\%$.}
\comprobacion para que una diferencia de tres puntos fuera significativa con estas exactitudes haría falta que $0{,}03/\sqrt{0{,}18\cdot2/n}\ge1{,}96$, es decir, $n\ge1{,}96^{2}\cdot0{,}36/0{,}0009\approx1537$ preguntas por modelo.
\end{solucion}

\begin{solucion}{eje:t5-multiples}
\datos $m=20$; $\alpha=0{,}05$; todas las nulas ciertas; contrastes independientes.
\planteamiento \cref{prop:t5-multiples} y \cref{teo:t5-bonferroni}.
\pasos
\begin{enumerate}[label=\arabic*.]
  \item (a) $\PP(\text{al menos uno})=1-0{,}95^{20}\approx1-0{,}3585=0{,}6415$. Número esperado: $20\cdot0{,}05=1$.
  \item (b) Bonferroni: nivel $0{,}05/20=0{,}0025$ por contraste.
  \item Probabilidad exacta de algún falso positivo con contrastes independientes: $1-(1-0{,}0025)^{20}=1-0{,}9975^{20}\approx0{,}0488$.
\end{enumerate}
\resultado \textbf{(a) $\approx0{,}64$ y $1$ falso positivo esperado; (b) $0{,}0025$ por contraste y FWER $\approx0{,}049\le0{,}05$.}
\comprobacion el valor $0{,}0488$ es algo menor que $0{,}05$, como garantiza el \cref{teo:t5-bonferroni}; la diferencia es la holgura de la desigualdad de la unión.
\end{solucion}

\begin{solucion}{eje:t5-sidak}
\planteamiento (a) complemento e independencia; (b) igualar a $\alpha$; (c) desigualdad de Bernoulli.
\pasos
\begin{enumerate}[label=\arabic*.]
  \item (a) $\PP(\min_iP_i>t)=\PP(P_1>t,\dots,P_n>t)=\prod_i\PP(P_i>t)=(1-t)^{n}$, luego $\PP(\min_iP_i\le t)=1-(1-t)^{n}$.
  \item (b) Se produce algún falso positivo si y solo si $\min_iP_i\le t$. Con $t=1-(1-\alpha)^{1/n}$: $1-(1-t)^{n}=1-\bigl((1-\alpha)^{1/n}\bigr)^{n}=1-(1-\alpha)=\alpha$.
  \item (c) Con $x=-\alpha/n\ge-1$, Bernoulli da $(1-\alpha/n)^{n}\ge1-\alpha$. Tomando raíces $n$-ésimas (función creciente): $1-\alpha/n\ge(1-\alpha)^{1/n}$, es decir, $\alpha/n\le1-(1-\alpha)^{1/n}$.
  \item Un umbral mayor rechaza más hipótesis, luego el procedimiento es al menos tan potente como el de Bonferroni.
\end{enumerate}
\resultado \textbf{$\PP(\min P_i\le t)=1-(1-t)^{n}$; el umbral $1-(1-\alpha)^{1/n}$ da FWER exactamente $\alpha$ y es mayor o igual que $\alpha/n$.} $\blacksquare$
\comprobacion con $n=20$ y $\alpha=0{,}05$: $1-0{,}95^{1/20}\approx0{,}002\,561>0{,}0025$. La ganancia es pequeña, y a cambio exige independencia; Bonferroni no la exige.
\end{solucion}

\begin{solucion}{eje:t5-uniforme-p}
\planteamiento calcular la función de distribución de $p$ usando que $F_0$ es continua y estrictamente creciente, y por tanto invertible en su recorrido.
\pasos
\begin{enumerate}[label=\arabic*.]
  \item Sea $u\in(0,1)$. $p\le u\iff1-F_0(T)\le u\iff F_0(T)\ge1-u\iff T\ge F_0^{-1}(1-u)$.
  \item Bajo $H_0$, $\PP\bigl(T\ge F_0^{-1}(1-u)\bigr)=1-F_0\bigl(F_0^{-1}(1-u)\bigr)=1-(1-u)=u$ (por continuidad, $\PP(T=c)=0$).
  \item Luego $\PP_{H_0}(p\le u)=u$ para todo $u\in(0,1)$: $p\sim\Unif(0,1)$. Con $u=\alpha$, $\PP_{H_0}(p\le\alpha)=\alpha$.
\end{enumerate}
\resultado \textbf{Bajo $H_0$, el $p$-valor es uniforme en $(0,1)$ y la regla «rechazar si $p\le\alpha$» tiene nivel exactamente $\alpha$.} $\blacksquare$
\comprobacion este resultado justifica el apartado (a) del \cref{eje:t5-sidak}. Si el estadístico es discreto (por ejemplo, binomial), $p$ no es exactamente uniforme y el nivel real es $\le\alpha$.
\end{solucion}

\begin{solucion}{eje:t5-dualidad}
\planteamiento escribir la condición de no rechazo y despejar $\mu_0$.
\pasos
\begin{enumerate}[label=\arabic*.]
  \item Sea $c=z_{1-\alpha/2}$. El contraste no rechaza $H_0:\mu=\mu_0$ si y solo si $\abs{\xbar-\mu_0}\le c\,\sigma/\sqrt n$.
  \item $\abs{\xbar-\mu_0}\le c\sigma/\sqrt n\iff-c\sigma/\sqrt n\le\mu_0-\xbar\le c\sigma/\sqrt n\iff\xbar-c\sigma/\sqrt n\le\mu_0\le\xbar+c\sigma/\sqrt n$.
  \item La última condición es la pertenencia de $\mu_0$ al intervalo del \cref{teo:t4-ic-z}.
\end{enumerate}
\resultado \textbf{No se rechaza $\mu_0$ al nivel $\alpha$ si y solo si $\mu_0$ está en el intervalo de nivel $1-\alpha$.} El intervalo de confianza es el conjunto de valores del parámetro compatibles con los datos al nivel $\alpha$. $\blacksquare$
\comprobacion en el \cref{eje:t4-interpretacion}, el $0$ no está en $[10,20]$, luego un contraste bilateral al $5\,\%$ rechaza la mejora nula.
\end{solucion}

\begin{solucion}{eje:t5-potencia}
\datos $\pi(\mu)=1-\Phi\bigl(z_{1-\alpha}-(\mu-\mu_0)\sqrt n/\sigma\bigr)$.
\planteamiento monotonía de $\Phi$; resolver la desigualdad $\pi(\mu_1)\ge1-\beta$.
\pasos
\begin{enumerate}[label=\arabic*.]
  \item (a) El argumento $g(\mu)=z_{1-\alpha}-(\mu-\mu_0)\sqrt n/\sigma$ es estrictamente decreciente en $\mu$ y $\Phi$ es estrictamente creciente, luego $\pi=1-\Phi\circ g$ es estrictamente creciente. Además $\pi(\mu_0)=1-\Phi(z_{1-\alpha})=1-(1-\alpha)=\alpha$.
  \item (b) Si $\mu>\mu_0$, $g(\mu)\to-\infty$ cuando $n\to\infty$, y $\Phi(g)\to0$, luego $\pi(\mu)\to1$.
  \item (c) $\pi(\mu_1)\ge1-\beta\iff\Phi\bigl(z_{1-\alpha}-\delta\sqrt n/\sigma\bigr)\le\beta\iff z_{1-\alpha}-\delta\sqrt n/\sigma\le z_\beta=-z_{1-\beta}\iff\sqrt n\ge(z_{1-\alpha}+z_{1-\beta})\sigma/\delta$, y se elevan al cuadrado ambos miembros, positivos.
  \item (d) $n\ge\bigl((1{,}645+1{,}2816)\cdot80/20\bigr)^{2}=(11{,}706)^{2}\approx137{,}0$, luego $n=138$.
\end{enumerate}
\resultado \textbf{$\pi$ es creciente, vale $\alpha$ en $\mu_0$ y tiende a $1$; el tamaño muestral es el indicado; en el ejemplo, $n=138$ respuestas.} $\blacksquare$
\comprobacion con $n=138$: $\pi(220)=1-\Phi\bigl(1{,}645-20\sqrt{138}/80\bigr)=1-\Phi(1{,}645-2{,}937)=\Phi(1{,}292)\approx0{,}90$.
\end{solucion}

% ---------------------------------------------------------------------
\section{Tema 6}

\begin{solucion}{eje:t6-vf}
\planteamiento \cref{teo:t6-mc}, \cref{prop:t6-residuos}, \cref{cor:t6-r2} y \cref{obs:t6-errores}.
\pasos
\begin{enumerate}[label=(\alph*)]
  \item \textbf{Verdadera}: $\widehat b_0=\ybar-\widehat b_1\xbar$ equivale a $\ybar=\widehat b_0+\widehat b_1\xbar$.
  \item \textbf{Falsa}: $R^{2}=r^{2}=0{,}81\ge0$.
  \item \textbf{Falsa}: $R^{2}$ mide la calidad del ajuste lineal, no la causalidad.
  \item \textbf{Verdadera}, por la \cref{prop:t6-residuos}.
  \item \textbf{Falsa}: la pendiente de $x$ sobre $y$ es $r\,s_x/s_y$, que coincide con $1/\widehat b_1=s_x/(r\,s_y)$ solo si $r^{2}=1$.
\end{enumerate}
\resultado \textbf{(a) V; (b) F; (c) F; (d) V; (e) F.}
\end{solucion}

\begin{solucion}{eje:t6-basico}
\datos $n=5$; $\sum x_i=10$; $\sum y_i=15$; $\sum x_i^{2}=30$; $\sum y_i^{2}=55$; $\sum x_iy_i=38$.
\planteamiento sumas centradas con la fórmula de cálculo y \cref{teo:t6-mc}.
\pasos
\begin{enumerate}[label=\arabic*.]
  \item $\xbar=2$, $\ybar=3$; $S_{xx}=30-5\cdot4=10$; $S_{yy}=55-5\cdot9=10$; $S_{xy}=38-5\cdot2\cdot3=8$; $r=8/\sqrt{100}=0{,}8$.
  \item $\widehat b_1=8/10=0{,}8$; $\widehat b_0=3-0{,}8\cdot2=1{,}4$; $\widehat y=1{,}4+0{,}8x$.
  \item Ajustados: $1{,}4;\allowbreak\ 2{,}2;\allowbreak\ 3{,}0;\allowbreak\ 3{,}8;\allowbreak\ 4{,}6$. Residuos: $-0{,}4;\allowbreak\ 0{,}8;\allowbreak\ -1{,}0;\allowbreak\ 1{,}2;\allowbreak\ -0{,}6$.
  \item $\operatorname{SC}_E=0{,}16+0{,}64+1+1{,}44+0{,}36=3{,}6$; $R^{2}=1-3{,}6/10=0{,}64$; $s_e=\sqrt{3{,}6/3}=\sqrt{1{,}2}\approx1{,}095$.
  \item $\widehat y(2{,}5)=1{,}4+2=3{,}4$.
\end{enumerate}
\resultado \textbf{$\widehat y=1{,}4+0{,}8x$; $r=0{,}8$; $R^{2}=0{,}64$; $s_e\approx1{,}10$; predicción $3{,}4$.}
\comprobacion los residuos suman $0$ y $\operatorname{SC}_E=S_{yy}(1-r^{2})=10\cdot0{,}36=3{,}6$.
\end{solucion}

\begin{solucion}{eje:t6-gauss}
\datos $(1,8),(3,2),(5,1)$.
\planteamiento la log-verosimilitud total es una constante menos $\operatorname{SC}(b,a)/(2\sigma^{2})$; se maximiza minimizando $\operatorname{SC}$ (\cref{obs:t6-modelo}).
\pasos
\begin{enumerate}[label=\arabic*.]
  \item (a) $y_1\sim\Nor(ax_1+b,\sigma^{2})$: $\ln f(y_1)=-\ln\bigl(\sigma\sqrt{2\pi}\bigr)-\dfrac{(y_1-ax_1-b)^{2}}{2\sigma^{2}}$.
  \item (b) Sumando para los tres datos, $\ell(a,b)=-3\ln(\sigma\sqrt{2\pi})-\frac{1}{2\sigma^{2}}\sum_i(y_i-ax_i-b)^{2}$; para cualquier $\sigma$, el máximo en $(a,b)$ es el mínimo de la suma de cuadrados.
  \item $\xbar=3$; $\ybar=11/3$; $S_{xx}=35-27=8$; $S_{xy}=19-3\cdot3\cdot\frac{11}3=19-33=-14$.
  \item $\widehat a=-14/8=-1{,}75$; $\widehat b=\frac{11}{3}+1{,}75\cdot3=\frac{107}{12}\approx8{,}917$.
\end{enumerate}
\resultado \textbf{$\widehat a=-1{,}75$ y $\widehat b=107/12\approx8{,}92$}, independientes de $\sigma$.
\comprobacion residuos: $8-7{,}167=0{,}833$; $2-3{,}667=-1{,}667$; $1-0{,}167=0{,}833$; suman $0$ y $\sum x_ie_i=0{,}833-5+4{,}167=0$.
\end{solucion}

\begin{solucion}{eje:t6-examenes}
\datos $n=11$ y las sumas del enunciado.
\planteamiento fórmula de cálculo para las sumas centradas y \cref{teo:t6-mc}.
\pasos
\begin{enumerate}[label=\arabic*.]
  \item $\xbar=761/11\approx69{,}18$; $\ybar=1765/11\approx160{,}45$.
  \item $S_{xx}=52\,729-761^{2}/11\approx52\,729-52\,647{,}36=81{,}64$.
  \item $S_{xy}=122\,500-761\cdot1765/11\approx122\,500-122\,105{,}91=394{,}09$.
  \item $S_{yy}=287\,529-1765^{2}/11\approx287\,529-283\,202{,}27=4326{,}73$.
  \item $\widehat b_1=394{,}09/81{,}64\approx4{,}827$; $\widehat b_0=160{,}45-4{,}827\cdot69{,}18\approx-173{,}51$.
  \item $r=394{,}09/\sqrt{81{,}64\cdot4326{,}73}\approx394{,}09/594{,}31\approx0{,}663$.
  \item (b) $\widehat y(73)=-173{,}51+4{,}827\cdot73\approx178{,}9$.
  \item (c) $\widehat y(65)\approx-173{,}51+313{,}78=140{,}27$; residuo $175-140{,}27\approx34{,}7$.
\end{enumerate}
\resultado \textbf{$\widehat y=-173{,}51+4{,}83x$; $r\approx0{,}663$; predicción $\approx178{,}9$; residuo $\approx34{,}7$.}
\comprobacion coincide con la recta y la correlación de la fuente; $R^{2}\approx0{,}44$: el tercer examen explica menos de la mitad de la variabilidad del final.
\end{solucion}

\begin{solucion}{eje:t6-buceo}
\datos $n=6$; $\sum x_i=450$; $\sum y_i=262$; $\sum x_i^{2}=35\,500$; $\sum y_i^{2}=13\,784$; $\sum x_iy_i=17\,700$.
\planteamiento \cref{teo:t6-mc}; después, examinar los residuos.
\pasos
\begin{enumerate}[label=\arabic*.]
  \item $\xbar=75$; $\ybar\approx43{,}67$; $S_{xx}=35\,500-33\,750=1750$; $S_{xy}=17\,700-450\cdot262/6=17\,700-19\,650=-1950$; $S_{yy}=13\,784-11\,440{,}67=2343{,}33$.
  \item $\widehat b_1=-1950/1750\approx-1{,}114$; $\widehat b_0=43{,}67+1{,}114\cdot75\approx127{,}24$.
  \item $r=-1950/\sqrt{1750\cdot2343{,}33}\approx-1950/2025{,}0\approx-0{,}963$.
  \item (b) $\widehat y(85)\approx127{,}24-94{,}71=32{,}5$ minutos; $\widehat y(110)\approx127{,}24-122{,}57=4{,}7$ minutos.
  \item La primera es una interpolación (dentro de $[50,100]$) y es razonable; la segunda es una extrapolación y no es fiable.
  \item (c) Residuos: $8{,}5;\allowbreak\ -5{,}4;\allowbreak\ -4{,}2;\allowbreak\ -3{,}1;\allowbreak\ -2{,}0;\allowbreak\ 6{,}2$. Son positivos en los extremos y negativos en el centro: la relación es curva (convexa), y la recta subestima en los extremos.
\end{enumerate}
\resultado \textbf{$\widehat y=127{,}24-1{,}114x$; $r\approx-0{,}963$; predicciones $32{,}5$ (fiable) y $4{,}7$ (extrapolación); los residuos sugieren curvatura.}
\comprobacion los residuos suman $8{,}48-5{,}38-4{,}24-3{,}10-1{,}95+6{,}19\approx0$. A pesar de $R^{2}\approx0{,}93$, el patrón de residuos indica que un modelo curvo sería mejor (\cref{obs:t6-errores}(e)).
\end{solucion}

\begin{solucion}{eje:t6-helados}
\datos $n=9$ y las sumas del enunciado.
\planteamiento \cref{teo:t6-mc} y \cref{cor:t6-r2}.
\pasos
\begin{enumerate}[label=\arabic*.]
  \item $\xbar=654/9\approx72{,}67$; $\ybar=1288/9\approx143{,}11$.
  \item $S_{xx}=47\,988-654^{2}/9=47\,988-47\,524=464$.
  \item $S_{xy}=98\,132-654\cdot1288/9\approx98\,132-93\,594{,}67=4537{,}33$.
  \item $S_{yy}=231\,026-1288^{2}/9\approx231\,026-184\,327{,}11=46\,698{,}89$.
  \item $\widehat b_1=4537{,}33/464\approx9{,}779$; $\widehat b_0=143{,}11-9{,}779\cdot72{,}67\approx-567{,}48$.
  \item $r=4537{,}33/\sqrt{464\cdot46\,698{,}89}\approx4537{,}33/4654{,}9\approx0{,}975$; $R^{2}\approx0{,}950$.
  \item $\widehat y(74)\approx-567{,}48+9{,}779\cdot74\approx156{,}1$.
\end{enumerate}
\resultado \textbf{$\widehat y=-567{,}48+9{,}78x$; $r\approx0{,}975$; $R^{2}\approx0{,}95$; unos $156$ helados.}
\comprobacion la ordenada en el origen, $-567$, no tiene interpretación (no se venden helados negativos a $0$ grados): $x=0$ está muy fuera del rango de los datos.
\end{solucion}

\begin{solucion}{eje:t6-resumen}
\datos $\xbar=10$; $\ybar=50$; $s_x=2$; $s_y=8$; $r=0{,}75$.
\planteamiento \cref{teo:t6-mc} con $\widehat b_1=r\,s_y/s_x$; para $x$ sobre $y$, intercambiar los papeles.
\pasos
\begin{enumerate}[label=\arabic*.]
  \item (a) $\widehat b_1=0{,}75\cdot8/2=3$; $\widehat b_0=50-3\cdot10=20$; $\widehat y=20+3x$; $\widehat y(12)=56$.
  \item (b) Pendiente $r\,s_x/s_y=0{,}75\cdot2/8=0{,}1875$; ordenada $10-0{,}1875\cdot50=0{,}625$; $\widehat x=0{,}625+0{,}1875\,y$.
  \item (c) $3\cdot0{,}1875=0{,}5625=0{,}75^{2}$. Despejando $x$ en la recta de (a) se obtendría $x=(y-20)/3$, de pendiente $1/3\neq0{,}1875$.
\end{enumerate}
\resultado \textbf{$\widehat y=20+3x$ (predicción $56$) y $\widehat x=0{,}625+0{,}1875y$; el producto de pendientes es $r^{2}=0{,}5625$.}
\comprobacion ambas rectas pasan por $(10,50)$: $20+30=50$ y $0{,}625+9{,}375=10$.
\end{solucion}

\begin{solucion}{eje:t6-loglog}
\datos $u_i=\log_{10}N_i=6,7,8,9$; $v_i=\log_{10}L_i=0{,}6990;\allowbreak\ 0{,}6128;\allowbreak\ 0{,}5185;\allowbreak\ 0{,}4314$.
\planteamiento recta de mínimos cuadrados de $v$ sobre $u$; $\widehat\alpha=-\widehat b_1$ y $\widehat a=10^{\widehat b_0}$.
\pasos
\begin{enumerate}[label=\arabic*.]
  \item $\bar u=7{,}5$; $\bar v=2{,}2617/4=0{,}565\,43$; $S_{uu}=2{,}25+0{,}25+0{,}25+2{,}25=5$.
  \item $S_{uv}=\sum(u_i-\bar u)v_i=-1{,}5\cdot0{,}6990-0{,}5\cdot0{,}6128+0{,}5\cdot0{,}5185+1{,}5\cdot0{,}4314=-0{,}448\,55$.
  \item $\widehat b_1=-0{,}448\,55/5\approx-0{,}0897$; $\widehat b_0=0{,}565\,43+0{,}0897\cdot7{,}5\approx1{,}2382$.
  \item $\widehat\alpha\approx0{,}0897$; $\widehat a=10^{1{,}2382}\approx17{,}3$; $L\approx17{,}3\,N^{-0{,}0897}$.
  \item (b) $\log_{10}\widehat L=1{,}2382-0{,}0897\cdot10\approx0{,}3411$; $\widehat L\approx10^{0{,}3411}\approx2{,}19$.
  \item (c) $N=10^{10}$ está fuera del rango observado ($10^{6}$ a $10^{9}$); la predicción supone que la ley de potencia sigue valiendo un orden de magnitud más allá.
\end{enumerate}
\resultado \textbf{$L\approx17{,}3\,N^{-0{,}090}$; predicción para $N=10^{10}$: $L\approx2{,}19$ (extrapolación).}
\comprobacion con la ley ajustada, $L(10^{6})=17{,}3\cdot10^{-0{,}538}\approx17{,}3\cdot0{,}290\approx5{,}02$, muy próximo al dato $5{,}0$; $r\approx-0{,}9999$.
\end{solucion}

\begin{solucion}{eje:t6-simpson}
\datos A: $45/50$ fáciles, $75/250$ difíciles. B: $200/250$ fáciles, $10/50$ difíciles.
\planteamiento calcular proporciones por grupo y en total, y observar la composición de cada evaluación.
\pasos
\begin{enumerate}[label=\arabic*.]
  \item A: fáciles $0{,}90$; difíciles $0{,}30$; total $120/300=0{,}40$.
  \item B: fáciles $0{,}80$; difíciles $0{,}20$; total $210/300=0{,}70$.
  \item A es mejor en cada tipo de pregunta, pero peor en el total: A se ha evaluado sobre todo con preguntas difíciles ($250$ de $300$) y B sobre todo con fáciles. La dificultad es una variable de confusión.
  \item Una comparación justa usa la misma mezcla para ambos. Con una mezcla al $50\,\%$: A obtiene $(0{,}90+0{,}30)/2=0{,}60$ y B, $(0{,}80+0{,}20)/2=0{,}50$.
\end{enumerate}
\resultado \textbf{A es mejor} (lo es en ambos estratos); la exactitud total favorece a B solo porque los dos modelos se evaluaron con mezclas distintas de dificultad (paradoja de Simpson).
\comprobacion con cualquier mezcla común con proporción $w$ de fáciles, A obtiene $0{,}30+0{,}60w$ y B, $0{,}20+0{,}60w$: A gana por $0{,}10$ siempre.
\end{solucion}

\begin{solucion}{eje:t6-cuadratica}
\datos $x=-2,\dots,2$; $y=x^{2}$.
\planteamiento calcular $S_{xy}$ y los residuos.
\pasos
\begin{enumerate}[label=\arabic*.]
  \item $\xbar=0$; $\ybar=10/5=2$; $S_{xy}=\sum x_iy_i-0=\sum x_i^{3}=-8-1+0+1+8=0$. Luego $r=0$.
  \item $\widehat b_1=0$; $\widehat b_0=2$; la recta es $\widehat y=2$.
  \item Residuos: $2,-1,-2,-1,2$. Forman una «U»: positivos en los extremos, negativos en el centro.
\end{enumerate}
\resultado \textbf{$r=0$ y la recta es horizontal, pero $y$ está determinada exactamente por $x$}; el patrón de los residuos revela la relación cuadrática.
\comprobacion $\operatorname{SC}_E=4+1+4+1+4=14=S_{yy}$, coherente con $R^{2}=0$.
\end{solucion}

\begin{solucion}{eje:t6-normales}
\planteamiento derivadas parciales de $\operatorname{SC}(b_0,b_1)=\sum_i(y_i-b_0-b_1x_i)^{2}$.
\pasos
\begin{enumerate}[label=\arabic*.]
  \item $\dfrac{\partial\operatorname{SC}}{\partial b_0}=-2\sum_i(y_i-b_0-b_1x_i)$ y $\dfrac{\partial\operatorname{SC}}{\partial b_1}=-2\sum_ix_i(y_i-b_0-b_1x_i)$.
  \item Igualando a cero se obtienen las ecuaciones normales: $nb_0+b_1\sum x_i=\sum y_i$ y $b_0\sum x_i+b_1\sum x_i^{2}=\sum x_iy_i$.
  \item De la primera, $b_0=\ybar-b_1\xbar$. Sustituyendo en la segunda: $(\ybar-b_1\xbar)n\xbar+b_1\sum x_i^{2}=\sum x_iy_i$, es decir, $b_1\bigl(\sum x_i^{2}-n\xbar^{2}\bigr)=\sum x_iy_i-n\xbar\ybar$, y $b_1=S_{xy}/S_{xx}$.
  \item Hessiana: $H=\begin{pmatrix}2n&2\sum x_i\\2\sum x_i&2\sum x_i^{2}\end{pmatrix}$, constante. Su elemento $(1,1)$ es $2n>0$ y su determinante es $4\bigl(n\sum x_i^{2}-(\sum x_i)^{2}\bigr)=4nS_{xx}>0$. Por el criterio de Sylvester, $H$ es definida positiva.
  \item Una función cuadrática con hessiana constante definida positiva es estrictamente convexa, y su único punto crítico es el mínimo global.
\end{enumerate}
\resultado \textbf{Se recupera $\widehat b_1=S_{xy}/S_{xx}$ y $\widehat b_0=\ybar-\widehat b_1\xbar$, y el punto crítico es el mínimo global.} $\blacksquare$
\comprobacion con los datos del \cref{ej:t6-recta}: $5b_0+15b_1=20$ y $15b_0+55b_1=66$ tienen solución $b_1=0{,}6$, $b_0=2{,}2$.
\end{solucion}

\begin{solucion}{eje:t6-residuos}
\planteamiento las ecuaciones normales del \cref{eje:t6-normales} se escriben en términos de los residuos $e_i=y_i-\widehat b_0-\widehat b_1x_i$.
\pasos
\begin{enumerate}[label=\arabic*.]
  \item La primera ecuación normal es $\sum_i(y_i-\widehat b_0-\widehat b_1x_i)=0$, es decir, $\sum_ie_i=0$. (Directamente: $e_i=(y_i-\ybar)-\widehat b_1(x_i-\xbar)$ y ambas desviaciones suman cero.)
  \item La segunda es $\sum_ix_i(y_i-\widehat b_0-\widehat b_1x_i)=0$, es decir, $\sum_ix_ie_i=0$. (Directamente: $\sum_ix_ie_i=\sum_i(x_i-\xbar)e_i=S_{xy}-\widehat b_1S_{xx}=0$.)
  \item $\sum_i\widehat y_ie_i=\widehat b_0\sum_ie_i+\widehat b_1\sum_ix_ie_i=0$.
  \item Como $y_i=\widehat y_i+e_i$ y $\sum e_i=0$, $\frac1n\sum\widehat y_i=\ybar$.
\end{enumerate}
\resultado \textbf{Las tres identidades son ciertas, y la media de los ajustados es $\ybar$.} $\blacksquare$
\comprobacion en el \cref{ej:t6-recta}: $\sum x_ie_i=-0{,}8+1{,}2+3{,}0-2{,}4-1{,}0=0$.
\end{solucion}

\begin{solucion}{eje:t6-origen}
\planteamiento derivar $g(b)=\sum_i(y_i-bx_i)^{2}$.
\pasos
\begin{enumerate}[label=\arabic*.]
  \item $g'(b)=-2\sum_ix_i(y_i-bx_i)=0\iff\sum x_iy_i=b\sum x_i^{2}\iff b=\sum x_iy_i/\sum x_i^{2}$.
  \item $g''(b)=2\sum x_i^{2}>0$: $g$ es una parábola convexa y el punto crítico es su mínimo global.
  \item La ecuación $g'(\widehat b)=0$ es precisamente $\sum_ix_ie_i=0$.
  \item Ejemplo: $(1,2),(2,2)$. $\widehat b=(2+4)/(1+4)=1{,}2$; residuos $2-1{,}2=0{,}8$ y $2-2{,}4=-0{,}4$; $\sum e_i=0{,}4\ne0$, mientras que $\sum x_ie_i=0{,}8-0{,}8=0$.
\end{enumerate}
\resultado \textbf{$\widehat b=\sum x_iy_i/\sum x_i^{2}$; $\sum x_ie_i=0$, pero en general $\sum e_i\neq0$}, porque sin ordenada en el origen no hay ecuación normal asociada a $b_0$. $\blacksquare$
\comprobacion con ordenada libre, esos dos puntos dan la recta horizontal $y=2$ con residuos nulos.
\end{solucion}

\begin{solucion}{eje:t6-transformaciones}
\planteamiento por la \cref{prop:t1-lineal}, $x_i'-\xbar'=\alpha(x_i-\xbar)$ e $y_i'-\ybar'=\gamma(y_i-\ybar)$.
\pasos
\begin{enumerate}[label=\arabic*.]
  \item $S_{x'y'}=\alpha\gamma S_{xy}$ y $S_{x'x'}=\alpha^{2}S_{xx}$, luego $\widehat b_1'=\alpha\gamma S_{xy}/(\alpha^{2}S_{xx})=(\gamma/\alpha)\widehat b_1$.
  \item Por el \cref{eje:t1-r-invariante}, $r'=\operatorname{signo}(\alpha\gamma)\,r$, luego $R^{2}=r'^{2}=r^{2}$.
  \item Recta ajustada: $\widehat y_i'-\ybar'=\widehat b_1'(x_i'-\xbar')=(\gamma/\alpha)\widehat b_1\,\alpha(x_i-\xbar)=\gamma(\widehat y_i-\ybar)$, de modo que $\widehat y_i'=\gamma\widehat y_i+\delta$.
  \item Residuos: $e_i'=y_i'-\widehat y_i'=\gamma y_i+\delta-\gamma\widehat y_i-\delta=\gamma e_i$.
  \item De metros a centímetros: $\alpha=100$, $\gamma=1$; la pendiente se divide por $100$.
\end{enumerate}
\resultado \textbf{$\widehat b_1'=(\gamma/\alpha)\widehat b_1$, $R^{2}$ invariante y $e_i'=\gamma e_i$}; al pasar $x$ a centímetros, la pendiente se divide por $100$ (es el cambio de $y$ por centímetro, no por metro). $\blacksquare$
\comprobacion la predicción es equivariante: transformar los datos y ajustar da lo mismo que ajustar y transformar la predicción.
\end{solucion}

\begin{solucion}{eje:t6-regresion-media}
\planteamiento escribir la recta como $\widehat y-\ybar=\widehat b_1(x-\xbar)$ con $\widehat b_1=r\,s_y/s_x$.
\pasos
\begin{enumerate}[label=\arabic*.]
  \item (a) $\widehat y-\ybar=r\dfrac{s_y}{s_x}(x-\xbar)$; dividiendo por $s_y$: $\dfrac{\widehat y-\ybar}{s_y}=r\,\dfrac{x-\xbar}{s_x}$, es decir, $\widehat z_y=r\,z_x$.
  \item (b) $\abs{\widehat z_y}=\abs r\,\abs{z_x}<\abs{z_x}$ si $\abs r<1$ y $z_x\ne0$: la predicción está más cerca de la media, en unidades de desviación típica.
  \item (c) $\widehat z_y=0{,}5\cdot2=1$.
\end{enumerate}
\resultado \textbf{$\widehat z_y=r\,z_x$; con $r=0{,}5$, un modelo a $+2$ desviaciones en la primera evaluación se predice a $+1$ en la segunda.} $\blacksquare$
\comprobacion si $\abs r=1$ no hay regresión a la media ($\widehat z_y=\pm z_x$); si $r=0$, la mejor predicción es la media. Es una razón para desconfiar del modelo que encabeza una única evaluación ruidosa.
\end{solucion}

% =====================================================================
\chapter{Tablas estadísticas}
\label{ap:tablas}
% =====================================================================

Las tablas se han calculado numéricamente con cuatro decimales (tres en la tabla de la $\chi^{2}$). Para $z<0$ se usa la simetría $\Phi(-z)=1-\Phi(z)$; para la $t$ de Student, $t_{k;1-q}=-t_{k;q}$.

\section{Distribución normal estándar}

\begin{table}[ht]
\centering\footnotesize
\caption{Función de distribución $\Phi(z)=\PP(Z\le z)$ de $Z\sim\Nor(0,1)$. La fila da las unidades y la primera cifra decimal de $z$; la columna, la segunda cifra decimal.}\label{tab:phi}
\setlength{\tabcolsep}{3.6pt}
\begin{tabular}{@{}c|cccccccccc@{}}
\toprule
$z$ & 0{,}00 & 0{,}01 & 0{,}02 & 0{,}03 & 0{,}04 & 0{,}05 & 0{,}06 & 0{,}07 & 0{,}08 & 0{,}09 \\
\midrule
0{,}0 & 0{,}5000 & 0{,}5040 & 0{,}5080 & 0{,}5120 & 0{,}5160 & 0{,}5199 & 0{,}5239 & 0{,}5279 & 0{,}5319 & 0{,}5359 \\
0{,}1 & 0{,}5398 & 0{,}5438 & 0{,}5478 & 0{,}5517 & 0{,}5557 & 0{,}5596 & 0{,}5636 & 0{,}5675 & 0{,}5714 & 0{,}5753 \\
0{,}2 & 0{,}5793 & 0{,}5832 & 0{,}5871 & 0{,}5910 & 0{,}5948 & 0{,}5987 & 0{,}6026 & 0{,}6064 & 0{,}6103 & 0{,}6141 \\
0{,}3 & 0{,}6179 & 0{,}6217 & 0{,}6255 & 0{,}6293 & 0{,}6331 & 0{,}6368 & 0{,}6406 & 0{,}6443 & 0{,}6480 & 0{,}6517 \\
0{,}4 & 0{,}6554 & 0{,}6591 & 0{,}6628 & 0{,}6664 & 0{,}6700 & 0{,}6736 & 0{,}6772 & 0{,}6808 & 0{,}6844 & 0{,}6879 \\
0{,}5 & 0{,}6915 & 0{,}6950 & 0{,}6985 & 0{,}7019 & 0{,}7054 & 0{,}7088 & 0{,}7123 & 0{,}7157 & 0{,}7190 & 0{,}7224 \\
0{,}6 & 0{,}7257 & 0{,}7291 & 0{,}7324 & 0{,}7357 & 0{,}7389 & 0{,}7422 & 0{,}7454 & 0{,}7486 & 0{,}7517 & 0{,}7549 \\
0{,}7 & 0{,}7580 & 0{,}7611 & 0{,}7642 & 0{,}7673 & 0{,}7704 & 0{,}7734 & 0{,}7764 & 0{,}7794 & 0{,}7823 & 0{,}7852 \\
0{,}8 & 0{,}7881 & 0{,}7910 & 0{,}7939 & 0{,}7967 & 0{,}7995 & 0{,}8023 & 0{,}8051 & 0{,}8078 & 0{,}8106 & 0{,}8133 \\
0{,}9 & 0{,}8159 & 0{,}8186 & 0{,}8212 & 0{,}8238 & 0{,}8264 & 0{,}8289 & 0{,}8315 & 0{,}8340 & 0{,}8365 & 0{,}8389 \\
1{,}0 & 0{,}8413 & 0{,}8438 & 0{,}8461 & 0{,}8485 & 0{,}8508 & 0{,}8531 & 0{,}8554 & 0{,}8577 & 0{,}8599 & 0{,}8621 \\
1{,}1 & 0{,}8643 & 0{,}8665 & 0{,}8686 & 0{,}8708 & 0{,}8729 & 0{,}8749 & 0{,}8770 & 0{,}8790 & 0{,}8810 & 0{,}8830 \\
1{,}2 & 0{,}8849 & 0{,}8869 & 0{,}8888 & 0{,}8907 & 0{,}8925 & 0{,}8944 & 0{,}8962 & 0{,}8980 & 0{,}8997 & 0{,}9015 \\
1{,}3 & 0{,}9032 & 0{,}9049 & 0{,}9066 & 0{,}9082 & 0{,}9099 & 0{,}9115 & 0{,}9131 & 0{,}9147 & 0{,}9162 & 0{,}9177 \\
1{,}4 & 0{,}9192 & 0{,}9207 & 0{,}9222 & 0{,}9236 & 0{,}9251 & 0{,}9265 & 0{,}9279 & 0{,}9292 & 0{,}9306 & 0{,}9319 \\
1{,}5 & 0{,}9332 & 0{,}9345 & 0{,}9357 & 0{,}9370 & 0{,}9382 & 0{,}9394 & 0{,}9406 & 0{,}9418 & 0{,}9429 & 0{,}9441 \\
1{,}6 & 0{,}9452 & 0{,}9463 & 0{,}9474 & 0{,}9484 & 0{,}9495 & 0{,}9505 & 0{,}9515 & 0{,}9525 & 0{,}9535 & 0{,}9545 \\
1{,}7 & 0{,}9554 & 0{,}9564 & 0{,}9573 & 0{,}9582 & 0{,}9591 & 0{,}9599 & 0{,}9608 & 0{,}9616 & 0{,}9625 & 0{,}9633 \\
1{,}8 & 0{,}9641 & 0{,}9649 & 0{,}9656 & 0{,}9664 & 0{,}9671 & 0{,}9678 & 0{,}9686 & 0{,}9693 & 0{,}9699 & 0{,}9706 \\
1{,}9 & 0{,}9713 & 0{,}9719 & 0{,}9726 & 0{,}9732 & 0{,}9738 & 0{,}9744 & 0{,}9750 & 0{,}9756 & 0{,}9761 & 0{,}9767 \\
2{,}0 & 0{,}9772 & 0{,}9778 & 0{,}9783 & 0{,}9788 & 0{,}9793 & 0{,}9798 & 0{,}9803 & 0{,}9808 & 0{,}9812 & 0{,}9817 \\
2{,}1 & 0{,}9821 & 0{,}9826 & 0{,}9830 & 0{,}9834 & 0{,}9838 & 0{,}9842 & 0{,}9846 & 0{,}9850 & 0{,}9854 & 0{,}9857 \\
2{,}2 & 0{,}9861 & 0{,}9864 & 0{,}9868 & 0{,}9871 & 0{,}9875 & 0{,}9878 & 0{,}9881 & 0{,}9884 & 0{,}9887 & 0{,}9890 \\
2{,}3 & 0{,}9893 & 0{,}9896 & 0{,}9898 & 0{,}9901 & 0{,}9904 & 0{,}9906 & 0{,}9909 & 0{,}9911 & 0{,}9913 & 0{,}9916 \\
2{,}4 & 0{,}9918 & 0{,}9920 & 0{,}9922 & 0{,}9925 & 0{,}9927 & 0{,}9929 & 0{,}9931 & 0{,}9932 & 0{,}9934 & 0{,}9936 \\
2{,}5 & 0{,}9938 & 0{,}9940 & 0{,}9941 & 0{,}9943 & 0{,}9945 & 0{,}9946 & 0{,}9948 & 0{,}9949 & 0{,}9951 & 0{,}9952 \\
2{,}6 & 0{,}9953 & 0{,}9955 & 0{,}9956 & 0{,}9957 & 0{,}9959 & 0{,}9960 & 0{,}9961 & 0{,}9962 & 0{,}9963 & 0{,}9964 \\
2{,}7 & 0{,}9965 & 0{,}9966 & 0{,}9967 & 0{,}9968 & 0{,}9969 & 0{,}9970 & 0{,}9971 & 0{,}9972 & 0{,}9973 & 0{,}9974 \\
2{,}8 & 0{,}9974 & 0{,}9975 & 0{,}9976 & 0{,}9977 & 0{,}9977 & 0{,}9978 & 0{,}9979 & 0{,}9979 & 0{,}9980 & 0{,}9981 \\
2{,}9 & 0{,}9981 & 0{,}9982 & 0{,}9982 & 0{,}9983 & 0{,}9984 & 0{,}9984 & 0{,}9985 & 0{,}9985 & 0{,}9986 & 0{,}9986 \\
3{,}0 & 0{,}9987 & 0{,}9987 & 0{,}9987 & 0{,}9988 & 0{,}9988 & 0{,}9989 & 0{,}9989 & 0{,}9989 & 0{,}9990 & 0{,}9990 \\
3{,}1 & 0{,}9990 & 0{,}9991 & 0{,}9991 & 0{,}9991 & 0{,}9992 & 0{,}9992 & 0{,}9992 & 0{,}9992 & 0{,}9993 & 0{,}9993 \\
3{,}2 & 0{,}9993 & 0{,}9993 & 0{,}9994 & 0{,}9994 & 0{,}9994 & 0{,}9994 & 0{,}9994 & 0{,}9995 & 0{,}9995 & 0{,}9995 \\
3{,}3 & 0{,}9995 & 0{,}9995 & 0{,}9995 & 0{,}9996 & 0{,}9996 & 0{,}9996 & 0{,}9996 & 0{,}9996 & 0{,}9996 & 0{,}9997 \\
3{,}4 & 0{,}9997 & 0{,}9997 & 0{,}9997 & 0{,}9997 & 0{,}9997 & 0{,}9997 & 0{,}9997 & 0{,}9997 & 0{,}9997 & 0{,}9998 \\
\bottomrule
\end{tabular}
\end{table}

\begin{table}[ht]
\centering\small
\caption{Cuantiles de la normal estándar más usados: $\Phi(z_q)=q$.}\label{tab:zq}
\begin{tabular}{@{}lcccccccc@{}}
\toprule
$q$ & 0{,}80 & 0{,}90 & 0{,}95 & 0{,}96 & 0{,}975 & 0{,}98 & 0{,}99 & 0{,}995 \\
\midrule
$z_q$ & 0{,}8416 & 1{,}2816 & 1{,}6449 & 1{,}7507 & 1{,}9600 & 2{,}0537 & 2{,}3263 & 2{,}5758 \\
\bottomrule
\end{tabular}
\end{table}

\section{Distribución \texorpdfstring{$t$}{t} de Student}

\begin{table}[ht]
\centering\small
\caption{Cuantiles $t_{k;q}$ de la distribución $t$ con $k$ grados de libertad: $\PP(T_k\le t_{k;q})=q$. La última fila ($k=\infty$) son los cuantiles de la normal estándar.}\label{tab:t}
\begin{tabular}{@{}c|ccccc@{}}
\toprule
$k$ & $q=0{,}90$ & $0{,}95$ & $0{,}975$ & $0{,}99$ & $0{,}995$ \\
\midrule
1 & 3{,}078 & 6{,}314 & 12{,}706 & 31{,}821 & 63{,}657 \\
2 & 1{,}886 & 2{,}920 & 4{,}303 & 6{,}965 & 9{,}925 \\
3 & 1{,}638 & 2{,}353 & 3{,}182 & 4{,}541 & 5{,}841 \\
4 & 1{,}533 & 2{,}132 & 2{,}776 & 3{,}747 & 4{,}604 \\
5 & 1{,}476 & 2{,}015 & 2{,}571 & 3{,}365 & 4{,}032 \\
6 & 1{,}440 & 1{,}943 & 2{,}447 & 3{,}143 & 3{,}707 \\
7 & 1{,}415 & 1{,}895 & 2{,}365 & 2{,}998 & 3{,}499 \\
8 & 1{,}397 & 1{,}860 & 2{,}306 & 2{,}896 & 3{,}355 \\
9 & 1{,}383 & 1{,}833 & 2{,}262 & 2{,}821 & 3{,}250 \\
10 & 1{,}372 & 1{,}812 & 2{,}228 & 2{,}764 & 3{,}169 \\
11 & 1{,}363 & 1{,}796 & 2{,}201 & 2{,}718 & 3{,}106 \\
12 & 1{,}356 & 1{,}782 & 2{,}179 & 2{,}681 & 3{,}055 \\
13 & 1{,}350 & 1{,}771 & 2{,}160 & 2{,}650 & 3{,}012 \\
14 & 1{,}345 & 1{,}761 & 2{,}145 & 2{,}624 & 2{,}977 \\
15 & 1{,}341 & 1{,}753 & 2{,}131 & 2{,}602 & 2{,}947 \\
16 & 1{,}337 & 1{,}746 & 2{,}120 & 2{,}583 & 2{,}921 \\
17 & 1{,}333 & 1{,}740 & 2{,}110 & 2{,}567 & 2{,}898 \\
18 & 1{,}330 & 1{,}734 & 2{,}101 & 2{,}552 & 2{,}878 \\
19 & 1{,}328 & 1{,}729 & 2{,}093 & 2{,}539 & 2{,}861 \\
20 & 1{,}325 & 1{,}725 & 2{,}086 & 2{,}528 & 2{,}845 \\
21 & 1{,}323 & 1{,}721 & 2{,}080 & 2{,}518 & 2{,}831 \\
22 & 1{,}321 & 1{,}717 & 2{,}074 & 2{,}508 & 2{,}819 \\
23 & 1{,}319 & 1{,}714 & 2{,}069 & 2{,}500 & 2{,}807 \\
24 & 1{,}318 & 1{,}711 & 2{,}064 & 2{,}492 & 2{,}797 \\
25 & 1{,}316 & 1{,}708 & 2{,}060 & 2{,}485 & 2{,}787 \\
26 & 1{,}315 & 1{,}706 & 2{,}056 & 2{,}479 & 2{,}779 \\
27 & 1{,}314 & 1{,}703 & 2{,}052 & 2{,}473 & 2{,}771 \\
28 & 1{,}313 & 1{,}701 & 2{,}048 & 2{,}467 & 2{,}763 \\
29 & 1{,}311 & 1{,}699 & 2{,}045 & 2{,}462 & 2{,}756 \\
30 & 1{,}310 & 1{,}697 & 2{,}042 & 2{,}457 & 2{,}750 \\
40 & 1{,}303 & 1{,}684 & 2{,}021 & 2{,}423 & 2{,}704 \\
60 & 1{,}296 & 1{,}671 & 2{,}000 & 2{,}390 & 2{,}660 \\
120 & 1{,}289 & 1{,}658 & 1{,}980 & 2{,}358 & 2{,}617 \\
$\infty$ & 1{,}282 & 1{,}645 & 1{,}960 & 2{,}326 & 2{,}576 \\
\bottomrule
\end{tabular}
\end{table}

\section{Distribución \texorpdfstring{$\chi^{2}$}{chi cuadrado}}

\begin{table}[ht]
\centering\small
\caption{Cuantiles $\chi^{2}_{k;q}$ de la distribución $\chi^{2}$ con $k$ grados de libertad: $\PP(\chi^{2}_k\le\chi^{2}_{k;q})=q$.}\label{tab:chi2}
\begin{tabular}{@{}c|cccccccc@{}}
\toprule
$k$ & $q=0{,}005$ & $0{,}025$ & $0{,}05$ & $0{,}90$ & $0{,}95$ & $0{,}975$ & $0{,}99$ & $0{,}999$ \\
\midrule
1 & 0{,}000 & 0{,}001 & 0{,}004 & 2{,}706 & 3{,}841 & 5{,}024 & 6{,}635 & 10{,}828 \\
2 & 0{,}010 & 0{,}051 & 0{,}103 & 4{,}605 & 5{,}991 & 7{,}378 & 9{,}210 & 13{,}816 \\
3 & 0{,}072 & 0{,}216 & 0{,}352 & 6{,}251 & 7{,}815 & 9{,}348 & 11{,}345 & 16{,}266 \\
4 & 0{,}207 & 0{,}484 & 0{,}711 & 7{,}779 & 9{,}488 & 11{,}143 & 13{,}277 & 18{,}467 \\
5 & 0{,}412 & 0{,}831 & 1{,}145 & 9{,}236 & 11{,}070 & 12{,}833 & 15{,}086 & 20{,}515 \\
6 & 0{,}676 & 1{,}237 & 1{,}635 & 10{,}645 & 12{,}592 & 14{,}449 & 16{,}812 & 22{,}458 \\
7 & 0{,}989 & 1{,}690 & 2{,}167 & 12{,}017 & 14{,}067 & 16{,}013 & 18{,}475 & 24{,}322 \\
8 & 1{,}344 & 2{,}180 & 2{,}733 & 13{,}362 & 15{,}507 & 17{,}535 & 20{,}090 & 26{,}124 \\
9 & 1{,}735 & 2{,}700 & 3{,}325 & 14{,}684 & 16{,}919 & 19{,}023 & 21{,}666 & 27{,}877 \\
10 & 2{,}156 & 3{,}247 & 3{,}940 & 15{,}987 & 18{,}307 & 20{,}483 & 23{,}209 & 29{,}588 \\
11 & 2{,}603 & 3{,}816 & 4{,}575 & 17{,}275 & 19{,}675 & 21{,}920 & 24{,}725 & 31{,}264 \\
12 & 3{,}074 & 4{,}404 & 5{,}226 & 18{,}549 & 21{,}026 & 23{,}337 & 26{,}217 & 32{,}909 \\
13 & 3{,}565 & 5{,}009 & 5{,}892 & 19{,}812 & 22{,}362 & 24{,}736 & 27{,}688 & 34{,}528 \\
14 & 4{,}075 & 5{,}629 & 6{,}571 & 21{,}064 & 23{,}685 & 26{,}119 & 29{,}141 & 36{,}123 \\
15 & 4{,}601 & 6{,}262 & 7{,}261 & 22{,}307 & 24{,}996 & 27{,}488 & 30{,}578 & 37{,}697 \\
16 & 5{,}142 & 6{,}908 & 7{,}962 & 23{,}542 & 26{,}296 & 28{,}845 & 32{,}000 & 39{,}252 \\
17 & 5{,}697 & 7{,}564 & 8{,}672 & 24{,}769 & 27{,}587 & 30{,}191 & 33{,}409 & 40{,}790 \\
18 & 6{,}265 & 8{,}231 & 9{,}390 & 25{,}989 & 28{,}869 & 31{,}526 & 34{,}805 & 42{,}312 \\
19 & 6{,}844 & 8{,}907 & 10{,}117 & 27{,}204 & 30{,}144 & 32{,}852 & 36{,}191 & 43{,}820 \\
20 & 7{,}434 & 9{,}591 & 10{,}851 & 28{,}412 & 31{,}410 & 34{,}170 & 37{,}566 & 45{,}315 \\
25 & 10{,}520 & 13{,}120 & 14{,}611 & 34{,}382 & 37{,}652 & 40{,}646 & 44{,}314 & 52{,}620 \\
30 & 13{,}787 & 16{,}791 & 18{,}493 & 40{,}256 & 43{,}773 & 46{,}979 & 50{,}892 & 59{,}703 \\
\bottomrule
\end{tabular}
\end{table}



\backmatter
% =====================================================================
\begin{thebibliography}{99}
\addcontentsline{toc}{chapter}{Bibliografía}
\markboth{Bibliografía}{}
\small

\bibitem{bh}
J.~K. Blitzstein y J.~Hwang.
\newblock \emph{Introduction to Probability}, 2.ª ed.
\newblock CRC Press, 2019.
\newblock Texto del curso Stat~110 de Harvard. Capítulos~4 «Expectation», 5 «Continuous Random Variables», 6 «Moments», 7 «Joint Distributions» y 10 «Inequalities and Limit Theorems».

\bibitem{openstax}
B.~Illowsky y S.~Dean.
\newblock \emph{Introductory Statistics}, 2.ª ed.
\newblock OpenStax, Rice University, 2023. Libro de texto de acceso libre. Secciones 2.2--2.7, 7.1--7.2, 8.1--8.3, 9.1--9.5, 11.2 y 12.3--12.4.
\newblock \url{https://openstax.org/details/books/introductory-statistics-2e}

\bibitem{mit1805}
J.~Orloff y J.~Bloom.
\newblock 18.05: Introduction to Probability and Statistics.
\newblock Massachusetts Institute of Technology, MIT OpenCourseWare, primavera de 2022. Lecturas (\emph{readings}) 5c, 6b, 7b, 10a, 10b, 17b, 18, 19, 20, 22, 23a, 23b y 26; \emph{problem sets}~4, 5, 6, 8, 9, 10 y 11; problemas de la clase~23.
\newblock \url{https://ocw.mit.edu/courses/18-05-introduction-to-probability-and-statistics-spring-2022/}

\bibitem{rice}
J.~A. Rice.
\newblock \emph{Mathematical Statistics and Data Analysis}, 3.ª ed.
\newblock Duxbury (Thomson Brooks/Cole), 2007. Capítulos~8 «Estimation of Parameters and Fitting of Probability Distributions», 9 «Testing Hypotheses and Assessing Goodness of Fit», 10 «Summarizing Data» y 14 «Linear Least Squares».

\bibitem{stats200}
Z.~Fan.
\newblock Stats~200: Introduction to Statistical Inference.
\newblock Stanford University, otoño de 2016 (edición archivada \texttt{stats200.1172}). Clases~5, 7, 11, 12, 13, 18 y 25; \emph{homework}~3, 4 y~5.
\newblock \url{https://web.stanford.edu/class/archive/stats/stats200/stats200.1172/}

\bibitem{wasserman}
L.~Wasserman.
\newblock \emph{All of Statistics: A Concise Course in Statistical Inference}.
\newblock Springer, 2004. Capítulos~5 «Convergence of Random Variables», 6 «Models, Statistical Inference and Learning», 9 «Parametric Inference», 10 «Hypothesis Testing and p-values», 13 «Linear and Logistic Regression» y 16 «Causal Inference».

\end{thebibliography}


\end{document}
